\documentclass[11pt]{article}

\usepackage[margin=1in]{geometry}
\usepackage{amsmath,amssymb,amsthm}
\usepackage{mathtools}
\usepackage{booktabs}
\usepackage{microtype}
\usepackage{enumitem}
\usepackage[colorlinks=true,linkcolor=blue,citecolor=blue,urlcolor=blue]{hyperref}

\theoremstyle{plain}
\newtheorem{theorem}{Theorem}[section]
\newtheorem{proposition}[theorem]{Proposition}
\newtheorem{lemma}[theorem]{Lemma}
\newtheorem{corollary}[theorem]{Corollary}

\theoremstyle{definition}
\newtheorem{definition}[theorem]{Definition}

\theoremstyle{remark}
\newtheorem{remark}[theorem]{Remark}
\newtheorem{conjecture}[theorem]{Conjecture}

\newcommand{\N}{\mathbb{N}}
\newcommand{\im}{\operatorname{im}}
\newcommand{\Capability}{\operatorname{Cap}}
\newcommand{\Adm}{\operatorname{Adm}}
\newcommand{\Pol}{\mathsf{Pol}}
\newcommand{\realise}{\operatorname{Realise}}
\newcommand{\LocalAdm}{\mathcal{A}_{\mathrm{loc}}}
\newcommand{\eps}{\varepsilon}

\title{Effective Causal Descent:\\ Uniform Selection, Contextual Gluing, and Causal Histories}
\author{Anonymous}
\date{}

\begin{document}

\maketitle

\begin{abstract}
We study decision-making under partial observation through four
presentation-relative diagnostics: objectwise local validity, compatibility of a
declared diagram, effective realisability, and bounded realisability. The basic
constraint is observational uniformity: an interface cannot support different
responses to worlds it does not distinguish. The diagnostic labels are not
invariants of bare capability predicates; they record where success first fails
in a chosen local-to-global presentation. The static kernel is a fibrewise
selector lemma. Compactness, Helly, and join-tree results provide certificate and
gluing tools; parity, support contextuality, source integration, and Bell
coupling give finite representations. The obstruction datum records objectwise
local success and the nested global spaces
\(\Gamma_{\mathrm{set}}\supseteq\Gamma_{\mathrm{eff}}\supseteq
\Gamma_{\mathcal B}\); terminal compression preserves the nested global spaces
while erasing precisely the \(\mathsf D\)-gap. A recursive selector game yields
the binary DNR \(\Pi^0_1\) class and the PA oracle spectrum; uniform binary fibre
selection and nested-prefix path selection have Weihrauch degree \(\mathsf{WKL}\).
A primitive-recursive game language has a \(\Sigma^0_3\)-complete capability index
set, and counting finite source integrations is \(\#\mathrm P\)-complete. In the
causal tier, backward induction handles perfect information, activated-valid
partial strategies form sheaves, active sensing separates full-site local
failure from cover-relative matching failure, distributed selectors are
equivalent to blind-order one-shot arenas for pure-strategy capability, and
finite causal-poset protocols have confluent executions. Bell coupling failure
has a Farkas dual certificate, and mixed strategies on the Bell one-shot skeleton
realise exactly the finite local behaviours.
\end{abstract}

\tableofcontents

\section{Introduction}\label{sec:intro}

Consider a decision-maker acting from partial observations: a label on a box, a
sensor reading, a measurement record, or a local transcript. The hidden world may
contain more information than the record reveals. Once two worlds produce the
same record, however, the decision-maker has only one response available for
both. This observational uniformity is the elementary constraint from which the
selector problem begins.

The same structure appears in algorithms reading inputs through an interface,
physical instruments receiving measurement outcomes, distributed agents with
private signals, and sequential controllers acting through histories. In each
case, the policy is a function of the information presented to it, not of the
hidden world itself. Asking only whether a successful policy exists therefore
hides useful structure. A failure to find such a policy can occur at several
levels:
\begin{enumerate}[label=(\roman*)]
\item \emph{Local impossibility}: for a fixed observation, no single response is
valid for all worlds consistent with that observation.
\item \emph{Contextual incompatibility}: valid choices exist in each isolated
context, but cannot be made to agree across overlapping scopes.
\item \emph{Noncomputability}: compatible choices exist as a set-theoretic
family, but no computable implementation selects them.
\item \emph{Resource exhaustion}: computable implementations exist, but all of
them exceed the declared time, memory, communication, or other resource bound.
\end{enumerate}

Effective causal descent (ECD) separates these modes of failure inside a
declared typed presentation. The guiding question is not merely whether a valid
strategy exists, but at which structural layer an obstruction first appears. The
same underlying decision problem can be organized in more than one mathematical
way. A compressed presentation may preserve the yes--no capability predicate
while hiding the difference between a missing local choice and a failure of
compatibility. Conversely, a cover presentation may expose the incompatible
commitments that a full global object absorbs as local emptiness. The diagnostic
classification is therefore attached to a specific local-to-global
representation, not to a bare capability predicate.

A running quality-control example gives the intended reading. A single inspector
sees only the exterior record of a box and must choose a report--action pair. The
visible label is the observation; the hidden condition of the box is the world.
If the same label covers both acceptable and defective boxes, then no single
response is valid behind that label: the obstruction is local. If several
inspectors see different attributes and their individually consistent tables
disagree on shared attributes, the obstruction lies in combining the tables: it
is a compatibility failure. If the perfect table exists but encodes a
noncomputable set, the obstruction is effective. If the table is computable but
cannot be produced within the budget, the obstruction is resource-bounded. The
example is mnemonic rather than additional structure; the formal objects are
sets, maps, relations, diagrams, implementations, and loss functions.

Unless a theorem states otherwise, the set-theoretic development is carried out
in ZFC. Finite constructions use only finite choice, and the computability
examples provide explicit set-theoretic selectors when needed.

The paper proceeds in three stages. First, the minimal selector kernel isolates
observational uniformity. It consists of worlds, observable states, choices, a
validity relation, and a policy class. The fibrewise selector lemma shows that a
valid global policy is the same as a choice, at each realised observation, of an
action valid for every world in the corresponding fibre. The prediction-game
instantiation then explains where validity comes from: reports, disclosures,
actions, dynamics, queries, decoders, and losses.

Second, context diagrams and distributed selectors introduce overlap. Multiple
observation contexts are organized as assignment spaces with restriction maps.
Objectwise validity gives a constrained diagram; when validity is preserved by
restriction to smaller subcontexts, it gives a descent system in the
subfunctorial sense used in Definition~\ref{def:Av}. Set-theoretic, effective,
and resource-admissible global families are nested subspaces of one ambient
family space. The descent kernel uses ordinary restriction agreement; arbitrary
relational compatibility belongs to separate constraint extensions.

Third, the same typed language is used to locate standard finite, computable,
quantitative, and causal obstructions. The finite theory includes compactness
and Helly bounds for local certificates, join-tree gluing for acyclic context
structures, a six-object parity obstruction, support-level strong contextuality,
finite source integration, and Bell coupling. Constraint-datum isomorphisms
preserve compatible families; affine systems have a complete cokernel
obstruction; bijective constraint networks are governed by holonomy fixed
points; Bell coupling infeasibility has a Farkas dual certificate; and exact
counting of global integrations is $\#\mathrm P$-complete.

The computability tier gives a separate form of failure. The selector set of the
halting-stage game is a nonempty $\Pi^0_1$ class without computable members.
Nested countable distributed selectors give a second, overlapping-scope
$\mathsf K$ presentation of path selection. The DNR instance has the PA oracle
spectrum; its input-uniform path-selection operator, and uniform binary fibre
selection, have Weihrauch degree $\mathsf{WKL}$. Recursive capability is
$\Sigma^0_3$-complete for a concrete primitive-recursive presentation language.
Computational-class separation is distinguished from deterministic-time resource
separation. A quantitative interlude separates pointwise, observation-uniform,
effective, and bounded worst-case optima, with a Chebyshev-radius formula in
passive metric prediction.

The causal tier adds time while keeping the arena fixed before a strategy is
chosen. Backward induction proves that perfect-information node-local
nonemptiness already gives a global strategy. Activated patch semantics
represents partial contingent plans and cover-relative matching. Blind-order
one-shot arenas give an exact normal form for finite distributed selectors, and
every finite pure-strategy imperfect-information arena has an extensional
distributed compilation. On the Bell one-shot skeleton, probability measures over
pure strategies are exactly couplings of deterministic response tables. Finite
causal posets of universally executed cuts have unique confluent executions;
each linear extension gives a policy-independent arena preserving worldwise
loss. Contingent stage occurrence remains tree-indexed. A bounded-height
countable arena transfers the halting-stage effectivity obstruction to the
causal setting without asserting a general infinite-arena hierarchy.

Figure~\ref{fig:architecture} summarises the canonical presentations and the
nested realisation filters.

\subsection*{Relation to established frameworks}

The phrase ``effective causal descent'' denotes a presentation-level
classification of selection problems subject to restriction compatibility,
computability filters, and resource constraints. It is an analytical framework,
not effective descent in the Barr--Beck or Grothendieck sense.

The fibrewise selector problem is closely related to uniform constraint
satisfaction. Its multi-context form connects directly with relational database
joins: in acyclic database schemes, join trees and separator consistency explain
when local tables can be amalgamated into a single global table
\cite{beerifagin}. The contextual formulation follows the sheaf-theoretic
treatment of empirical models by Abramsky and Brandenburger, where measurement
contexts are local observation scopes and contextuality is the non-existence of a
global section compatible with the local data \cite{abramsky}. The Bell
presentation uses the deterministic-strategy coupling perspective: an empirical
behaviour is local precisely when it is represented by a probability measure over
deterministic response tables with prescribed marginals. Fine's theorem is the
historically central result for the standard $2\times2\times2$ scenario
\cite{fine}, while the general finite statement follows from the local polytope
representation \cite{pitowsky}.

The computability arguments use standard admissible numberings, index sets, and
the theory of $\Pi^0_1$ classes \cite{kleene,jockuschsoare,rogers,soare}. Viewing
policies as choices from represented multivalued maps connects the framework to
computable analysis and Weihrauch complexity \cite{brattkagherardi,weihrauch},
where reducibility calibrates the uniform algorithmic difficulty of mathematical
choice principles. The causal tier uses the standard extensive-form distinction
between nature moves, information sets, and contingent strategies
\cite{osbornerubinstein}. The finite-poset protocol semantics is related to
partial-order models of concurrency, including event structures and Mazurkiewicz
traces; the specific contribution here is the loss-preserving translation from
universally executed cut posets to imperfect-information arenas. The
least-fixed-point appendix uses standard domain theory \cite{abramskyjung}. Past
factorisation is deliberately weaker than a structural causal model; comparison
with structural causal semantics concerns possible extensions rather than an
identification \cite{pearl}.

Many component theorems used here are standard in their native settings. Their
role is to supply structured reference cases for one presentation-relative
selector formalism. The contribution is the common typed comparison: local
validity, compatibility, effective realisation, resource bounds, and their
preservation or failure under translation are represented in one language.

\begin{figure}[htbp]
\centering
\[
\begin{array}{rclcl}
\textsf{single observer}
&\longmapsto&
\textsf{fibre-indexed selector}
&\longmapsto&
\textsf{local validity}\;\mathsf L
\\[1mm]
\textsf{distributed observers}
&\longmapsto&
\textsf{CSP incidence diagram}
&\longmapsto&
\textsf{compatibility}\;\mathsf D
\\[1mm]
\textsf{extensive arena}
&\longmapsto&
\textsf{causal patch sheaf}
&\longmapsto&
\textsf{backward induction / matching}
\\[2mm]
&&
\Gamma_{\mathrm{set}}
&\supseteq&
\Gamma_{\mathrm{eff}}\supseteq\Gamma_{\mathcal B}
\end{array}
\]
\caption{Canonical presentations of the static, distributed, and causal tiers.
The first two arrows indicate how modelling data are packaged as local objects;
the last column records the diagnostic level exposed by that packaging.
Effectivity and resource bounds are imposed afterwards by nested realisation
images in the common global-family space; the diagnostic exposed by a row is
presentation-relative.}
\label{fig:architecture}
\end{figure}

\section{Minimal selector kernel}\label{sec:kernel}

At its most fundamental level, a selection problem asks whether a
decision-maker can choose valid responses while observing only partial
information about the underlying world. The minimal kernel contains no causality,
computation, probability, or resource bound. It records only observational
uniformity: a policy is a function of what is seen, so all worlds in one
observation fibre must be handled by one choice.

\begin{definition}[Selector problem]\label{def:selector}
A selector problem is a tuple $(W,Z,C,I,V,\Pi)$ consisting of:
\begin{enumerate}[label=(\roman*)]
\item a nonempty set $W$ of underlying worlds;
\item a nonempty set $Z$ of observable states;
\item a nonempty set $C$ of available choices;
\item an observation interface $I:W\to Z$;
\item a validity relation $V\subseteq W\times C$, where $V(w,c)$ means that
choice $c$ is valid in world $w$;
\item a declared policy class $\Pi\subseteq C^Z$.
\end{enumerate}
For $z\in\im(I)$ define the common admissible-choice set
\[
\Adm_V(z)=\{c\in C:\forall w\in I^{-1}(z),\ V(w,c)\}.
\]
The problem is capable in $\Pi$ when some $\pi\in\Pi$ satisfies
$V(w,\pi(I(w)))$ for every $w\in W$.
\end{definition}

\begin{remark}[Operational reading]
In plain terms, $\Adm_V(z)$ is the menu of choices that work for every world
that looks like $z$. If this menu is empty, the observation has already lost too
much information. If every menu is nonempty, the next question is whether the
policy class can choose one item from each menu.
\end{remark}

\begin{lemma}[Fibrewise selector lemma]\label{lem:selector}
For every selector problem,
\[
\exists\pi\in\Pi\ \forall w\in W\ V(w,\pi(I(w)))
\iff
\exists\pi\in\Pi\ \forall z\in\im(I)\ \pi(z)\in\Adm_V(z).
\]
\end{lemma}

\begin{proof}
Suppose the left side holds through $\pi$. Fix $z\in\im(I)$ and
$w\in I^{-1}(z)$. Then $I(w)=z$, so $V(w,\pi(z))$. Since $w$ was arbitrary,
$\pi(z)\in\Adm_V(z)$. This proves the right side.

Conversely, suppose the right side holds through $\pi$, and let $w\in W$.
Putting $z=I(w)$ gives $w\in I^{-1}(z)$ and
$\pi(z)\in\Adm_V(z)$, hence $V(w,\pi(I(w)))$. This holds for every world.
\end{proof}

\begin{remark}
Lemma~\ref{lem:selector} is a representation lemma. It says that checking a
policy world by world is the same as checking it fibre by fibre: at each
observable state, the policy must pick a choice that works for all hidden worlds
behind that state. The common-choice sets $\Adm_V(z)$ are therefore the local
objects on which the rest of the paper builds.
\end{remark}

\section{Prediction-game instantiation}\label{sec:games}

To model interactive systems, the minimal selector kernel is instantiated inside
a typed prediction game. In this setting, validity is no longer primitive. It is
computed from a report, a disclosure channel, an action, an environmental update,
a query of the resulting history, a decoder, and a loss function.

\begin{definition}[Typed prediction game]\label{def:game}
A typed prediction game is a tuple
\[
G=(W,H,Z,R,A,M,Y,I,D,U,q,e,\ell),
\]
where the carriers $W,Z,R,A,M,H,Y$ are nonempty. They represent, respectively,
worlds, histories, observations, reports, actions, disclosed messages, and query
values. The structural maps are
\[
I:W\to Z,\quad D:R\to M,\quad U:W\times M\times A\to H,
\]
\[
q:H\to Y,\quad e:R\to Y,\quad \ell:R\times Y\to[0,\infty].
\]
Thus $I$ observes the world, $D$ records what a report discloses, $U$ updates the
world-history after disclosure and action, $q$ asks the relevant question of the
history, $e$ decodes the report, and $\ell$ measures loss. The loss acts directly
on a report and a query value; in closed loop it is evaluated as $\ell(r,q(h))$.
The game is called \emph{exact} when $\ell(r,y)=0$ if and only if $e(r)=y$; this
condition constrains only the zero level of the loss.

A policy is $\pi=(\pi_R,\pi_A)$ with
\[
\pi_R:Z\to R,\qquad \pi_A:Z\times R\to A.
\]
Its induced choice and closed-loop history are
\[
s_\pi(z)=(\pi_R(z),\pi_A(z,\pi_R(z))),
\]
\[
h_\pi(w)=U(w,D(\pi_R(I(w))),\pi_A(I(w),\pi_R(I(w)))).
\]
Only the values of $\pi_A$ on the graph of $\pi_R$ affect capability. The larger
action type is retained because it records how an implementation would react to
counterfactual reports; the selector kernel uses only $s_\pi$.
\end{definition}

\begin{remark}[Report, action, and observation]
The closed-loop order matters. At observation $z$, the policy reports
$r=\pi_R(z)$ and then acts by $a=\pi_A(z,r)$. The disclosure map records what the
report reveals to the system, and the dynamics records how the world responds to
that disclosed report and action. Exactness identifies zero loss with a correct
report; positive and infinite losses record quantitative error and hard failure.
\end{remark}

\begin{definition}[Policy classes]\label{def:realise}
Write
\[
\Pi_{\mathrm{all}}=\Pol(Z;R,A)=R^Z\times A^{Z\times R}.
\]
Let $\mathrm{Impl}_{\mathrm{eff}}(G)$ be a represented class of effective
implementations and
\[
\realise_{\mathrm{eff}}:\mathrm{Impl}_{\mathrm{eff}}(G)
\to\Pi_{\mathrm{all}}
\]
its realisation map. If $\mathcal B$ is a resource predicate on implementations,
put
\[
\mathrm{Impl}_{\mathcal B}
=\{\iota\in\mathrm{Impl}_{\mathrm{eff}}(G):\mathcal B(\iota)\},\quad
\Pi_{\mathrm{eff}}=\im(\realise_{\mathrm{eff}}),\quad
\Pi_{\mathcal B}=\im(\realise_{\mathrm{eff}}|_{\mathrm{Impl}_{\mathcal B}}).
\]
Thus $\Pi_{\mathcal B}\subseteq\Pi_{\mathrm{eff}}\subseteq\Pi_{\mathrm{all}}$.
Resource predicates are imposed on implementations; extensionally equal policies
may have implementations with different costs.
\end{definition}

\begin{remark}[Policies versus implementations]
The hierarchy separates what a policy does from how it is realised. A successful
input--output table may exist as a set, may be computable by some implementation,
and may still fail a declared time, memory, or architectural bound. These are
different filters on the same ambient space of possible policy behaviours.
\end{remark}



\begin{definition}[Observational policy equivalence]\label{def:gauge}
For policies $\pi,\pi'\in\Pi_{\mathrm{all}}$, write
$\pi\sim_{\mathrm{obs}}\pi'$ when
\[
s_\pi(z)=s_{\pi'}(z)\qquad\text{for every }z\in\im(I).
\]
\end{definition}

\begin{proposition}[Observational policy quotient]\label{prop:gauge}
For every tolerance, the closed-loop loss and validity of an individual policy
depend only on its $\sim_{\mathrm{obs}}$-class, and
\[
\Pi_{\mathrm{all}}/\!\sim_{\mathrm{obs}}
\cong (R\times A)^{\im(I)}.
\]
For every policy class $\Pi$, capability in $\Pi$ is determined by its image in
the quotient; saturation is needed only to reconstruct $\Pi$ itself as a union
of complete equivalence classes.
\end{proposition}

\begin{proof}
Closed-loop loss uses a policy only through
$(\pi_R(z),\pi_A(z,\pi_R(z)))$ at realised interface values, proving the first
claim. The fibres of the map
$\pi\mapsto s_\pi|_{\im(I)}$ are exactly the
$\sim_{\mathrm{obs}}$-classes. It is surjective: given
$s(z)=(r_z,a_z)$, define $\pi_R(z)=r_z$ and
$\pi_A(z,r_z)=a_z$ on the image, and use fixed elements of the nonempty sets
$R,A$ to extend the maps elsewhere. The quotient is therefore isomorphic to the
displayed function space. A class contains a valid policy exactly when its
quotient image contains the corresponding valid principal selector, regardless
of whether the class is saturated.
\end{proof}

\begin{definition}[Admissible report--action pairs]\label{def:admissible-pairs}
For $z\in\im(I)$ define
\[
\Adm(z)=\{(r,a)\in R\times A:
\forall w\in I^{-1}(z),\ e(r)=q(U(w,D(r),a))\}.
\]
Thus $\Adm(z)$ contains report--action pairs, not reports alone. A policy is an
exact valid selector when $s_\pi(z)\in\Adm(z)$ for every $z\in\im(I)$.
\end{definition}

\begin{definition}[Worst-case capability]\label{def:cap}
For $\eps\in[0,\infty)$ and $\Pi\subseteq\Pi_{\mathrm{all}}$, write
$\Capability^{wc}(G,\Pi;q,\eps)$ when some $\pi\in\Pi$ satisfies
\[
\ell(\pi_R(I(w)),q(h_\pi(w)))\le\eps
\qquad\text{for every }w\in W.
\]
\end{definition}

\begin{proposition}[Prediction-game selector representation]\label{prop:game-selector}
If $G$ is exact,
\[
\Capability^{wc}(G,\Pi;q,0)
\iff
\exists\pi\in\Pi\ \forall z\in\im(I)\ s_\pi(z)\in\Adm(z).
\]
\end{proposition}

\begin{proof}
Let
\[
S_\Pi=\{s_\pi:Z\to R\times A:\pi\in\Pi\}
\subseteq (R\times A)^Z
\]
be the principal-selector image of the policy class. Apply
Lemma~\ref{lem:selector} with $C=R\times A$, policy class $S_\Pi$, and
\[
V(w,(r,a))\iff e(r)=q(U(w,D(r),a)).
\]
A witness in $S_\Pi$ is, by definition, $s_\pi$ for some $\pi\in\Pi$.
Exactness identifies this relation with zero closed-loop loss.
\end{proof}

\subsection{Passive factorisation}

\begin{proposition}[Set-theoretic fibre factorisation]\label{prop:fibrefactor}
Let $\bar I:W\to\im(I)$ be the corestriction of $I$. For a map $f:W\to Y$,
the following are equivalent:
\begin{enumerate}[label=(\roman*)]
\item there is $g:\im(I)\to Y$ with $f=g\circ\bar I$;
\item $f$ is constant on every fibre of $I$.
\end{enumerate}
The factor $g$ is unique.
\end{proposition}

\begin{proof}
If $f=g\circ\bar I$, equal interface values give equal values of $f$. Conversely,
if $f$ is constant on fibres, define $g(z)$ to be $f(w)$ for any
$w\in I^{-1}(z)$. Fibre constancy makes this independent of the chosen world,
and $f=g\circ\bar I$. Every $z\in\im(I)$ has a preimage, so this equation also
forces uniqueness.
\end{proof}

\begin{proposition}[Passive prediction]\label{prop:passive}
Assume exact loss. Suppose $U(w,m,a)=U_0(w)$ for all $m,a$, and put $f=q\circ U_0$. Let
$g:\im(I)\to Y$ be the factor of Proposition~\ref{prop:fibrefactor}, when it
exists. Exact capability in $\Pi$ holds if and only if $f$ is constant on the
fibres of $I$ and there is $\pi\in\Pi$ such that
\[
e(\pi_R(z))=g(z)\qquad\text{for every }z\in\im(I).
\]
\end{proposition}

\begin{proof}
In the passive game the action and disclosure do not change the queried value,
so exact validity of $\pi$ is
$e(\pi_R(I(w)))=f(w)$ for every $w$. If it holds, $f$ is constant on fibres and
the displayed lifting equation follows. Conversely, fibre constancy and the
lifting equation give $e(\pi_R(I(w)))=g(I(w))=f(w)$ for every world; exact loss
then gives capability. Exact validity imposes no further condition on the action
component in this passive game, although membership in the declared policy class
$\Pi$ may still constrain it.
\end{proof}

\begin{remark}[Passive prediction]
In a passive game the action does not change what is being predicted. Exact
prediction is possible precisely when the target value is already determined by
the observation, up to the decoder. Thus two indistinguishable worlds with
different target values are the simplest information-theoretic obstruction.
\end{remark}

\subsection{Anti-validity and the uniform diagonal}

\begin{definition}[Anti-validity]\label{def:antivalid}
A relevant fibre $I^{-1}(z)$ is \emph{anti-valid} when for every $(r,a)$ there
is $w\in I^{-1}(z)$ such that
\[
e(r)\ne q(U(w,D(r),a)).
\]
Equivalently, $\Adm(z)=\varnothing$.
\end{definition}

\begin{definition}[Uniform diagonal]\label{def:diag}
A query has \emph{uniform diagonal form} when there is a fixed-point-free
$\delta:Y\to Y$ such that
\[
q(U(w,D(r),a))=\delta(e(r))
\]
for every $w,r,a$.
\end{definition}

\begin{proposition}[Elementary diagonal obstruction]\label{prop:diagonal}
Uniform diagonal form makes every relevant fibre anti-valid. Consequently exact
capability fails for every policy class.
\end{proposition}

\begin{proof}
For any $z\in\im(I)$, choose $w\in I^{-1}(z)$. For every $(r,a)$,
$q(U(w,D(r),a))=\delta(e(r))\ne e(r)$. Thus $\Adm(z)=\varnothing$ and
Proposition~\ref{prop:game-selector} precludes capability.
\end{proof}

\begin{remark}[Diagonal form]
A uniform diagonal is the extreme local obstruction: every possible report is
turned into a different queried value. The obstruction is local because it
already empties each relevant fibre. The computability examples below use a
diagonal idea differently: their fibres remain nonempty, but no computable
selector can choose valid values for all fibres at once.
\end{remark}

\begin{proposition}[Distinct local failure shapes]\label{prop:failure-shapes}
There is a two-world obstruction that is not a uniform diagonal, and a uniform
diagonal obstruction with no two-world witness.
\end{proposition}

\begin{proof}
For the first, take two worlds with one interface value, binary reports,
singleton actions, passive dynamics, and distinct query values $0$ and $1$ on
the worlds. They form the pair obstruction of
Corollary~\ref{cor:worldpair}, but the query cannot be one fixed function of the
report on both worlds. For the second, take one world and set
$q(U(w,D(r),\ast))=1-r$ with identity decoder. This is a uniform diagonal, while
a two-world witness is impossible.
\end{proof}


\subsection{Robust validity and display-map semantics}

\begin{definition}[Robust admissibility]\label{def:eps}
For $\eps\in[0,\infty)$ and $z\in\im(I)$ put
\[
\Adm_\eps(z)=\{(r,a):\forall w\in I^{-1}(z),
\ \ell(r,q(U(w,D(r),a)))\le\eps\}.
\]
\end{definition}

\begin{lemma}[Robust selector representation]\label{lem:robust-selector}
For every $\eps\ge0$,
\[
\Capability^{wc}(G,\Pi;q,\eps)
\iff
\exists\pi\in\Pi\ \forall z\in\im(I)\ s_\pi(z)\in\Adm_\eps(z).
\]
\end{lemma}

\begin{proof}
If capability holds, restrict its inequality to each fibre and substitute
$I(w)=z$. This shows $s_\pi(z)\in\Adm_\eps(z)$. Conversely, if these memberships
hold and $w\in W$, then $w\in I^{-1}(I(w))$, so the defining inequality for
$\Adm_\eps(I(w))$ is exactly the capability inequality at $w$.
\end{proof}

\begin{definition}[Validity family and realised sections]\label{def:bundle}
The \emph{validity family} at tolerance $\eps$ has display map
\[
p^\eps:\mathrm{Tot}^\eps\to\im(I),\qquad
\mathrm{Tot}^\eps=
\{(z,r,a):z\in\im(I),\ (r,a)\in\Adm_\eps(z)\},
\]
with $p^\eps(z,r,a)=z$. No topology or local triviality is intended. A section
is a map $s:\im(I)\to\mathrm{Tot}^\eps$ with
$p^\eps\circ s=\mathrm{id}$, equivalently a principal part
$\widetilde s:\im(I)\to R\times A$ satisfying
$\widetilde s(z)\in\Adm_\eps(z)$. It is $\Pi$-realised when
$\widetilde s=s_\pi|_{\im(I)}$ for some $\pi\in\Pi$. Write
$\Gamma_\Pi(p^\eps)$ for these realised sections.
\end{definition}

\begin{lemma}[Section representation]\label{lem:section-representation}
\[
\Capability^{wc}(G,\Pi;q,\eps)
\iff
\Gamma_\Pi(p^\eps)\ne\varnothing.
\]
\end{lemma}

\begin{proof}
By definition, $\Gamma_\Pi(p^\eps)$ is nonempty exactly when some
$\pi\in\Pi$ has $s_\pi(z)\in\Adm_\eps(z)$ for every $z\in\im(I)$. Apply
Lemma~\ref{lem:robust-selector}.
\end{proof}

\begin{remark}[Section language]
The display map turns the admissible choices into a fibre bundle in the bare
set-theoretic sense: over each observation sits the set of choices that work
there. Capability asks for a section of this family, and realised capability asks
that the section come from the declared policy class. This is the
single-observer prototype of the multi-context gluing problems.
\end{remark}

\subsection{Monotonicity and interface refinement}

\begin{proposition}[Monotonicity]\label{prop:monotone}
For a fixed typed game:
\begin{enumerate}[label=(\roman*)]
\item enlarging the policy class preserves capability;
\item enlarging the tolerance preserves capability;
\item if $\varnothing\ne W'\subseteq W$, the same policy that is valid on $W$
is valid after restricting the game to $W'$.
\end{enumerate}
The third claim concerns a fixed policy class. Transport of an implementation
class under world restriction is separate modelling data.
\end{proposition}

\begin{proof}
The first claim keeps the same witness in the larger policy class. The second
keeps the same witness and uses $\eps\le\eps'$. For the third, every fibre of the
restricted interface is a subset of the corresponding original fibre, so every
report--action pair valid on the original fibre remains valid on the restricted
one. The policy itself is unchanged.
\end{proof}

\begin{proposition}[Interface refinement]\label{prop:refine}
Suppose $I=k\circ I'$ with $I':W\to Z'$. For a policy $\pi$ on $Z$, define
\[
k^*\pi=(\pi_R\circ k,\pi_A'),\qquad
\pi_A'(z',r)=\pi_A(k(z'),r).
\]
Let $\Pi_Z\subseteq\Pol(Z;R,A)$ and
$\Pi_{Z'}\subseteq\Pol(Z';R,A)$. If
$k^*(\Pi_Z)\subseteq\Pi_{Z'}$, capability for $(I,\Pi_Z)$ implies capability
for $(I',\Pi_{Z'})$.
\end{proposition}

\begin{proof}
For each $z'\in\im(I')$,
$I'^{-1}(z')\subseteq I^{-1}(k(z'))$. Hence every pair admissible on the latter
fibre is admissible on the former. Pulling a witnessing policy back by $k^*$
therefore yields a valid policy for the refined interface, and the image
condition places it in $\Pi_{Z'}$.
\end{proof}

\begin{remark}[Refinement]
A refined interface distinguishes worlds that the coarse interface may have
identified. Since each refined fibre is contained in a coarse fibre, refinement
cannot destroy an already valid observation-indexed policy once the policy class
admits the corresponding pullback.
\end{remark}

\section{Static local obstructions and past factorisation}\label{sec:static-local}

This section examines how local obstructions can be certified. A fibre is
locally impossible when the worlds behind one observation admit no common valid
choice. This is expressed as emptiness of an intersection of worldwise validity
sets. Compactness and convexity turn such failures into finite certificates;
past factorisation and finite propagation explain how causal limits generate
instances of the same obstruction.



\begin{definition}[Obstructing family]\label{def:obfamily}
For $w\in W$ let
\[
\operatorname{Val}(w)=\{(r,a):e(r)=q(U(w,D(r),a))\}.
\]
A subset $S\subseteq I^{-1}(z)$ is obstructing when
$\bigcap_{w\in S}\operatorname{Val}(w)=\varnothing$.
\end{definition}

\begin{lemma}[Fibre intersection identity]\label{lem:Wfam}
For every $z\in\im(I)$,
\[
\Adm(z)=\varnothing
\iff
I^{-1}(z)\text{ contains an obstructing family}.
\]
\end{lemma}

\begin{proof}
By definition,
$\Adm(z)=\bigcap_{w\in I^{-1}(z)}\operatorname{Val}(w)$. If it is empty, the
whole fibre is an obstructing family. If an obstructing
$S\subseteq I^{-1}(z)$ exists, the intersection over the full fibre is contained
in the empty intersection over $S$, and is therefore empty.
\end{proof}

\begin{definition}[Robust singleton-validity sets]\label{def:valeps}
For $w\in W$ and $\eps\ge0$ put
\[
\operatorname{Val}_\eps(w)=
\{(r,a)\in R\times A:
\ell(r,q(U(w,D(r),a)))\le\eps\}.
\]
\end{definition}

\begin{lemma}[Robust fibre intersection]\label{lem:robust-intersection}
\[
\Adm_\eps(z)=
\bigcap_{w\in I^{-1}(z)}\operatorname{Val}_\eps(w).
\]
\end{lemma}

\begin{proof}
Both sides consist exactly of the report--action pairs whose loss is at most
$\eps$ for every world in the fibre.
\end{proof}

\begin{theorem}[Compact local-obstruction certificate]\label{thm:compact-obstruction}
Suppose $R\times A$ is compact and every sublevel set
$\operatorname{Val}_\eps(w)$ is closed. This holds, for example, if for every
world the map
\[
(r,a)\longmapsto\ell(r,q(U(w,D(r),a)))
\]
is lower semicontinuous. Then $\Adm_\eps(z)=\varnothing$ if and only if some finite
$S\subseteq I^{-1}(z)$ satisfies
\[
\bigcap_{w\in S}\operatorname{Val}_\eps(w)=\varnothing.
\]
\end{theorem}

\begin{proof}
The reverse implication follows because the full intersection is contained in the
intersection over $S$. Conversely, if
every finite subfamily had nonempty
intersection, the closed sets would have the finite-intersection property.
Compactness of $R\times A$ would then make their total intersection nonempty,
contrary to $\Adm_\eps(z)=\varnothing$.
\end{proof}

\begin{remark}[Finite certificates]
The theorem turns an infinite-looking local failure into finite evidence. If no
choice works for all worlds behind one observation, then compactness supplies a
finite subfamily of worlds already witnessing the failure. Under convexity,
Helly's theorem bounds the size of such a certificate.
\end{remark}

\begin{corollary}[Helly bound]\label{cor:helly}
Suppose $R\times A$ is identified with a compact convex set of affine dimension
$d$ and every $\operatorname{Val}_\eps(w)$ is closed and convex. If
$\Adm_\eps(z)=\varnothing$, there is an obstructing
$S\subseteq I^{-1}(z)$ with $|S|\le d+1$.
\end{corollary}

\begin{proof}
Theorem~\ref{thm:compact-obstruction} gives a finite obstructing subfamily.
Helly's theorem applied in the affine hull yields an empty subfamily of size at
most $d+1$ \cite{eckhoff}.
\end{proof}

\begin{remark}[Randomised-choice instance]
In the selector-kernel specialisation, or in the typed game after taking the
report component singleton and $A=\Delta(C_0)$, let a finite primitive choice
set $C_0$ be replaced by the simplex of randomised choices. If primitive losses
$L_0(w,c)$ are finite and expected loss is
$L(w,p)=\sum_{c\in C_0}p(c)L_0(w,c)$, then the robust validity sets
$\{p\in\Delta(C_0):L(w,p)\le\eps\}$ are closed convex polytopes in a
simplex of affine dimension $|C_0|-1$. Corollary~\ref{cor:helly} gives an
obstruction certificate involving at most $|C_0|$ worlds.
\end{remark}

\begin{corollary}[World-pair obstruction]\label{cor:worldpair}
Suppose $I(w_0)=I(w_1)$ and
\[
q(U(w_0,D(r),a))\ne q(U(w_1,D(r),a))
\]
for every $(r,a)$. Then their common fibre has empty admissibility and exact
capability fails for every policy class.
\end{corollary}

\begin{proof}
No pair $(r,a)$ can have decoded report equal to both distinct query values, so
$\{w_0,w_1\}$ is obstructing. Apply Lemma~\ref{lem:Wfam} and
Proposition~\ref{prop:game-selector}.
\end{proof}

A cut $\tau$ is a declared boundary between accessible records and future
evolution. The following factorisation is additional structure. In the causal
tier, Proposition~\ref{prop:trace-relative-past} derives it from independently
supplied exogenous records after a causal predecessor trace is fixed.

\begin{definition}[Past-factored observation]\label{def:past-factor}
At a cut $\tau$, a past-factored observation for a cut interface
$I_\tau:W\to Z$ consists of independently specified accessible past
records $\mathsf P_\tau$ and maps
\[
P_\tau:W\to\mathsf P_\tau,\qquad J_\tau:\mathsf P_\tau\to Z
\]
with $I_\tau=J_\tau\circ P_\tau$. When this is meant to factor the original
game interface, one assumes $I_\tau=I$. Thus equal accessible past records imply
equal interface values; the converse is not assumed. The factorisation is causally
informative only because $P_\tau$ is supplied independently of $I_\tau$; the
extensive-form construction of Section~\ref{sec:feedback} provides such an
independent history structure.
\end{definition}

\begin{corollary}[Past-cone obstruction]\label{cor:pastcone}
For the prediction game whose interface is the cut interface $I_\tau$ of Definition~\ref{def:past-factor}
(or for the original game when $I_\tau=I$), if $P_\tau(w_0)=P_\tau(w_1)$ and
the separation condition of Corollary~\ref{cor:worldpair} holds, exact
capability fails for every policy class.
\end{corollary}

\begin{proof}
Past factorisation gives $I_\tau(w_0)=I_\tau(w_1)$, so
Corollary~\ref{cor:worldpair} applies.
\end{proof}

\begin{proposition}[Finite-propagation factorisation]\label{prop:lightcone}
Consider a deterministic cellular automaton of interaction radius $\rho$, an
observer window $V$, and horizon $T$. Let
$V^{+\rho T}$ be the cells within graph distance $\rho T$ of $V$. Any interface
computed from the states of $V$ through time $T$ factors through the initial
restriction
\[
P_T(w)=w|_{V^{+\rho T}}.
\]
\end{proposition}

\begin{proof}
By induction on time, the state of a cell after $t$ updates depends only on
initial cells within distance $\rho t$. Therefore every observer-window state
through time $T$, and every function of that record, is determined by
$w|_{V^{+\rho T}}$.
\end{proof}

\begin{corollary}[Light-cone prediction obstruction]\label{cor:lightcone}
For a passive target $f:W\to Y$, suppose two worlds agree on
$V^{+\rho T}$ but satisfy $f(w_0)\ne f(w_1)$. Then exact capability fails for
every policy class whose interface is confined to the observer record through
$T$.
\end{corollary}

\begin{proof}
Proposition~\ref{prop:lightcone} gives the same interface value in both worlds,
while the passive query separates them under every report--action pair.
Corollary~\ref{cor:worldpair} applies.
\end{proof}

\begin{remark}[Causal reach]
The light-cone obstruction is a fibre obstruction with causal content supplied by
the factorisation. If the observer record through time $T$ is determined by the
initial data in $V^{+\rho T}$, then targets depending on differences outside that
region cannot be exactly predicted from that record.
\end{remark}

\section{Context diagrams and effective descent}\label{sec:descent}

The discussion now broadens from one observation channel to many overlapping
contexts. A local assignment is a partial description of a possible global
policy, and compatibility means that two partial descriptions agree wherever
they overlap. Category theory supplies the bookkeeping: contexts are objects,
restriction maps are arrows, and a global assignment is a compatible family of
local choices. Relational CSP constraints can be added as a separate extension,
but the descent kernel itself uses ordinary restriction agreement. Arrows are
oriented from larger contexts to smaller contexts, so restriction is covariant;
reversing all arrows gives the usual presheaf convention.

\begin{definition}[Constrained assignment diagram]\label{def:restr}
A constrained assignment diagram consists of a small category $\mathcal C$, a
functor
\[
\mathcal A:\mathcal C\to\mathbf{Set},
\]
and objectwise admissibility sets
$\LocalAdm(c)\subseteq\mathcal A(c)$. Standard compatibility of a family
$s\in\prod_c\mathcal A(c)$ is
\[
\mathcal A(f)(s_c)=s_{c'}
\qquad(f:c\to c').
\]
The set of compatible admissible families is
\[
\Gamma_{\mathrm{set}}(\LocalAdm)=
\left\{s\in\prod_c\mathcal A(c):
 s_c\in\LocalAdm(c)\ \forall c,
 \ \mathcal A(f)(s_c)=s_{c'}\ \forall f:c\to c'\right\}.
\]
\end{definition}

\begin{remark}[Gluing interpretation]
A family $s$ chooses one local assignment at every context. It is compatible
when every overlap check passes. The diagnostic distinction between $\mathsf L$
and $\mathsf D$ is exactly the distinction between an empty local menu and a
collection of nonempty local menus whose choices cannot be stitched into one
consistent global assignment.
\end{remark}

\begin{definition}[Descent system]\label{def:Av}
A constrained assignment diagram is a \emph{descent system} when the
admissibility sets form a subfunctor: for every $f:c\to c'$,
\[
s_c\in\LocalAdm(c)
\Longrightarrow
\mathcal A(f)(s_c)\in\LocalAdm(c').
\]
\end{definition}

\begin{proposition}[Limit semantics of a descent system]\label{prop:limit-semantics}
For every descent system,
\[
\Gamma_{\mathrm{set}}(\LocalAdm)
=\varprojlim_{\mathcal C}\LocalAdm
=\left(\varprojlim_{\mathcal C}\mathcal A\right)
\cap\prod_c\LocalAdm(c),
\]
where all sets are viewed in the ambient product.
\end{proposition}

\begin{proof}
A family belongs to the limit of the subfunctor exactly when each component lies
in $\LocalAdm(c)$ and the same restriction equalities that define the limit of
$\mathcal A$ hold. These are precisely the conditions defining the displayed
intersection.
\end{proof}

\begin{theorem}[Initial-object collapse]\label{thm:initial-collapse}
Let $\mathcal C$ have an initial object $c_0$, with unique arrow
$!_c:c_0\to c$. Then
\[
\Gamma_{\mathrm{set}}(\LocalAdm)
\cong
\{u\in\LocalAdm(c_0):
\mathcal A(!_c)(u)\in\LocalAdm(c)\ \forall c\}.
\]
If $\LocalAdm$ is a subfunctor, this reduces to
\[
\Gamma_{\mathrm{set}}(\LocalAdm)\cong\LocalAdm(c_0).
\]
Consequently a subfunctorial diagram with an initial object cannot have outcome
$\mathsf D_\eps$ while all local sets are nonempty.
\end{theorem}

\begin{proof}
Compatibility forces $s_c=\mathcal A(!_c)(s_{c_0})$ for every object, so a
family is determined by its initial component. Conversely, any initial component
whose images are locally admissible defines a compatible family: for
$f:c\to c'$, uniqueness gives $f\circ !_c=!_{c'}$, and functoriality gives the
required restriction equality. If admissibility is a subfunctor, every image of
an admissible initial component is automatically admissible.
\end{proof}

\begin{remark}[Why initial objects remove compatibility failure]
An initial object acts as a complete context whose restrictions determine all
others. If validity is stable under restriction, any admissible assignment at
that complete context automatically supplies the matching local assignments.
Thus compatibility failure can appear only when the chosen presentation lacks
such a determining object or when admissibility is not restriction-stable.
\end{remark}

In this presentation, ``local'' means local to an object of the patch category;
the maximal full-tree patch is itself an object. Emptiness at that object is not
the same notion as node-local emptiness.

\begin{definition}[Realised family spaces]\label{def:ECD}
Let
\[
\sigma_{\mathrm{eff}}:\mathrm{Impl}_{\mathrm{eff}}
\to\prod_c\mathcal A(c)
\]
be a represented realisation map, and let $\sigma_{\mathcal B}$ be its
restriction to a declared bounded implementation class. Define
\[
\Gamma_{\mathrm{eff}}
=\im(\sigma_{\mathrm{eff}})\cap\Gamma_{\mathrm{set}},
\qquad
\Gamma_{\mathcal B}
=\im(\sigma_{\mathcal B})\cap\Gamma_{\mathrm{set}}.
\]
\end{definition}



\begin{proposition}[Fibre-indexed static diagram]\label{prop:fibre-diagram}
For a typed prediction game at tolerance $\eps$, take the discrete category with
objects $z\in\im(I)$, put
\[
\mathcal A(z)=R\times A,
\qquad
\LocalAdm_\eps(z)=\Adm_\eps(z),
\]
and send an implementation to the family
$(s_{\realise_{\mathrm{eff}}(\iota)}(z))_z$. Then effective, respectively
bounded, global-family existence is equivalent to capability in
$\Pi_{\mathrm{eff}}$, respectively $\Pi_{\mathcal B}$.
\end{proposition}

\begin{proof}
A global family on the discrete category is a map
$s:\im(I)\to R\times A$ with $s(z)\in\Adm_\eps(z)$. Requiring that family to be
realised gives exactly the selector condition of
Lemma~\ref{lem:robust-selector}.
\end{proof}

\begin{proposition}[Compressed terminal-category diagram]\label{prop:reduction}
For a typed prediction game at tolerance $\eps$, take the terminal category, with one object and only its identity arrow, and
put
\[
\mathcal A(c)=\operatorname{Maps}(\im(I),R\times A),
\]
\[
\LocalAdm_\eps(c)=
\{s:s(z)\in\Adm_\eps(z)\text{ for every }z\in\im(I)\}.
\]
Define
$\sigma_{\mathrm{eff}}(\iota)=
s_{\realise_{\mathrm{eff}}(\iota)}|_{\im(I)}$ and restrict it for
$\sigma_{\mathcal B}$. Then effective and bounded family existence are
respectively equivalent to capability in $\Pi_{\mathrm{eff}}$ and
$\Pi_{\mathcal B}$.
\end{proposition}

\begin{proof}
The terminal-category compatibility condition is vacuous. A realised ambient map lies
in $\LocalAdm_\eps(c)$ exactly when its policy is an $\eps$-valid selector.
Apply Lemma~\ref{lem:robust-selector}.
\end{proof}

\begin{proposition}[Discrete compatibility and choice]\label{prop:discreteVanishing}
On a discrete context category,
\[
\Gamma_{\mathrm{set}}(\LocalAdm)=\prod_c\LocalAdm(c).
\]
For finitely many objects, nonempty local sets give a global family by finite
choice. For an arbitrary small discrete category, the corresponding implication
uses the axiom of choice in the ambient ZFC foundation.
\end{proposition}

\begin{proof}
There are no nonidentity restriction equations, so a global family is exactly a
choice of one local element at every object.
\end{proof}

\begin{remark}[Presentation relativity]\label{rem:presentation}
The labels local, compatibility, effective, and bounded are relative to the
declared diagram and are not invariants of a bare capability predicate. Table~\ref{tab:presentations} summarises the structural cases used in the article; several rows list constructions appearing in Sections~\ref{sec:distributed}, \ref{sec:Dtri}, and \ref{sec:feedback}. Repackaging all constraints into one object can move a failure
between descriptions. The terminal-category construction of Proposition~\ref{prop:reduction} and the extensive-form
construction of Section~\ref{sec:feedback} are canonical relative to their input
data; externally supplied contextual diagrams carry their stated packaging as
part of the theorem.
\end{remark}







\begin{remark}[Principal specialisations]
The fibre-indexed discrete diagram gives the primary static prediction
specialisation; the terminal-category diagram is its compressed form. Allowing all
set-theoretic families with standard restriction agreement gives the contextual
specialisation. Intersecting global families with a represented implementation
image gives the effective specialisation. The downward-closed incidence diagram
of Definition~\ref{def:packaging} gives the empirical-support specialisation.
These are uses of one ambient section formalism, not claims that their native
mathematical theories are identical.
\end{remark}

\begin{definition}[Tolerance-filtered constrained diagram]\label{def:epsEC}
A tolerance-filtered constrained diagram has subsets
$\LocalAdm_\eps(c)\subseteq\mathcal A(c)$ satisfying, for every object,
\[
\eps\le\eps'
\Longrightarrow
\LocalAdm_\eps(c)\subseteq\LocalAdm_{\eps'}(c).
\]
It is a \emph{filtered descent system} when, in addition, every tolerance level
is restriction-stable:
\[
\mathcal A(f)(\LocalAdm_\eps(c))
\subseteq\LocalAdm_\eps(c')
\qquad(f:c\to c').
\]
Use the same realisation maps at every tolerance and define
\[
\Gamma_{\mathrm{set}}^\eps
=\Gamma_{\mathrm{set}}(\LocalAdm_\eps),\quad
\Gamma_{\mathrm{eff}}^\eps
=\im(\sigma_{\mathrm{eff}})\cap\Gamma_{\mathrm{set}}^\eps,
\quad
\Gamma_{\mathcal B}^\eps
=\im(\sigma_{\mathcal B})\cap\Gamma_{\mathrm{set}}^\eps,
\]
For a static prediction game also define the fibre-local success set
\[
T_{\mathrm{fib}}
=\{\eps:\forall z\in\im(I),\ \Adm_\eps(z)\ne\varnothing\}.
\]
The diagram-level success sets are
\begin{align*}
T_{\mathsf L}&=\{\eps:\forall c\ \LocalAdm_\eps(c)\ne\varnothing\},\\
T_{\mathsf D}&=\{\eps:\Gamma_{\mathrm{set}}^\eps\ne\varnothing\},\\
T_{\mathsf K}&=\{\eps:\Gamma_{\mathrm{eff}}^\eps\ne\varnothing\},\\
T_{\mathsf R}&=\{\eps:\Gamma_{\mathcal B}^\eps\ne\varnothing\}.
\end{align*}
Put $\eps_\chi=\inf T_\chi\in[0,\infty]$, with the convention
$\inf\varnothing=+\infty$.
For the terminal-category static compression, $T_{\mathrm{fib}}=T_{\mathsf L}$ in ZFC;
the equality uses choice when the observation image is infinite. In the
fibre-indexed discrete presentation, $T_{\mathrm{fib}}$ is objectwise local
success by definition.
\end{definition}

\section{Well-formedness, semantic outcomes, and transport}\label{sec:ladder}

This section turns the preceding distinctions into a diagnostic partition. The
questions are asked in order. First, does every declared local object have an
admissible element? Second, do these local elements glue to a compatible global
family? Third, is some global family effectively realised? Fourth, is some
effective realisation within the resource bound? The local-success bit separates
local failure from compatibility failure; the nested global spaces separate
set-theoretic, effective, and bounded realisation.

\begin{definition}[Well-formed presentation]\label{def:levels}
For each presentation language used in the article, a raw presentation is a code for the
corresponding carrier sets, maps, category, admissibility predicates, and
realisation maps. It is \emph{well formed}, written $\mathsf{WF}(P)$, when the
schema-specific type checks, category laws, functor laws, admissibility
predicates, and realisation maps have the required types. The diagnostic code
\[
\mathsf U(P)\iff\neg\mathsf{WF}(P)
\]
records malformed input. Semantic outcomes are defined only for
$\mathsf{WF}(P)$.
\end{definition}

\begin{definition}[Tolerance-relative semantic outcomes]\label{def:first}
Fix a well-formed presentation $P$, a tolerance $\eps$, an effective
realisation image, and a bounded implementation model. Use the filtered local
and global spaces of Definition~\ref{def:epsEC}, and define
\begin{align*}
\mathsf L_\eps(P)&:\quad
 \exists c\ [\LocalAdm_\eps(c)=\varnothing],\\
\mathsf D_\eps(P)&:\quad
 (\forall c\ \LocalAdm_\eps(c)\ne\varnothing)
 \land\Gamma_{\mathrm{set}}^\eps=\varnothing,\\
\mathsf K_\eps(P)&:\quad
 \Gamma_{\mathrm{set}}^\eps\ne\varnothing
 \land\Gamma_{\mathrm{eff}}^\eps=\varnothing,\\
\mathsf R_\eps(P)&:\quad
 \Gamma_{\mathrm{eff}}^\eps\ne\varnothing
 \land\Gamma_{\mathcal B}^\eps=\varnothing,\\
\checkmark_\eps(P)&:\quad
 \Gamma_{\mathcal B}^\eps\ne\varnothing.
\end{align*}
\end{definition}

\begin{proposition}[Invariance under represented natural isomorphism]
\label{prop:presentation-invariance}
Fix a tolerance and let two constrained diagrams on the same context category
be related by a natural isomorphism $\alpha:\mathcal A\Rightarrow\mathcal B$ satisfying
$\alpha_c(\LocalAdm_{\mathcal A}(c))=
\LocalAdm_{\mathcal B}(c)$ for every object. Then $\alpha$ induces a bijection
of compatible admissible-family spaces and preserves the $\mathsf L$ and
$\mathsf D$ outcomes. If $\alpha$ and its inverse are computable and intertwine
the effective realisation images, they also preserve $\mathsf K$. If they
intertwine the bounded images, they preserve $\mathsf R$.
\end{proposition}

\begin{proof}
Apply $\alpha$ componentwise. Naturality preserves every restriction equation,
and the local-set hypothesis preserves objectwise admissibility. The inverse
natural transformation gives the inverse map on global families. The additional
computability and image hypotheses carry membership in the effective and bounded
subsets in both directions.
\end{proof}

\begin{proposition}[Functoriality of global-family spaces]
\label{prop:global-family-functoriality}
Fix a context category and a tolerance. Let $P$ and $Q$ be constrained diagrams
with assignment functors $\mathcal A^P$ and $\mathcal A^Q$. If
$\alpha:\mathcal A^P\Rightarrow\mathcal A^Q$ is a natural transformation such
that
\[
\alpha_c(\LocalAdm^P(c))\subseteq\LocalAdm^Q(c)
\qquad(c\in\mathcal C),
\]
then componentwise application of $\alpha$ defines a map
\[
\Gamma_{\mathrm{set}}(P)\longrightarrow\Gamma_{\mathrm{set}}(Q).
\]
If $\alpha$ also carries the effective, respectively bounded, realised images
of $P$ into those of $Q$, this map restricts to
\[
\Gamma_{\mathrm{eff}}(P)\to\Gamma_{\mathrm{eff}}(Q),
\qquad
\Gamma_{\mathcal B}(P)\to\Gamma_{\mathcal B}(Q).
\]
\end{proposition}

\begin{proof}
For $s\in\Gamma_{\mathrm{set}}(P)$, put
$(\alpha_*s)_c=\alpha_c(s_c)$. Local admissibility is preserved by the displayed
hypothesis. For an arrow $f:c\to c'$, naturality gives
\[
\mathcal A^Q(f)(\alpha_c(s_c))
=\alpha_{c'}(\mathcal A^P(f)(s_c))
=\alpha_{c'}(s_{c'}),
\]
so $\alpha_*s$ is compatible. The effective and bounded restrictions are exactly
the stated image-preservation hypotheses.
\end{proof}

\begin{definition}[Extensional obstruction datum]
\label{def:obstruction-datum}
For a fixed well-formed presentation $P$ at tolerance $\eps$, put
\[
\lambda_\eps(P)=
\begin{cases}
1,&\LocalAdm_\eps(c)\ne\varnothing\text{ for every }c,\\
0,&\text{otherwise},
\end{cases}
\]
and define
\[
\mathfrak O_\eps(P)=
\left(
\lambda_\eps(P),
\Gamma_{\mathrm{set}}^\eps(P)
\supseteq\Gamma_{\mathrm{eff}}^\eps(P)
\supseteq\Gamma_{\mathcal B}^\eps(P)
\right).
\]
An isomorphism $\mathfrak O_\eps(P)\cong\mathfrak O_\eps(Q)$ is a bijection of
the set-theoretic global-family spaces carrying the effective and bounded
subsets onto their counterparts and preserving $\lambda_\eps$.
\end{definition}

\begin{definition}[Capability, global-family, and diagnostic equivalence]
\label{def:equivalence-levels}
A \emph{capability equivalence} between two presentations preserves the relevant
policy or strategy space and the worldwise loss functional, and therefore
preserves the capability predicate at every tolerance. A \emph{global-family
equivalence} preserves the nested global spaces
$\Gamma_{\mathrm{set}}\supseteq\Gamma_{\mathrm{eff}}\supseteq
\Gamma_{\mathcal B}$. A \emph{diagnostic equivalence} at tolerance $\eps$
preserves the full obstruction datum $\mathfrak O_\eps$, including the
objectwise local-success bit. Capability equivalence need not imply diagnostic
equivalence.
\end{definition}

\begin{proposition}[Obstruction-data transport]
\label{prop:obstruction-transport}
An isomorphism $\mathfrak O_\eps(P)\cong\mathfrak O_\eps(Q)$ preserves and
reflects each outcome
\[
\mathsf L_\eps,\quad\mathsf D_\eps,\quad
\mathsf K_\eps,\quad\mathsf R_\eps,\quad\checkmark_\eps.
\]
If only the three nested global spaces are preserved, then
$\mathsf K_\eps$, $\mathsf R_\eps$, and $\checkmark_\eps$ are preserved and
reflected, because they are Boolean combinations of emptiness facts about those
spaces. The distinction between $\mathsf L_\eps$ and $\mathsf D_\eps$ need not
be preserved without the local-success bit.
\end{proposition}

\begin{proof}
The local outcome is exactly $\lambda_\eps=0$. When $\lambda_\eps=1$, the
compatibility outcome is emptiness of $\Gamma_{\mathrm{set}}^\eps$. After that
space is nonempty, the effectivity outcome is emptiness of
$\Gamma_{\mathrm{eff}}^\eps$; after the effective space is nonempty, the
resource outcome is emptiness of $\Gamma_{\mathcal B}^\eps$. A bijection
carrying each distinguished subset onto its counterpart preserves all these
emptiness and nonemptiness conditions. Without the local-success bit, emptiness
of the set-theoretic global space does not record whether some local set was
already empty.
\end{proof}

\begin{remark}[Diagnostic data and certificate data]
The obstruction datum $\mathfrak O_\eps$ is sufficient for transporting the five
diagnostic outcomes because the outcomes use only local nonemptiness and the
emptiness pattern of the three nested global spaces. It is not enough to
transport the shape of a certificate. For that one must also remember the object
set, the local admissibility sets, the restriction maps, and the way local
elements occur as restrictions of compatible global families. Minimal diagnostic
transport and certificate transport are therefore distinct requirements.
\end{remark}

\begin{definition}[Terminal extensional compression]
\label{def:extensional-compression}
For a constrained diagram $P$ at tolerance $\eps$, let
$\operatorname{ExtComp}_\eps(P)$ be the terminal-category diagram with
\[
\mathcal A(*)=\prod_c\mathcal A(c),
\qquad
\LocalAdm_\eps(*)=\Gamma_{\mathrm{set}}^\eps(P).
\]
Use as effective and bounded realisation images the original images in the same
ambient product.
\end{definition}

\begin{theorem}[Compression boundary for $\mathsf L$ and $\mathsf D$]
\label{thm:compression-boundary}
For every constrained diagram $P$,
\[
\Gamma_\chi^\eps(\operatorname{ExtComp}_\eps(P))
=\Gamma_\chi^\eps(P)
\qquad
(\chi\in\{\mathrm{set},\mathrm{eff},\mathcal B\}).
\]
If $P$ has outcome $\mathsf D_\eps$, its extensional compression has outcome
$\mathsf L_\eps$. Consequently no equivalence that retains only the three
nested global-family spaces can preserve the distinction between local and
compatibility failure.
\end{theorem}

\begin{proof}
A terminal-category diagram has no nontrivial compatibility equations, so its global
set is its sole local set, namely $\Gamma_{\mathrm{set}}^\eps(P)$. Intersecting
with the unchanged effective and bounded images gives the other two displayed
equalities. If $P$ has outcome $\mathsf D_\eps$, all its local sets are nonempty
but $\Gamma_{\mathrm{set}}^\eps(P)=\varnothing$. The sole local set of the
compression is therefore empty, which is outcome $\mathsf L_\eps$.
\end{proof}

\begin{definition}[Obstruction spectrum and $\mathsf D$-gap]
\label{def:obstruction-spectrum}
For a tolerance-filtered presentation, its obstruction spectrum is
\[
\operatorname{Spec}(P)=
(T_{\mathsf L}(P),T_{\mathsf D}(P),T_{\mathsf K}(P),T_{\mathsf R}(P)),
\]
with the success sets of Definition~\ref{def:epsEC}. The compatibility gap is
\[
G_{\mathsf D}(P)=T_{\mathsf L}(P)\setminus T_{\mathsf D}(P).
\]
Thus $\eps\in G_{\mathsf D}(P)$ exactly when the presentation has outcome
$\mathsf D_\eps$.
\end{definition}

\begin{corollary}[Compression spectrum]
\label{cor:compression-spectrum}
For a tolerance-filtered presentation, define $\operatorname{ExtComp}(P)$
levelwise by applying Definition~\ref{def:extensional-compression} at every
$\eps$. Then
\[
\operatorname{Spec}(\operatorname{ExtComp}(P))=
(T_{\mathsf D}(P),T_{\mathsf D}(P),T_{\mathsf K}(P),T_{\mathsf R}(P)).
\]
Consequently $G_{\mathsf D}(\operatorname{ExtComp}(P))=\varnothing$, and for
each tolerance the only diagnostic state changed by terminal compression is
$\mathsf D_\eps$, which becomes $\mathsf L_\eps$.
\end{corollary}

\begin{proof}
The compressed presentation has one local set at tolerance $\eps$, namely
$\Gamma_{\mathrm{set}}^\eps(P)$; hence its local-success set is
$T_{\mathsf D}(P)$. The terminal category has no nontrivial compatibility
equations, so its set-theoretic global-family space is the same set
$\Gamma_{\mathrm{set}}^\eps(P)$. The effective and bounded images are unchanged
by construction, giving the $T_{\mathsf K}$ and $T_{\mathsf R}$ identities. The
statement about $G_{\mathsf D}$ and the fixed-tolerance outcomes follows from
the definitions of the five semantic outcomes.
\end{proof}

\begin{remark}
The compression boundary says exactly what is lost when all constraints are
packed into one object. The compression-invariant part of the spectrum is
$(T_{\mathsf D},T_{\mathsf K},T_{\mathsf R})$. The part erased by compression is
$G_{\mathsf D}$, the tolerance region where local validity holds but gluing
fails. Recovering that region requires local incidence and admissibility data in
addition to the nested global-family spaces.
\end{remark}

\begin{corollary}[Compression thresholds and idempotence]
\label{cor:compression-thresholds}
For levelwise terminal compression,
\[
\eps_{\mathsf L}(\operatorname{ExtComp}(P))=\eps_{\mathsf D}(P),\quad
\eps_{\mathsf D}(\operatorname{ExtComp}(P))=\eps_{\mathsf D}(P),\quad
\eps_{\mathsf K}(\operatorname{ExtComp}(P))=\eps_{\mathsf K}(P),\quad
\eps_{\mathsf R}(\operatorname{ExtComp}(P))=\eps_{\mathsf R}(P).
\]
Moreover, terminal compression is idempotent at the level of spectra and at each
fixed tolerance:
\[
\operatorname{Spec}(\operatorname{ExtComp}(\operatorname{ExtComp}(P)))
=\operatorname{Spec}(\operatorname{ExtComp}(P)).
\]
If the finite difference
\(\delta_{\mathsf D}(P)=\eps_{\mathsf D}(P)-\eps_{\mathsf L}(P)\) is defined,
then the compressed presentation has \(\delta_{\mathsf D}=0\).
\end{corollary}

\begin{proof}
The threshold identities follow by taking infima in the four success-set
identities of Corollary~\ref{cor:compression-spectrum}. Applying
Corollary~\ref{cor:compression-spectrum} a second time leaves the spectrum
unchanged because the first two components of the compressed spectrum are both
$T_{\mathsf D}(P)$. The final statement follows from the first two threshold
identities.
\end{proof}

\begin{remark}[Fixed-tolerance compression signature]
At a fixed tolerance, write
\[
\lambda=\mathbf 1[\text{all local sets are nonempty}],\quad
s=\mathbf 1[\Gamma_{\mathrm{set}}\ne\varnothing],\quad
e=\mathbf 1[\Gamma_{\mathrm{eff}}\ne\varnothing],\quad
b=\mathbf 1[\Gamma_{\mathcal B}\ne\varnothing].
\]
Nestedness gives $s\ge e\ge b$. Terminal compression acts on this Boolean
signature by
\[
(\lambda,s,e,b)\longmapsto(s,s,e,b).
\]
Thus the unique changed diagnostic state is
$(1,0,0,0)$, the $\mathsf D$ state, which is sent to the $\mathsf L$ state
$(0,0,0,0)$.
\end{remark}


\begin{proposition}[Diagnostic partition]\label{prop:diagnostic}
Every raw presentation has exactly one outcome among
$\mathsf U$ and, when well formed,
$\mathsf L_\eps,\mathsf D_\eps,\mathsf K_\eps,
\mathsf R_\eps,\checkmark_\eps$.
\end{proposition}

\begin{proof}
A malformed presentation has outcome $\mathsf U$. Assume $\mathsf{WF}(P)$.
If some local set is empty, $\mathsf L_\eps$ holds. Otherwise all local sets are
nonempty. Then either $\Gamma_{\mathrm{set}}^\eps$ is empty, giving
$\mathsf D_\eps$, or it is nonempty. In the latter case either
$\Gamma_{\mathrm{eff}}^\eps$ is empty, giving $\mathsf K_\eps$, or it is
nonempty. Finally either $\Gamma_{\mathcal B}^\eps$ is empty, giving
$\mathsf R_\eps$, or it is nonempty, giving $\checkmark_\eps$. The alternatives
are pairwise disjoint at the first condition on which they differ.
\end{proof}

\begin{remark}\label{rem:ladder}
The sequence is a diagnostic partition of a presentation at fixed tolerance and
fixed realisation models, not an intrinsic order on physical mechanisms. The
letter records the first failed filter in the chosen presentation; different
mechanism-specific certificates may witness the same failed filter.
Table~\ref{tab:ladder} summarises the outcomes, certificates, and representative inhabitants; representative results are cited in the table.
\end{remark}

\begin{table}[htbp]
\centering
\caption{Tolerance-relative outcomes and representative certificates.}
\label{tab:ladder}
\begin{tabular}{@{}p{0.09\linewidth}p{0.27\linewidth}p{0.28\linewidth}p{0.26\linewidth}@{}}
\toprule
Code & Defining condition & Certificate form & Representative result\\
\midrule
$\mathsf U$ & malformed presentation & failed typing or functor law & well-formedness check\\
$\mathsf L_\eps$ & some local set empty & finite or Helly witness & Theorem~\ref{thm:compact-obstruction}\\
$\mathsf D_\eps$ & locals nonempty; no matching family & incompatible scope constraints & Corollary~\ref{cor:team-decision}\\
$\mathsf K_\eps$ & global families exist; no effective one & nonempty $\Pi^0_1$ class without computable point & Theorems~\ref{thm:Pi01} and~\ref{thm:nested-distributed-K}\\
$\mathsf R_\eps$ & effective families exist; no bounded one & cost-class separation & Theorems~\ref{thm:primitive-class-separation} and~\ref{thm:timeR}\\
$\checkmark_\eps$ & bounded effective family exists & realised bounded section & capability\\
\bottomrule
\end{tabular}
\end{table}

\begin{table}[htbp]
\centering
\caption{Structural effect of canonical presentations on compatibility failure.}
\label{tab:presentations}
\begin{tabular}{@{}p{0.27\linewidth}p{0.14\linewidth}p{0.22\linewidth}p{0.25\linewidth}@{}}
\toprule
Presentation & Initial object? & Admissibility subfunctor? & Compatibility outcome\\
\midrule
fibre-indexed single observer & generally no & yes & no $\mathsf D$ in ZFC\\
terminal-category compression & yes & yes & no $\mathsf D$\\
full causal patch site & yes & yes & no $\mathsf D$\\
sensing cover incidence & no & yes & cover-relative $\mathsf D$ possible\\
distributed CSP incidence & generally no & not generally & $\mathsf D$ possible\\
nested reduced prefix chain & no & yes & set-theoretic gluing; $\mathsf K$ possible\\
empirical support presheaf & generally no & yes & $\mathsf D$ possible\\
Bell star diagram & yes & no & $\mathsf D$ possible\\
\bottomrule
\end{tabular}
\end{table}


\section{Distributed selectors and canonical compatibility presentations}\label{sec:distributed}

A single interface generates independent fibre constraints and cannot produce a
compatibility obstruction in its canonical discrete presentation in ZFC.
Distributed observation changes the situation. Several decision-makers may see
different slices of the same world, and the same local response rule must be
reused whenever the same local observation reappears. The obstruction is no
longer that one fibre has no admissible value; it is that locally admissible
response rules cannot be made to agree on their shared variables.

\begin{definition}[Distributed selector problem]\label{def:distributed}
A distributed selector problem consists of a nonempty world set $W$, a finite
agent set $J$, interfaces $I_j:W\to Z_j$, nonempty choice sets $C_j$, a nonempty
active scope $S(w)\subseteq J$ for each world, and relations
\[
V_w\subseteq\prod_{j\in S(w)}C_j.
\]
A strategy is a tuple $\pi=(\pi_j)_{j\in J}$ with
$\pi_j:Z_j\to C_j$. It is valid when
\[
(\pi_j(I_j(w)))_{j\in S(w)}\in V_w
\qquad\text{for every }w\in W.
\]
\end{definition}

\begin{definition}[Canonical distributed CSP and incidence diagram]
\label{def:distributed-csp}
For each agent define its active observation set
\[
Z_j^{\mathrm{act}}=\{I_j(w):w\in W,\ j\in S(w)\},
\qquad
\mathcal X_{\mathrm{act}}
=\{(j,z):j\in J,\ z\in Z_j^{\mathrm{act}}\}.
\]
The domain at $(j,z)$ is $C_j$, and the scope of $w$ is
\[
\operatorname{sc}(w)=\{(j,I_j(w)):j\in S(w)\}.
\]
Let the incidence category have as objects exactly the nonempty subsets of
realised scopes, with a unique arrow from $D$ to $E$ precisely when
$E\subseteq D$, and put
\[
\mathcal A(D)=\prod_{(j,z)\in D}C_j.
\]
For $w$ with $\operatorname{sc}(w)=D$, let
\[
\theta_w:\prod_{j\in S(w)}C_j\to\mathcal A(D)
\]
be the coordinate relabelling
$(c_j)_{j\in S(w)}\mapsto(a_{(j,I_j(w))})_{(j,I_j(w))\in D}$ with
$a_{(j,I_j(w))}=c_j$. For each realised scope define
\[
R_D=\bigcap_{\operatorname{sc}(w)=D}\theta_w(V_w)
\subseteq\mathcal A(D),
\]
and set
\[
\LocalAdm(D)=
\begin{cases}
R_D,&D\text{ is a realised scope},\\
\mathcal A(D),&\text{otherwise}.
\end{cases}
\]
Compatibility is coordinate restriction agreement. Every unrealised object has
$\LocalAdm(D)=\mathcal A(D)\ne\varnothing$; hence objectwise local nonemptiness
is equivalent to $R_D\ne\varnothing$ for every realised scope. This is generally
a constrained diagram rather than a subfunctor.
\end{definition}

\begin{remark}[Team interpretation]
In the distributed setting, each agent writes one table from observations to
choices. That table is reused across all worlds in which the agent sees the same
observation. Each world then tests a tuple of table entries. Compatibility can
fail because the same table entry is pulled into several tests and the tests
force contradictory values for it.
\end{remark}

\begin{theorem}[Distributed selector--CSP--section equivalence]
\label{thm:distributed-equivalence}
For unrestricted local policy classes, the following are equivalent:
\begin{enumerate}[label=(\roman*)]
\item a valid distributed strategy exists;
\item the canonical CSP has a solution;
\item the incidence diagram has a compatible admissible family.
\end{enumerate}
If some realised-scope relation $R_D$ is empty, the presentation has local
outcome $\mathsf L_0$. If all objectwise sets are nonempty but no global family
exists, it has compatibility outcome $\mathsf D_0$.
\end{theorem}

\begin{proof}
A strategy tuple restricts to an assignment on active variables by
$\alpha(j,z)=\pi_j(z)$. Conversely, since every $C_j$ is nonempty, an assignment
on $\mathcal X_{\mathrm{act}}$ extends by arbitrary values to observations never
used while agent $j$ is active, producing a full policy tuple. Validity at $w$
is exactly the constraint
$\alpha|_{\operatorname{sc}(w)}\in R_{\operatorname{sc}(w)}$, proving
(i)$\Leftrightarrow$(ii).

A compatible family is determined by singleton components on realised scopes,
which define an assignment on $\mathcal X_{\mathrm{act}}$. Compatibility forces
every scope component to be its restriction, and local admissibility forces
membership in $R_D$. Conversely, a CSP solution restricts to a compatible
admissible family. This proves (ii)$\Leftrightarrow$(iii). The outcome statements
are the definitions at tolerance zero.
\end{proof}


\begin{definition}[Relational constraint datum]\label{def:constraint-datum}
The relational constraint datum of a finite distributed selector problem is
\[
\mathfrak R(P)=
\bigl(\mathcal X_{\mathrm{act}},(C_x)_x,
\mathcal D,(R_D)_{D\in\mathcal D}\bigr),
\]
where $\mathcal D$ is the set of realised scopes and $C_{(j,z)}=C_j$.
An isomorphism $\mathfrak R(P)\to\mathfrak R(P')$ consists of a bijection
$\eta:\mathcal X_{\mathrm{act}}\to\mathcal X'_{\mathrm{act}}$ and domain
bijections $\delta_x:C_x\to C'_{\eta(x)}$ such that
$D\mapsto\eta(D)$ bijects the realised scopes and the product bijection
$\delta_D=\prod_{x\in D}\delta_x$ satisfies
\[
\delta_D(R_D)=R'_{\eta(D)}
\qquad(D\in\mathcal D).
\]
Thus the datum includes the relations, not merely the scope hypergraph or its
nerve.
\end{definition}

\begin{proposition}[Constraint-datum invariance]\label{prop:constraint-invariance}
An isomorphism $\mathfrak R(P)\cong\mathfrak R(P')$ induces canonical
bijections between the two CSP solution sets and between the compatible-family
spaces of their canonical incidence diagrams. It preserves and reflects
$\mathsf D_0$ whenever the realised-scope relations are nonempty.
\end{proposition}

\begin{proof}
For a CSP assignment $a$, define
\[
(\eta_*a)(\eta(x))=\delta_x(a(x)).
\]
For every realised scope $D$,
\[
a|_D\in R_D
\iff
(\eta_*a)|_{\eta(D)}=\delta_D(a|_D)\in R'_{\eta(D)}.
\]
Hence $a\mapsto\eta_*a$ bijects the solution sets, with inverse induced by
$\eta^{-1}$ and the inverse domain bijections. Theorem~\ref{thm:distributed-equivalence}
identifies each solution set with the compatible-family space of the canonical
incidence diagram. Objectwise nonemptiness is preserved by the same domain
bijections, so absence of a solution is a $\mathsf D_0$ outcome on one side
exactly when it is on the other.
\end{proof}

\begin{remark}
Proposition~\ref{prop:constraint-invariance} is an isomorphism theorem for the
full relational constraint datum. Scope incidence or homotopy type alone does
not determine compatibility, since changing the relations while retaining the
same scopes can change solvability.
\end{remark}


\begin{definition}[Solution-saturated constraint datum]
\label{def:saturated-constraint-datum}
Let $P$ be a finite distributed selector problem with nonempty CSP solution set
$S(P)$. For every active variable $x$, put
\[
C_x^{S}=\{a(x):a\in S(P)\},
\]
and, for every realised scope $D$, put
\[
R_D^{S}=\{a|_D:a\in S(P)\}
\subseteq\prod_{x\in D}C_x^{S}.
\]
The \emph{solution-saturated constraint datum}
$\mathfrak R^{S}(P)$ retains the active variables and realised-scope family of
$\mathfrak R(P)$ and replaces its domains and relations by
$(C_x^{S})_x$ and $(R_D^{S})_D$.

A coordinate isomorphism from $S(P)$ to $S(P')$ consists of a variable
bijection $\eta$ that bijects the realised-scope family onto the realised-scope
family, domain bijections
$\delta_x:C_x^{S}\to C_{\eta(x)}^{S'}$, and a bijection
$\Phi:S(P)\to S(P')$ satisfying
\[
\Phi(a)(\eta(x))=\delta_x(a(x))
\qquad(a\in S(P),\ x\in\mathcal X_{\mathrm{act}}).
\]
\end{definition}

\begin{proposition}[Recovery of saturated constraint data]
\label{prop:saturated-completeness}
For finite distributed selector problems with nonempty solution sets, the
following are equivalent:
\begin{enumerate}[label=(\roman*)]
\item their solution sets are coordinate-isomorphic in the sense of
Definition~\ref{def:saturated-constraint-datum};
\item their solution-saturated relational constraint data are isomorphic.
\end{enumerate}
Moreover, replacing $\mathfrak R(P)$ by $\mathfrak R^{S}(P)$ gives a
solution set canonically identified with the original one under the inclusions
$C_x^S\subseteq C_x$.
\end{proposition}

\begin{proof}
First consider the final assertion. Every $a\in S(P)$ has
$a|_D\in R_D^{S}$ for each realised scope, so it solves the saturated system.
Conversely, $R_D^{S}\subseteq R_D$ for every $D$. Hence every assignment solving
the saturated system solves the original system and belongs to $S(P)$. The inclusion of the support-reduced product into
the original product therefore canonically identifies the two solution sets.

Suppose (i) holds. For every realised scope $D$, coordinate preservation gives
\[
\delta_D(R_D^{S})
=\{\Phi(a)|_{\eta(D)}:a\in S(P)\}
=R_{\eta(D)}^{S'},
\]
where $\delta_D=\prod_{x\in D}\delta_x$. Thus $(\eta,(\delta_x)_x)$ is an
isomorphism of the saturated data, proving (ii).

Conversely, an isomorphism of the saturated data sends assignments
coordinatewise to assignments and preserves every saturated scope relation.
Proposition~\ref{prop:constraint-invariance} therefore gives a bijection of the
saturated solution sets. By the first paragraph these are $S(P)$ and $S(P')$,
and the induced bijection is coordinate-preserving. This proves (i).
\end{proof}

\begin{remark}
Proposition~\ref{prop:saturated-completeness} is a recovery theorem for nonempty,
support-reduced solution data. It is not a complete invariant of arbitrary
presentations or of unsolvable systems: an empty solution space cannot recover
unused domain elements or the local relations responsible for failure.
\end{remark}

\begin{corollary}[One-observer recovery]\label{cor:distributed-one}
If $J=\{j\}$, nonemptiness of $S(w)$ forces $S(w)=\{j\}$ for every world.
Hence every realised scope is a singleton, the incidence category is discrete,
and Theorem~\ref{thm:distributed-equivalence} reduces to fibrewise selection. Nonempty common-choice sets $\Adm_V(z)$ give a selector in ZFC, so $\mathsf D$ is absent.
\end{corollary}

\begin{proof}
Variables are the realised observations and constraints are precisely the
common-choice sets $\Adm_V(z)$. Apply discrete compatibility and choice.
\end{proof}

\begin{theorem}[Finite distributed capability is NP-complete]
\label{thm:distributed-NP}
For finite tabular distributed selector problems, deciding capability is
NP-complete.
\end{theorem}

\begin{proof}
Membership in NP follows by guessing one value for every variable $(j,z)$ and
checking every tabulated world relation.

For hardness, reduce 3SAT. Use one agent per Boolean variable, a singleton
observation set, and choice set $\{0,1\}$. For each clause create a world whose
active scope is the set of its variables and whose relation consists of the
assignments satisfying that clause (with signs interpreted in the
relation table). A distributed strategy is then a truth assignment, valid
exactly when it satisfies every clause.
\end{proof}

\begin{remark}[Compatibility is not effectivity]
Nonempty local relations do not imply a global distributed strategy. The team
problem of Corollary~\ref{cor:team-decision} has nonempty local relations but no global member, so its outcome
is $\mathsf D$, not $\mathsf K$. Section~\ref{sec:effect} supplies a distinct infinite represented
$\mathsf K$ example with nested distributed scopes.
\end{remark}



\begin{corollary}[Private-information team obstruction]
\label{cor:team-decision}
Let two agents receive private bits $x,y$ and respond with
$u=f(x),v=g(y)$. The four relations
\[
u\oplus v=x\wedge y
\]
are nonempty, but no distributed strategy satisfies all four. The induced
incidence diagram has outcome $\mathsf D_0$.
\end{corollary}

\begin{proof}
The inputs $(0,0),(0,1),(1,0)$ imply
$f(0)=g(0)=g(1)=f(1)$, while $(1,1)$ requires $f(1)\ne g(1)$.
\end{proof}

\begin{remark}[Deterministic CHSH core and Bell coupling]
Identifying $x,y$ with CHSH settings and $f(x),g(y)$ with deterministic local
responses makes these the CHSH-game winning equations. A deterministic response
table is a pair $(f,g)$, and Corollary~\ref{cor:team-decision} proves that no such
table wins all four inputs. This is the standard XOR-game form of CHSH, up to the conventional relabelling
between Boolean outputs and $\{\pm1\}$ outcomes. The Bell coupling theorem allows probability measures over the same deterministic tables. Thus the team theorem is the
deterministic compatibility obstruction, while Bell locality is its
distributional coupling relaxation. The one-shot extensive-form realisation is
Theorem~\ref{thm:oneshot}. Theorem~\ref{thm:finite-ii-compilation} gives
an extensional repeated-move compilation; preservation of temporal structure
beyond the one-shot normal form remains conjectural.
\end{remark}



The deterministic strategies of a finite Bell scenario are likewise distributed
strategies with parties as agents and settings as local observations. The XOR
team example is therefore the support-level deterministic core of the CHSH game:
it has locally satisfiable input relations but no single deterministic response
table satisfying all of them. The probability-valued coupling diagram of Definition~\ref{def:bellpack} is
their convex relaxation. The join-tree theorem gives the acyclic,
separator-consistent positive case for the scope hypergraph.




\section{Quantitative decomposition}\label{sec:quant}

A single gap between unrestricted and realised policies conflates information,
effectivity, and resources. Four optima are separated. Let
$C=R\times A$ and, for $c=(r,a)$, put
\[
L(w,c)=\ell(r,q(U(w,D(r),a))).
\]

\begin{definition}[Four worst-case degrees]\label{def:degree}
Define
\begin{align*}
d_{\mathrm{pt}}
&=\inf_{c:W\to C}\sup_{w\in W}L(w,c(w)),\\
d_{\mathrm{obs}}
&=\inf_{s:\im(I)\to C}\sup_{w\in W}L(w,s(I(w))),\\
d_{\mathrm{eff}}
&=\inf_{\pi\in\Pi_{\mathrm{eff}}}\sup_{w\in W}L(w,s_\pi(I(w))),\\
d_{\mathcal B}
&=\inf_{\pi\in\Pi_{\mathcal B}}\sup_{w\in W}L(w,s_\pi(I(w))),
\end{align*}
with $\inf\varnothing=+\infty$.
\end{definition}

\begin{proposition}[Nested degrees]\label{prop:tax}
\[
d_{\mathrm{pt}}
\le d_{\mathrm{obs}}
\le d_{\mathrm{eff}}
\le d_{\mathcal B}.
\]
\end{proposition}

\begin{proof}
A policy depending only on the observation is a special case of a world-dependent
choice after precomposition with $I$. Effective policy selectors form a subset
of all observation-indexed selectors, and bounded effective policies form a
subset of effective policies. Taking infima over successively smaller classes
gives the inequalities.
\end{proof}

\begin{proposition}[Chebyshev-radius formula]\label{prop:chebyshev}
Assume passive dynamics, singleton actions, report and target space equal to a
metric space $(Y,d)$, identity decoder, and loss $d$. Put $f=q\circ U_0$ and
\[
\operatorname{rad}(S)=\inf_{y\in Y}\sup_{x\in S}d(y,x)
\in[0,\infty],
\]
with value $\infty$ for unbounded $S$. Then
\[
d_{\mathrm{pt}}=0,
\qquad
d_{\mathrm{obs}}
=\sup_{z\in\im(I)}
\operatorname{rad}(f(I^{-1}(z))).
\]
Moreover
\[
\frac12\operatorname{diam}(S)
\le\operatorname{rad}(S)
\le\operatorname{diam}(S)
\]
for every nonempty bounded $S\subseteq Y$; for bounded subsets of $\mathbb R$,
the first inequality is an equality. If $I=k\circ I'$, then
\[
d_{\mathrm{obs}}(I')\le d_{\mathrm{obs}}(I).
\]
\end{proposition}

\begin{proof}
World-dependent reports choose $f(w)$ and give zero pointwise loss. For an
observation-indexed report $s$, the worst loss on fibre $z$ is
$\sup_{w\in I^{-1}(z)}d(s(z),f(w))$. Taking the infimum independently at each
fibre gives the displayed supremum of radii; in the declared ZFC foundation,
choose a $\delta$-approximate centre simultaneously in every fibre to establish
the upper direction when centres are not attained. Finite observation images
use only finite choice. The triangle inequality gives
$\operatorname{diam}(S)\le2\operatorname{rad}(S)$, while choosing a centre in
$S$ gives the upper bound. On $\mathbb R$, the midpoint of the infimum and
supremum is a centre. Refinement replaces each fibre image by a subset, whose
radius cannot increase.
\end{proof}



\begin{definition}[Adjacent gaps]\label{def:tax}
When the adjacent values are finite, define
\[
\Delta_{\mathrm{uniform}}=d_{\mathrm{obs}}-d_{\mathrm{pt}},\quad
\Delta_{\mathrm{effective}}=d_{\mathrm{eff}}-d_{\mathrm{obs}},\quad
\Delta_{\mathrm{resource}}=d_{\mathcal B}-d_{\mathrm{eff}}.
\]
These are respectively the observation-uniformity, effectivity, and resource
gaps.
\end{definition}



\begin{remark}[Quantitative gaps and diagnostic outcomes]
The numerical gaps in Definition~\ref{def:tax} are not the same data as the
semantic outcomes $\mathsf L,\mathsf D,\mathsf K,\mathsf R$. The gaps compare
optimal worst-case losses over nested classes of selectors, whereas the outcomes
classify emptiness of presentation-level local and global spaces at a fixed
tolerance. In particular, a single-observer discrete presentation can have a
positive observation-uniformity gap while $\mathsf D$ is impossible in ZFC.
\end{remark}

\begin{proposition}[Nested upward-closed success sets]\label{prop:spectrum}
Each $T_\chi$ is upward closed, and
\[
T_{\mathsf R}\subseteq T_{\mathsf K}\subseteq
T_{\mathsf D}\subseteq T_{\mathsf L},
\qquad
\eps_{\mathsf L}\le\eps_{\mathsf D}\le
\eps_{\mathsf K}\le\eps_{\mathsf R}.
\]
Status at an endpoint is determined by membership, not merely by comparison with
the infimum.
\end{proposition}

\begin{proof}
Increasing tolerance preserves every local validity inequality and hence every
witness at each successive level. The class inclusions give the nesting. Infima
reverse the displayed success-set inclusions in the stated order. An infimum
need not be attained, so endpoint membership must be checked separately.
\end{proof}

\begin{corollary}[Outcome at a fixed tolerance]\label{cor:tolerance-outcome}
For a well-formed problem and fixed $\eps$:
\begin{enumerate}[label=(\roman*)]
\item $\eps\notin T_{\mathsf L}$ gives outcome $\mathsf L_\eps$;
\item $\eps\in T_{\mathsf L}\setminus T_{\mathsf D}$ gives $\mathsf D_\eps$;
\item $\eps\in T_{\mathsf D}\setminus T_{\mathsf K}$ gives $\mathsf K_\eps$;
\item $\eps\in T_{\mathsf K}\setminus T_{\mathsf R}$ gives $\mathsf R_\eps$;
\item $\eps\in T_{\mathsf R}$ gives bounded capability.
\end{enumerate}
\end{corollary}

\begin{proof}
Apply the diagnostic partition of Proposition~\ref{prop:diagnostic} to the four
nested success conditions at tolerance $\eps$. Each case records the first
success set in the chain that does not contain $\eps$.
\end{proof}

\begin{proposition}[Robustness under class enlargement]\label{prop:robust}
Local emptiness is invariant under enlargement of the policy or implementation
class. Effectivity and resource outcomes are relative to the corresponding
realisation images.
\end{proposition}

\begin{proof}
The local sets $\LocalAdm_\eps(c)$ are defined before an implementation class is
imposed, so changing that class cannot populate an empty local set. By contrast,
$\Gamma_{\mathrm{eff}}^\eps$ and $\Gamma_{\mathcal B}^\eps$ are intersections
with class-dependent images; enlarging either image can turn the corresponding
empty intersection into a nonempty one.
\end{proof}

\section{Finite gluing, contextuality, and coupling}\label{sec:Dtri}

This section studies finite gluing from both sides. Some local data cannot be
assembled, and the obstruction can be seen in a parity cycle or in a strongly
contextual support model. Other local data glue for structural reasons: join
trees give the acyclic positive case, affine systems reduce to linear algebra,
and bijective networks reduce to holonomy fixed points. At the probabilistic
level, Bell locality becomes a coupling problem: the local probability tables
must be marginals of one probability measure on complete response tables.

\subsection{Parity as a compatibility obstruction}

\begin{definition}[Triangle assignment diagram]\label{def:triangle}
Let $\mathcal C^\triangle$ have as objects the nonempty proper subsets of
$\{0,1,2\}$ and an arrow $S\to T$ when $T\subseteq S$. Thus arrows point from
edges to their vertices. Put $\mathcal A(S)=\{0,1\}^S$ with ordinary
restriction, and define
\[
\LocalAdm(\{i\})=\{0,1\}^{\{i\}},\qquad
\LocalAdm(\{i,j\})=
\{v:v(i)+v(j)\equiv1\pmod2\}.
\]
Use standard compatibility.
\end{definition}

\begin{theorem}[Triangle compatibility obstruction]\label{thm:Dtri}
Every local admissibility set in the triangle is nonempty, but
$\Gamma_{\mathrm{set}}(\LocalAdm)=\varnothing$. Hence its semantic obstruction
is $\mathsf D$.
\end{theorem}

\begin{proof}
Each singleton admits two assignments and each edge admits the two odd-parity
assignments. Suppose a compatible family existed. Agreement along edge-to-vertex
arrows would determine bits $v(0),v(1),v(2)$ satisfying
\[
v(0)+v(1)\equiv v(1)+v(2)\equiv v(2)+v(0)\equiv1\pmod2.
\]
Adding the three equations gives
$2(v(0)+v(1)+v(2))\equiv1\pmod2$, while the left side is zero modulo two. This
contradiction excludes a compatible family.
\end{proof}

\begin{remark}\label{rem:triangle-lean}
As a constraint-satisfaction instance, the example is an inconsistent
odd-parity cycle. The six-object diagram adds explicit overlap objects and
restriction arrows, thereby exhibiting which local assignments must agree and
why the failure is one of gluing rather than local satisfiability. The Lean
artifact checks the parity core and local non-emptiness, not the full categorical
packaging.
\end{remark}

\subsection{Finite empirical support models}

\begin{definition}[Finite empirical model]\label{def:empirical}
Let $X$ be finite and let $M$ be a finite antichain of nonempty maximal
contexts with $X=\bigcup M$. Each measurement $m$ has a nonempty finite outcome set $O_m$.
A possibilistic empirical model assigns a nonempty support
$e_C\subseteq\prod_{m\in C}O_m$ to each $C\in M$ and satisfies overlap
compatibility:
\[
\{s|_{C\cap C'}:s\in e_C\}
=
\{s'|_{C\cap C'}:s'\in e_{C'}\}
\qquad(C,C'\in M)
\]
\cite{abramsky}. Strong contextuality means that no
$g\in\prod_{m\in X}O_m$ has $g|_C\in e_C$ for every $C\in M$.
\end{definition}

\begin{definition}[Downward-closed packaging]\label{def:packaging}
Let
\[
\downarrow M=\{D\subseteq X:D\subseteq C\text{ for some }C\in M\}.
\]
Use objects $D\in\downarrow M$, arrows $C\to D$ for $D\subseteq C$, and
\[
\mathcal A(D)=\prod_{m\in D}O_m
\]
with coordinate restriction. For any maximal $C\supseteq D$, define the
restricted support
\[
e_D=\{s|_D:s\in e_C\}.
\]
Overlap compatibility makes this independent of the choice of maximal $C$.
Set $\LocalAdm(D)=e_D$ for every $D\in\downarrow M$. For $D=\varnothing$ the
assignment space and support are the singleton empty assignment. Standard
compatibility defines $\mathrm{Pack}(e)$.
\end{definition}

\begin{proposition}[Restriction stability of empirical admissibility]\label{prop:empirical-subfunctor}
The admissibility sets of $\mathrm{Pack}(e)$ form a subfunctor of the assignment
functor.
\end{proposition}

\begin{proof}
Let $D'\subseteq D$ and $s\in e_D$. Choose a maximal context $C\supseteq D$ and
$t\in e_C$ with $t|_D=s$. Then
$s|_{D'}=t|_{D'}\in e_{D'}$. Thus coordinate restriction sends
$\LocalAdm(D)$ into $\LocalAdm(D')$, and identity and composition laws are
inherited from ordinary restriction.
\end{proof}

\begin{theorem}[Strong contextuality as gluing failure]\label{thm:KSfin}
A finite empirical model is strongly contextual if and only if
$\mathrm{Pack}(e)$ has nonempty local admissibility sets and no compatible
global family.
\end{theorem}

\begin{proof}
A compatible family determines one value at every singleton context. Coverage
combines these values into a global assignment $g$ on $X$. Compatibility then
forces the component at every $D$ to be $g|_D$, and maximal-context validity
gives $g|_C\in e_C$ for every $C\in M$. Conversely, any such global assignment
defines a compatible family by restriction. Thus a compatible family exists
exactly when the model is not strongly contextual. Since every local set is
nonempty, its absence is a $\mathsf D$ obstruction.
\end{proof}

\begin{proposition}[Empirical supports as distributed selectors]
\label{prop:empirical-distributed}
A finite empirical support model is the distributed selector problem whose
agents are measurements with singleton interfaces $Z_m=\{\ast\}$, whose world
set is the set of maximal contexts, whose active scope at context $C$ is $C$,
whose outcome choices are $O_m$, and whose
world relation is $e_C$. The restriction-stable support packaging of
Definition~\ref{def:packaging} refines its canonical incidence diagram and has
the same global families. Thus distributed capability holds exactly when the model
is not strongly contextual.
\end{proposition}

\begin{proof}
The variables are measurement--outcome coordinates, and the scope relation at
$C$ is exactly $e_C$. The canonical diagram leaves proper subscopes
unconstrained, while the support packaging records their projected supports.
Every global assignment satisfying the maximal supports automatically satisfies
those projections, and conversely maximal-context validity is unchanged. Thus
the two global-family spaces coincide. Apply
Theorem~\ref{thm:distributed-equivalence}.
\end{proof}

\begin{remark}
The result represents the established global-section criterion in the same typed
assignment spaces used for effectivity and resources. The restriction-stability
proved in Proposition~\ref{prop:empirical-subfunctor} is internal to a fixed empirical model; functoriality of the
construction $e\mapsto\mathrm{Pack}(e)$ on morphisms of empirical models remains
Conjecture~\ref{conj:context}.
\end{remark}

\subsection{Acyclic gluing and finite solution count}

\begin{theorem}[Join-tree gluing]\label{thm:jointree}
Let $V$ be a finite tree vertex set, let every $C_i$ be a finite context, and
let every outcome set $O_x$ be nonempty. Assume the running-intersection
property: for every variable $x$, the vertices $\{i:x\in C_i\}$ induce a
connected subtree. Let $R_i\subseteq\prod_{x\in C_i}O_x$ be nonempty finite
relations and define
\[
\Join_iR_i=
\left\{t\in\prod_{x\in\bigcup_iC_i}O_x:
 t|_{C_i}\in R_i\ \forall i\right\}.
\]
If for every tree edge $i{-}j$,
\[
\operatorname{proj}_{C_i\cap C_j}(R_i)
=
\operatorname{proj}_{C_i\cap C_j}(R_j),
\]
then the natural join $\Join_iR_i$ is nonempty.
\end{theorem}

\begin{proof}
Root the tree and process vertices in root-to-leaf order. Choose a tuple in the
root relation. Suppose a tuple has been chosen at a parent $i$ and let $j$ be a
child. Its value on
$C_i\cap C_j$ lies in the parent projection and therefore, by projection
equality, in the child projection. Choose a tuple in $R_j$ with that separator
value. Continue recursively through the finite tree. The induction invariant is that
the union of tuples chosen at processed vertices is a well-defined partial
assignment. When a child $j$ is added, any variable shared with an earlier
processed vertex lies, by running intersection, in the parent separator
$C_i\cap C_j$, where agreement was enforced.

If a variable occurs in two chosen contexts, the running-intersection property
places it in every context and separator on the path between them. The recursive
choices therefore give it the same value throughout. The chosen tuples hence
have a well-defined union, and that union restricts into every $R_i$.
\end{proof}

The theorem is the positive acyclic counterpart of the parity cycle and is the
local-to-global principle behind join-tree database schemes \cite{beerifagin}.

\begin{theorem}[Finite solution count]\label{thm:finite-solution-count}
Suppose a finite category has passed the well-formedness check, and all objects,
arrows, assignment spaces, admissibility predicates, equality tests, and
restriction maps are given by finite effective tables satisfying identity,
composition, and functor laws. Define
\[
N(P)=\#\{s\in\textstyle\prod_c\mathcal A(c):
 s_c\in\LocalAdm(c)\ \forall c,
 \ \mathcal A(f)(s_c)=s_{c'}\ \forall f:c\to c'\}.
\]
Call $N(P)$ the \emph{finite solution count}. It is computable, and a
well-formed problem has obstruction
$\mathsf D$ exactly when every $\LocalAdm(c)$ is nonempty and $N(P)=0$.
\end{theorem}

\begin{proof}
Enumerate the finite product of assignment spaces. The supplied decision tables
determine local admissibility and every compatibility equation, so counting the
families that pass all tests computes $N(P)$. Local non-emptiness excludes
$\mathsf L$, and $N(P)=0$ is precisely absence of a compatible admissible family.
\end{proof}

The count is presentation-dependent and is not called an invariant. Under an
explicit relational encoding, existence is a finite CSP problem and counting is
a finite counting-CSP problem.

\subsection{Application: finite source integration and counting}
\label{subsec:source-integration}

Finite source integration provides a concrete tabular instance of the
compatibility diagnostic. A source-integration instance consists of a finite
attribute set $X$, nonempty finite domains $(O_x)_{x\in X}$, and local sources
$(C_i,R_i)_{i=1}^m$ with $C_i\subseteq X$ and
\[
R_i\subseteq\prod_{x\in C_i}O_x.
\]
A global record is an element $t\in\prod_{x\in X}O_x$, and it integrates the
sources when $t|_{C_i}\in R_i$ for every $i$. The integration count is
\[
N_{\mathrm{int}}=\#\{t\in\prod_{x\in X}O_x:t|_{C_i}\in R_i\text{ for all }i\}.
\]

The source relations are the local tables. An integrating record is a global row
whose projections lie in all of those tables. The $\mathsf D$ outcome captures
the familiar situation in which each source table is nonempty but the natural
join is empty.

\begin{proposition}[Source integration as a distributed selector]
\label{prop:source-integration-distributed}
Every finite source-integration instance has a canonical distributed-selector
presentation whose valid distributed strategies are exactly the integrating
global records.
\end{proposition}

\begin{proof}
Use one distributed agent for each attribute $x\in X$, give it singleton
observation set $\{\ast\}$ and choice set $O_x$, and create one world $w_i$ for
each source. The active scope at $w_i$ is $C_i$, and the world relation is
$R_i$ under the evident coordinate relabelling. A distributed strategy is a
global record $(a_x)_{x\in X}$, and validity at $w_i$ is exactly
$(a_x)_{x\in C_i}\in R_i$. Theorem~\ref{thm:distributed-equivalence} identifies
these strategies with compatible admissible families of the canonical incidence
diagram.
\end{proof}

\begin{corollary}[Source-integration diagnostics]
\label{cor:source-integration-diagnostics}
For the canonical integration presentation at tolerance zero, empty realised
source relations give outcome $\mathsf L_0$; nonempty source relations with no
integrating global record give outcome $\mathsf D_0$; and a global record gives
set-theoretic capability. Effective and bounded diagnostics are obtained by
intersecting the global-record space with the declared realisation images.
\end{corollary}

\begin{proof}
This is the diagnostic partition applied to the distributed-selector
presentation of Proposition~\ref{prop:source-integration-distributed}.
\end{proof}

\begin{theorem}[Counting source integrations is $\#\mathrm P$-complete]
\label{thm:source-integration-sharpP}
For explicitly tabulated finite source-integration instances, the function
$N_{\mathrm{int}}$ is $\#\mathrm P$-complete. Hardness holds with Boolean
domains, source contexts of size at most three, and nonempty source relations.
\end{theorem}

\begin{proof}
Membership in $\#\mathrm P$ follows by nondeterministically guessing one value
for each attribute and checking all source tables in polynomial time.

For hardness, reduce parsimoniously from $\#3\mathrm{SAT}$. Given a 3CNF formula,
use its variables as Boolean attributes. For each clause, make a source whose
context is the set of variables appearing in the clause and whose relation is
the set of satisfying assignments to those variables. Each relation is nonempty
and has arity at most three. A global record satisfies all sources exactly when
the corresponding truth assignment satisfies all clauses, so the number of
integrating records is the number of satisfying assignments. Thus exact counting
is $\#\mathrm P$-hard, and membership gives completeness.
\end{proof}

\begin{corollary}[Counting compatible families]
\label{cor:counting-compatible-families}
Counting compatible admissible families in canonical incidence diagrams of
finite distributed selector problems is $\#\mathrm P$-complete, even for Boolean
choice sets and ternary realised scopes.
\end{corollary}

\begin{proof}
Theorem~\ref{thm:distributed-equivalence} gives a cardinality-preserving
identification between valid distributed strategies, CSP solutions, and
compatible admissible families. Apply
Theorem~\ref{thm:source-integration-sharpP}.
\end{proof}

\begin{remark}[Compression of source integration]
Terminal compression replaces the source-level incidence diagram by a single
object whose local set is the global join. When the join is empty, the compressed
presentation has outcome $\mathsf L_0$ regardless of whether the original source
relations were nonempty. Thus source integration instantiates the compression
boundary between local inconsistency and compatibility failure.
\end{remark}

\begin{theorem}[Affine compatibility certificate]\label{thm:affine-certificate}
Let $X$ be a finite global variable set over a fixed finite field. Suppose each
context is a subset $c\subseteq X$, its assignment space is $\mathbb F^c$,
arrows are coordinate projections, and every local admissibility relation is
specified by affine linear equations in those coordinates. Assume compatibility is ordinary agreement on overlaps, equivalently that a
compatible family is induced by a single global assignment $x\in\mathbb F^X$.
Pad each local matrix with zero columns outside its context and stack the
equations into one global system $Ax=b$. Such a compatible global assignment
exists iff
\[
\operatorname{rank}(A)=\operatorname{rank}([A\mid b]).
\]
Consequently compatibility on this subclass is decidable in polynomial time by
Gaussian elimination.
\end{theorem}

\begin{proof}
A global assignment is compatible and locally admissible exactly when it
satisfies every local affine equation, hence exactly when it solves $Ax=b$.
The standard rank criterion characterises consistency of a finite linear system,
and Gaussian elimination computes both ranks in polynomial time.
\end{proof}


\begin{proposition}[Affine cokernel obstruction]\label{prop:affine-cokernel}
In the setting of Theorem~\ref{thm:affine-certificate}, regard the stacked
matrix as a linear map
\[
A:\mathbb F^X\longrightarrow\mathbb F^m
\]
and put $H_A^1=\operatorname{coker}(A)$. Then
\[
Ax=b\text{ is solvable}
\iff [b]=0\text{ in }H_A^1.
\]
If every local affine relation is nonempty, the canonical compatibility
presentation has outcome $\mathsf D_0$ exactly when $[b]\ne0$.
\end{proposition}

\begin{proof}
The class $[b]$ vanishes in the cokernel precisely when
$b\in\operatorname{im}(A)$, which is equivalent to the existence of
$x\in\mathbb F^X$ with $Ax=b$. By Theorem~\ref{thm:affine-certificate}, such an
$x$ is exactly a compatible globally admissible assignment. Under the additional
local-nonemptiness hypothesis, failure of this global equation is
$\mathsf D_0$ rather than $\mathsf L_0$.
\end{proof}

\begin{theorem}[Bijective-network holonomy]
\label{thm:network-holonomy}
Let $G=(X,E)$ be a finite connected graph. Give each vertex $x$ a nonempty
finite domain $C_x$. For every oriented edge $e:u\to v$, let its constraint be
the graph of a bijection $\rho_e:C_u\to C_v$, with
$\rho_{\bar e}=\rho_e^{-1}$. Fix a base vertex $x_0$ and a spanning tree $K$.
For each vertex $v$, let $\theta_v:C_{x_0}\to C_v$ be transport along the
unique path in $K$. For every oriented non-tree edge $e:u\to v$, put
\[
h_e=\theta_v^{-1}\circ\rho_e\circ\theta_u
\in\operatorname{Aut}(C_{x_0}).
\]
Then the global solution set is canonically in bijection with
\[
\bigcap_{e\in E\setminus K}\operatorname{Fix}(h_e),
\]
with the empty intersection understood as $C_{x_0}$.
Equivalently, it is the common fixed-point set of all closed-walk holonomies
based at $x_0$.
\end{theorem}

\begin{proof}
A value $c_0\in C_{x_0}$ determines values
$c_v=\theta_v(c_0)$ satisfying every tree-edge constraint. A non-tree edge
$e:u\to v$ is satisfied exactly when
\[
\rho_e(\theta_u(c_0))=\theta_v(c_0),
\]
which is equivalent to $h_e(c_0)=c_0$. Hence common fixed points give, and are
given by, global solutions. The fundamental closed walks determined by the
non-tree edges generate every closed walk, and transport respects concatenation
and reversal, giving the equivalent formulation.
\end{proof}

\begin{corollary}[Bijective-cycle fixed-point certificate]
\label{cor:cycle-holonomy}
For a cycle with edge bijections
$\rho_i:C_i\to C_{i+1}$, the global solution set is in bijection with
\[
\operatorname{Fix}(\rho_{n-1}\circ\cdots\circ\rho_0).
\]
Its canonical incidence diagram has outcome $\mathsf D_0$ exactly when this
fixed-point set is empty.
\end{corollary}

\begin{proof}
Choose the cycle with one edge deleted as the spanning tree in
Theorem~\ref{thm:network-holonomy}. The sole non-tree edge produces the displayed
cycle composite. Every edge relation and singleton projection is nonempty, so
absence of a global solution is a compatibility outcome.
\end{proof}

\begin{remark}
For binary odd parity on a cycle, each edge transports a bit either by the
identity or by complementation; an odd number of complementations gives the
fixed-point-free bit flip. The deterministic CHSH equations similarly form a
four-cycle with three equality transports and one complementation. These results
concern graphs of bijective binary constraints. General relations need not
define transports, a nonidentity permutation may have fixed points, and no
higher-arity, homotopy, or cohomological completeness statement is inferred.
\end{remark}

The cokernel class of Proposition~\ref{prop:affine-cokernel} is complete for the
stated affine subclass. General nonlinear contextual relations require
additional certificate theories beyond the affine cokernel case.

\subsection{Finite Bell coupling}

The Bell presentation uses the same local-to-global idea with probabilities.
The local data are observed distributions for different measurement settings.
Compatibility asks whether these distributions can be coupled into a single
probability measure over deterministic response tables. This is the finite local
polytope viewpoint, expressed in the same diagrammatic language as the set-valued
examples.

\begin{definition}[Distributional coupling diagram]\label{def:bellpack}
For a finite Bell scenario, let $X_i$ and $O_i$ be the settings and outcomes of
party $i$, and fix a behaviour $p=(p_{\mathbf x})_{\mathbf x}$ with
$p_{\mathbf x}\in\Delta(\prod_iO_i)$ for every setting tuple. Put
$\Omega=\prod_iO_i^{X_i}$, and for a setting tuple
$\mathbf x=(x_i)_i$ define
\[
\pi_{\mathbf x}:\Omega\to\prod_iO_i,
\qquad
\pi_{\mathbf x}(\omega)=(\omega_i(x_i))_i.
\]
Let $\mathcal C_{\mathrm{Bell}}$ be the star-shaped restriction category with a
global object $\top$, one object for each setting tuple $\mathbf x$, and arrows
$\top\to\mathbf x$. Define
\[
\mathcal A(\top)=\Delta(\Omega),
\qquad
\mathcal A(\mathbf x)=\Delta\!\left(\prod_iO_i\right),
\]
with arrow map $(\pi_{\mathbf x})_*$. Put
\[
\LocalAdm(\top)=\Delta(\Omega),
\qquad
\LocalAdm(\mathbf x)=\{p_{\mathbf x}\}.
\]
This constrained assignment diagram is a descent subfunctor exactly in the
special case where every measure at $\top$ has the prescribed marginals. Its
compatible-family space is determined by
\[
\operatorname{Coupl}(p)=
\{\mu\in\Delta(\Omega):
(\pi_{\mathbf x})_*\mu=p_{\mathbf x}
\text{ for every }\mathbf x\}.
\]
\end{definition}

\begin{theorem}[Finite local-behaviour coupling]\label{thm:bell}
For a finite Bell behaviour, the following are equivalent:
\begin{enumerate}[label=(\roman*)]
\item the behaviour has a local hidden-variable model, equivalently lies in the
finite convex hull of deterministic response tables;
\item $\operatorname{Coupl}(p)\ne\varnothing$;
\item its distributional coupling diagram has a compatible global family.
\end{enumerate}
Thus a nonlocal finite behaviour gives a distributional $\mathsf D$ obstruction
in this presentation. No nonsignalling assumption is needed for this equivalence;
signalling behaviours are simply outside the local convex hull as well.
\end{theorem}

\begin{proof}
The proof treats either standard formulation of locality.

If locality is defined by deterministic tables, then for some
$\mu\in\Delta(\Omega)$,
\[
p(\mathbf o|\mathbf x)=
\sum_{\omega\in\Omega}\mu(\omega)
\prod_i\mathbf 1_{\{o_i=\omega_i(x_i)\}},
\]
which is exactly $(\pi_{\mathbf x})_*\mu(\mathbf o)$.

For the stochastic-kernel formulation, suppose the maps
$\lambda\mapsto p_i(o|x,\lambda)$ are measurable and
\[
p(\mathbf o|\mathbf x)=
\int_\Lambda\prod_i p_i(o_i|x_i,\lambda)\,d\nu(\lambda).
\]
For each $\lambda$ define
\[
\mu_\lambda(\omega)=
\prod_i\prod_{x\in X_i}p_i(\omega_i(x)|x,\lambda),
\qquad
\mu(\omega)=\int_\Lambda\mu_\lambda(\omega)\,d\nu(\lambda).
\]
The coordinate sets are finite, and summing $\mu_\lambda$ over all tables gives
one, so $\mu\in\Delta(\Omega)$. For fixed $\mathbf x,\mathbf o$, interchange
the finite sum with the integral. Summation over unconstrained table coordinates
contributes normalized factors one, leaving
\[
(\pi_{\mathbf x})_*\mu(\mathbf o)
=\int_\Lambda\prod_i p_i(o_i|x_i,\lambda)\,d\nu(\lambda)
=p(\mathbf o|\mathbf x).
\]
Thus every stochastic local model gives a coupling.

Conversely, if $\mu\in\operatorname{Coupl}(p)$, take hidden variable
$\omega\in\Omega$ with law $\mu$ and deterministic responses
\[
p_i(o|x,\omega)=
\begin{cases}1,&o=\omega_i(x),\\0,&\text{otherwise}.
\end{cases}
\]
Their induced distributions are $(\pi_{\mathbf x})_*\mu=p_{\mathbf x}$.
Finally, a compatible family in the star diagram is a measure at $\top$ whose
pushforward at every setting object is $p_{\mathbf x}$, exactly the coupling
condition.
\end{proof}

\begin{remark}[Presentation of Bell nonlocality]
The $\mathsf D$ classification in Theorem~\ref{thm:bell} belongs to the star
presentation of Definition~\ref{def:bellpack}. In that presentation the top
object allows all probability measures on deterministic response tables, and the
setting objects prescribe the observed marginals. If the top local admissibility
set is instead replaced by $\operatorname{Coupl}(p)$, coupling failure becomes
local emptiness at that object. The invariant content is coupling feasibility,
equivalently membership in the finite local polytope, not the diagnostic letter
after arbitrary repackaging.
\end{remark}

\begin{definition}[Bell coupling matrix]\label{def:bell-matrix}
Let
\[
Q=\left\{(\mathbf x,\mathbf o):
\mathbf x\in\prod_iX_i,\ 
\mathbf o\in\prod_iO_i\right\}.
\]
For $(\mathbf x,\mathbf o)\in Q$ and $\omega\in\Omega$, define
\[
A_{(\mathbf x,\mathbf o),\omega}
=\mathbf 1_{\{\pi_{\mathbf x}(\omega)=\mathbf o\}},
\]
and regard a behaviour as the vector
$p=(p_{\mathbf x}(\mathbf o))_{(\mathbf x,\mathbf o)\in Q}$.
\end{definition}

\begin{theorem}[Bell coupling dual certificate]\label{thm:bell-dual}
For every finite Bell behaviour $p$, exactly one of the following alternatives
holds:
\begin{enumerate}[label=(\roman*)]
\item there is $\mu\in\Delta(\Omega)$ with
$(\pi_{\mathbf x})_*\mu=p_{\mathbf x}$ for every $\mathbf x$;
\item there is $\beta\in\mathbb R^Q$ such that
\[
A^{\mathsf T}\beta\ge0,
\qquad
\langle\beta,p\rangle<0.
\]
\end{enumerate}
Equivalently, coupling failure has a linear inequality certificate which is
nonnegative on every finite local behaviour and negative on $p$. If $p$ is
rational, $\beta$ may be chosen rational.
\end{theorem}

\begin{proof}
For a nonnegative vector $\mu\in\mathbb R^\Omega$,
\[
(A\mu)_{(\mathbf x,\mathbf o)}
=\sum_{\omega:\pi_{\mathbf x}(\omega)=\mathbf o}\mu(\omega)
=(\pi_{\mathbf x})_*\mu(\mathbf o).
\]
Thus the coupling equations are $A\mu=p$ with $\mu\ge0$. Summing these
equations over $\mathbf o$ for any fixed $\mathbf x$ gives
$\sum_\omega\mu(\omega)=\sum_{\mathbf o}p_{\mathbf x}(\mathbf o)=1$, so
normalisation follows automatically.

Farkas' lemma gives exactly one of
\[
\exists\mu\ge0\ (A\mu=p)
\quad\text{and}\quad
\exists\beta\ (A^{\mathsf T}\beta\ge0,
\ \langle\beta,p\rangle<0)
\]
\cite{schrijver}. If $q=A\nu$ is local, where
$\nu\in\Delta(\Omega)$, then
\[
\langle\beta,q\rangle
=\langle A^{\mathsf T}\beta,\nu\rangle\ge0.
\]
Conversely, nonnegativity on every deterministic local table is precisely the
coordinate inequality $A^{\mathsf T}\beta\ge0$. For rational $A$ and $p$, the
rational form of Farkas' lemma supplies a rational separating vector.
\end{proof}

\begin{corollary}[Finite Bell certificates]\label{cor:bell-certificates}
A finite nonlocal behaviour has a finite Bell-inequality certificate. For
rational explicit input tables, coupling feasibility is decidable by linear
programming, and an infeasible instance has a certificate verifiable by finite
rational arithmetic. The polynomial bound is measured in the encoded size of
the explicit matrix $A$; that matrix may be exponentially larger than a
compressed scenario description.
\end{corollary}

\begin{proof}
Apply Theorems~\ref{thm:bell} and~\ref{thm:bell-dual}. Standard rational linear
programming decides feasibility in time polynomial in the explicit matrix and
bit lengths, and the displayed inequalities verify the dual witness.
\end{proof}

\begin{remark}[Stochastic kernels and deterministic tables]
For finite Bell scenarios, stochastic local response kernels and the finite
convex hull of deterministic response tables define the same local polytope. The
product-table construction in the proof gives the conversion. Since the local
behaviour set is a finite-dimensional polytope, Carath\'eodory's theorem gives
an equivalent finite convex decomposition into deterministic response tables.
\end{remark}

\begin{remark}
Fine's theorem gives the celebrated equivalences for the standard
$2\times2\times2$ Bell scenario \cite{fine}. Theorem~\ref{thm:bell} uses the
standard finite local-polytope representation for arbitrary finite scenarios
\cite{pitowsky}. A CHSH behaviour violating the local bound has no global coupling
\cite{chsh}.
\end{remark}

\begin{remark}[Support versus probability]
Strong contextuality excludes every deterministic global assignment from the
support. Bell nonlocality excludes a probability measure on deterministic
response tables with the required marginals. A behaviour may have support-level
global assignments and nevertheless have no coupling reproducing its
probabilities.
\end{remark}

\section{Effectivity, computability, and resources}\label{sec:effect}

The higher layers of the diagnostic distinguish mathematical existence from
realisation. A compatible valid section may exist in the set-theoretic sense and
still be unavailable to any implementation in the declared class. Abstract
effectivity is therefore relative to a realisation map. Ordinary computability is
the special case in which the domains carry effective representations and the
realisation image is a class of computable maps. The central distinction is
between nonemptiness of the valid-section space and the existence of a uniform
procedure selecting a valid section.

\begin{definition}[Uniform binary fibre selection]
\label{def:fibsel-weihrauch}
Throughout the Weihrauch statements in this section, finite spaces carry their
standard discrete representations, countable products carry the product
representation, and closed subsets of Cantor space carry negative information.
The symbol $\mathsf{WKL}$ denotes path selection for nonempty infinite
prefix-closed binary trees under the corresponding negative tree representation,
equivalently closed choice on nonempty effectively closed subsets of Cantor
space. Give a closed subset of the discrete space $2=\{0,1\}$ by negative
information, namely by an enumeration of its complement. The multivalued
problem $\mathrm{FibSel}_2$ takes a sequence $(A_n)_{n\in\N}$ of nonempty closed
subsets $A_n\subseteq2$, uniformly given by negative information, and returns
any $f\in2^\N$ satisfying
\[
f(n)\in A_n\qquad\text{for every }n.
\]
The domain representation includes the promise that no $A_n$ is empty.
\end{definition}

\begin{theorem}[Weihrauch degree of uniform binary fibre selection]
\label{thm:fibsel-weihrauch}
With the representation of Definition~\ref{def:fibsel-weihrauch},
\[
\mathrm{FibSel}_2
\equiv_{\mathrm W}
\widehat{\mathsf C_2}
\equiv_{\mathrm W}
\mathsf C_{2^\N}
\equiv_{\mathrm W}
\mathsf{WKL},
\]
where $\mathsf C_X$ denotes closed choice on the represented space $X$ and the
hat denotes countable parallelisation.
\end{theorem}

\begin{proof}
A negatively named nonempty subset of $2$ is precisely an instance of
$\mathsf C_2$. The input to $\mathrm{FibSel}_2$ is a sequence of such instances,
and its output selects one point from every instance. Thus
$\mathrm{FibSel}_2$ is, under the stated product representation, exactly the
countable parallelisation $\widehat{\mathsf C_2}$.

The standard represented-space equivalences
\[
\widehat{\mathsf C_2}
\equiv_{\mathrm W}\mathsf C_{2^\N}
\equiv_{\mathrm W}\mathsf{WKL}
\]
are the parallel binary-choice and closed-choice formulation of
Brattka and Gherardi \cite{brattkagherardi}. They are nontrivial:
an arbitrary closed subset of Cantor space need not be a coordinate product.
Applying that cited theorem to the exact parallel-choice problem identified in
the first paragraph proves the displayed chain.
\end{proof}

\begin{remark}[Scope of the classification]
Theorem~\ref{thm:fibsel-weihrauch} classifies the input-uniform operator on
negatively represented binary fibres. It does not classify arbitrary validity
representations or countably valued fibres. The halting-stage family of
Definition~\ref{def:Kgame} is a single computable instance of the same selection
pattern. Its solution set has the PA-degree spectrum of
Corollary~\ref{cor:PAdegree}, but this nonuniform degree statement is different
from Weihrauch-completeness of the input-uniform problem.
\end{remark}

\begin{definition}[Represented realisation classes]\label{def:represented-realisation}
Fix representations of the implementation and policy spaces and a realisation
map
\[
\realise_{\mathrm{eff}}:\mathrm{Impl}_{\mathrm{eff}}
\to\Pi_{\mathrm{all}}.
\]
The effective and bounded policy classes are its full image and the image of its
restriction to $\mathrm{Impl}_{\mathcal B}$.
\end{definition}

\begin{definition}[Cost-bounded implementation classes]\label{def:cost}
A quantitative resource model consists of a cost map
\[
\operatorname{cost}:\mathrm{Impl}_{\mathrm{eff}}(G)\times W
\to\overline{\N}=\N\cup\{\infty\}.
\]
For an input-dependent bound $B:W\to\overline{\N}$ define
\[
\mathrm{Impl}_{\le B}
=\{\iota:\forall w\in W,
\operatorname{cost}(\iota,w)\le B(w)\}.
\]
Bounds are ordered pointwise; $B\le B'$ implies
$\mathrm{Impl}_{\le B}\subseteq\mathrm{Impl}_{\le B'}$. A policy is bounded
when it has at least one realising implementation in $\mathrm{Impl}_{\le B}$.
\end{definition}

\begin{remark}[Dependence on the implementation model]\label{rem:abstractimage}
Set-theoretically, any $P\subseteq\Pi_{\mathrm{all}}$ can be presented as an
image by taking implementation set $P$ and the inclusion map. Thus ``effective''
has mathematical content only after representations and an implementation model
are fixed. The computability results in this section use standard acceptable numberings.
\end{remark}

\begin{proposition}[External-policy special case]\label{prop:nonembedded}
Let $\Pi_{\mathrm{ext}}\subseteq\Pi_{\mathrm{all}}$ be a specified external
policy class. If
\[
\im\!\left(
\realise_{\mathrm{eff}}|_{\mathrm{Impl}_{\mathcal B}}
\right)=\Pi_{\mathrm{ext}},
\]
then $\Pi_{\mathcal B}=\Pi_{\mathrm{ext}}$, and the selector results reduce to
the corresponding external-policy problem.
\end{proposition}

\begin{proof}
The displayed equality is the definition of the bounded realised policy class.
All selector statements depend on implementations only through that image.
\end{proof}

\subsection{A recursive co-c.e. effectivity obstruction}

The first effectivity example has no local failure and no compatibility failure.
There is a valid selector. The problem is that every computable candidate can be
turned against itself at its own index. This is the standard diagonally
nonrecursive construction, expressed as a selector game.

Fix an admissible numbering $(\varphi_n)$ of all partial binary-valued recursive
functions, so every total computable binary function has an index in the
numbering \cite{rogers}. Let $\operatorname{Comp}(n,s,b)$ be a fixed
primitive-recursive predicate saying that $\varphi_n(n)$ halts within $s$ steps
with output $b\in\{0,1\}$.

\begin{definition}[Halting-stage selector game]\label{def:Kgame}
Take
\[
W=\N\times\N,\quad Z=\N,\quad I(n,s)=n,
\quad R=M=H=Y=\{0,1\},\quad A=\{\ast\},
\]
with $D,q,e$ the identity maps and zero--one loss. Define the total recursive
dynamics
\[
U((n,s),m,\ast)=
\begin{cases}
1-m,&\operatorname{Comp}(n,s,m),\\
m,&\text{otherwise}.
\end{cases}
\]
Let $\Pi_{\mathrm{eff}}$ be the total computable binary-valued policies.
\end{definition}

\begin{lemma}[Uniform co-c.e. validity]\label{lem:Kfibres}
Under the standard pairing code for $(n,r)$,
\[
(r,\ast)\in\Adm(n)
\iff
\forall s\,\neg\operatorname{Comp}(n,s,r).
\]
Hence the complement of
$\{(n,r):(r,\ast)\in\Adm(n)\}$ is computably enumerable uniformly in
$(n,r)$, and every fibre is nonempty.
\end{lemma}

\begin{proof}
If a computation witness with output $r$ exists, the corresponding witness world
makes the dynamics return $1-r$, so report $r$ is invalid. If no such witness
exists, the dynamics returns $r$ on every world in the fibre. The complement is
therefore defined by the existential primitive-recursive predicate
$\exists s\,\operatorname{Comp}(n,s,r)$. A partial function has at most one
binary output, so at most one report is excluded.
\end{proof}

\begin{theorem}[Co-c.e. effectivity obstruction]\label{thm:Kgame}
The halting-stage selector game has a set-theoretic valid selector but no total
computable valid selector. Its semantic obstruction is $\mathsf K$.
\end{theorem}

\begin{proof}
Define set-theoretically
\[
f(n)=
\begin{cases}
1-b,&\varphi_n(n)\downarrow=b\in\{0,1\},\\
0,&\varphi_n(n)\uparrow.
\end{cases}
\]
If the computation diverges, both reports are valid; if it halts with value
$b$, Lemma~\ref{lem:Kfibres} excludes only $b$, so $1-b$ is valid. Thus $f$ is
a valid selector. This explicit definition uses classical excluded middle for
the halting statement but no choice principle.

Suppose a total computable valid selector $g$ existed. Choose an index $k$ in the
fixed binary numbering with $g=\varphi_k$, and put $b=g(k)$. A computation
witness shows $b\notin\Adm(k)$ by Lemma~\ref{lem:Kfibres}, contradicting
validity. Local sets and the set-theoretic global family are nonempty, so the
first semantic failure is effectivity.
\end{proof}

\begin{theorem}[$\Pi^0_1$ selector class]\label{thm:Pi01}
Let
\[
\mathcal S=\{g\in2^{\N}:\forall n\ (g(n),\ast)\in\Adm(n)\}.
\]
Then $\mathcal S$ is a nonempty $\Pi^0_1$ class with no computable member.
\end{theorem}

\begin{proof}
The complement of $\mathcal S$ is the c.e. open set
\[
\bigcup_{\operatorname{Comp}(n,s,b)}
\{g\in2^{\N}:g(n)=b\}.
\]
Equivalently, $\mathcal S$ is the path set of the computable tree
\[
T=\left\{\sigma\in2^{<\N}:
\forall n<|\sigma|\;\forall s<|\sigma|,
\neg\operatorname{Comp}(n,s,\sigma(n))\right\}.
\]
Hence $\mathcal S$ is effectively closed. The explicit selector constructed in
Theorem~\ref{thm:Kgame} proves nonemptiness. If a computable $g$ had index $k$,
a computation witness for $\varphi_k(k)=g(k)$ would violate membership at
$n=k$.
\end{proof}

\begin{corollary}[Oracle degrees of capability]\label{cor:PAdegree}
For the fixed acceptable numbering $(\varphi_n)$ introduced in Section~\ref{sec:effect}, let $\Pi_X$ be the policies computable relative
to oracle $X$. The halting-stage selector game is capable in $\Pi_X$ exactly
when $X$ has PA degree, equivalently when $X$ computes a binary diagonally
nonrecursive function. The resulting oracle-degree class is invariant under
replacement by another acceptable numbering. In
particular, the game has a low noncomputable selector and a selector of
hyperimmune-free degree.
\end{corollary}

\begin{proof}
The class $\mathcal S$ is exactly the class of binary diagonally nonrecursive
functions for the fixed acceptable numbering. The degrees computing such a
function are the PA degrees for every acceptable numbering, and the Low Basis and Hyperimmune-Free Basis Theorems supply low and hyperimmune-free
members of every nonempty $\Pi^0_1$ class
\cite{jockusch1989,jockuschsoare,rogers}.
\end{proof}

The index argument is not the local uniform diagonal of
Definition~\ref{def:diag}: no admissibility fibre is empty.

\subsection{Nested distributed effectivity}

The finite distributed theory separates local and compatibility failure. The
following countable finitary extension supplies an overlapping-scope
effectivity obstruction.

\begin{definition}[Countable finitary distributed selector]
\label{def:countable-distributed}
A countable finitary distributed selector problem has the data of
Definition~\ref{def:distributed}, except that $W$ and $J$ may be countable,
every active scope $S(w)$ is finite, and all carriers and relations have fixed
effective representations when computability is discussed.
\end{definition}

\begin{definition}[Nested-prefix distributed selector]
\label{def:nested-prefix}
Let $T\subseteq2^{<\N}$ be a prefix-closed infinite binary tree. Define a
countable finitary distributed selector $P_T$ by
\[
J=W=\N,\qquad Z_j=\{\ast\},\qquad C_j=\{0,1\},
\]
\[
S(n)=D_n=\{0,\ldots,n\},
\qquad
V_n=T\cap2^{n+1}.
\]
Its reduced prefix diagram has objects $D_n$, arrows $D_n\to D_m$ for
$m\le n$, assignment sets $2^{n+1}$, local sets $V_n$, and prefix-restriction
maps. The full canonical incidence category additionally contains every
nonempty finite subset of a realised prefix. When $T$ is computable, take
$\mathrm{Impl}_{\mathrm{eff}}$ to be indices of total computable binary
functions, use the standard evaluation realisation map into $2^\N$, and define
the effective global-family space by intersecting its image with $[T]$.
\end{definition}

\begin{theorem}[Nested distributed path representation]
\label{thm:nested-distributed-K}
For every prefix-closed infinite binary tree $T$:
\begin{enumerate}[label=(\roman*)]
\item every local set in the reduced prefix diagram is finite and nonempty;
\item the set-theoretic global-family space and the valid distributed-strategy
space are both canonically $[T]$;
\item the reduced prefix category has no initial object.
\end{enumerate}
If membership in $T$ is decidable, then $[T]$ is a nonempty effectively closed
subset of Cantor space, hence a $\Pi^0_1$ class. If it has no computable member,
computable strategies as the effective image give outcome $\mathsf K_0$. For the
binary DNR tree of Theorem~\ref{thm:Pi01}, an oracle computes a global strategy
exactly when it has PA degree.
\end{theorem}

\begin{proof}
Since $T$ is infinite, finitely branching, and prefix closed, König's lemma gives
a path through $T$; in particular, $T$ has nodes of every finite length, and
every $V_n$ is nonempty. A distributed strategy is a binary
sequence $x$, and validity at world $n$ is precisely
$x|_{n+1}\in T$. Thus validity for every world is equivalent to $x\in[T]$.
Likewise, a compatible family $(\sigma_n)_n$ in the reduced prefix diagram
satisfies $\sigma_m=\sigma_n|_{m+1}$ for $m\le n$, so its union is a path in
$T$; restrictions of a path give the inverse construction.

An initial object would require
some finite prefix $D_n$ to have an arrow to every $D_m$, which would require
$n\ge m$ for every $m$, impossible. If membership in $T$ is computable, its
path set is effectively closed; absence of a computable path then gives the
$\mathsf K_0$ classification. For the stated DNR tree,
Corollary~\ref{cor:PAdegree} gives the final oracle characterisation.
\end{proof}

\begin{corollary}[Pointwise versus uniform prefix sections]
\label{cor:nested-nonuniform}
If $T$ has no computable path, there is a matching family on the reduced prefix
diagram such that every finite component has a computable name, but no single
computable procedure produces these names uniformly in $n$.
\end{corollary}

\begin{proof}
Choose a path $x\in[T]$ set-theoretically and take
$\sigma_n=x|_{n+1}$. Every finite string $\sigma_n$ is computable, and the
family is matching. A uniform procedure computing $\sigma_n$ from $n$ would
compute their union $x$, contrary to the absence of computable paths.
\end{proof}

The preceding effective-outcome classification concerns a fixed tree with
decidable membership. In Theorem~\ref{thm:nested-WKL}, the tree itself is
represented input and may be noncomputable.

\begin{theorem}[Uniform degree of nested-prefix selection]
\label{thm:nested-WKL}
Let the input be a nonempty path set represented by a prefix-closed binary tree
$T$ given by negative information for its complement, equivalently by a negative
name for the closed set $[T]\subseteq2^\N$. The multivalued operator taking $T$
to a valid strategy of $P_T$ is Weihrauch-equivalent to $\mathsf{WKL}$. The
input need not be computable.
\end{theorem}

\begin{proof}
By Theorem~\ref{thm:nested-distributed-K}(ii), the valid strategies of $P_T$ are
exactly the paths through the input tree $T$, and both translations are the
identity on the tree code and on the output binary sequence. The represented
problem is therefore the standard path-selection formulation of
$\mathsf{WKL}$.
\end{proof}

\begin{remark}
The full canonical incidence category of $P_T$ is generally not a chain, and adjacent
projection equality can fail when $T$ has dead ends; the finite join-tree
theorem therefore does not apply. The exact Weihrauch statement concerns the input-uniform
operator $T\mapsto[T]$, not any fixed DNR instance.
\end{remark}


\subsection{Computational-class and resource separations}

\begin{theorem}[Primitive-recursive class separation]\label{thm:primitive-class-separation}
Let $a:\N\to\N$ be a total recursive function that is not primitive recursive,
for example a standard diagonal Ackermann function \cite{rogers}. There is a
passive recursive
game whose unique valid principal selector has report component $a$. With total recursive policies as
$\Pi_{\mathrm{eff}}$ and primitive-recursive policies as $\Pi_{\mathcal B}$,
the bounded-class obstruction is $\mathsf R$ relative to that declared class.
This is a computational-class separation rather than a quantitative time or
space lower bound. The presentation here is recursive, not an instance of the
primitive-recursive syntax $\mathcal L_{\mathrm{PR}}$ used in
Theorem~\ref{thm:AH}.
\end{theorem}

\begin{proof}
Take $W=Z=R=M=H=Y=\N$, $I=D=e=q=\mathrm{id}$, singleton actions, zero--one
loss, and passive total recursive dynamics
$U(n,m,\ast)=a(n)$. At interface value $n$, exact validity is
$r=a(n)$, so $\Adm(n)=\{(a(n),\ast)\}$. Consequently every valid selector has
report component $a$. The function $a$ is total recursive, hence gives an
effective admissible section, but by hypothesis no primitive-recursive policy
can realise it. Thus $\Gamma_{\mathrm{eff}}\ne\varnothing$ and
$\Gamma_{\mathcal B}=\varnothing$.
\end{proof}

\begin{definition}[Deterministic-time policy class]\label{def:dtime-class}
For the binary-string specialisation and a time-constructible bound $t$, let
implementations be deterministic Turing machines, let
$\operatorname{cost}(\iota,x)$ be running time, put
$B_c(x)=c\,t(|x|)$ for $c\ge1$, and define
\[
\Pi_{\mathrm{DTIME}(t)}
=\bigcup_{c\ge1}
\im\!\left(
\realise_{\mathrm{eff}}|_{\mathrm{Impl}_{\le B_c}}
\right).
\]
\end{definition}

\begin{theorem}[Deterministic-time resource separation]\label{thm:timeR}
Let $t,g:\N\to\N$ be time-constructible bounds with $t(n)\ge n$ and
$t(n)\log t(n)=o(g(n))$. There is a passive exact game on binary-string inputs
with a computable valid selector but no valid selector in
$\Pi_{\mathrm{DTIME}(t)}$.
\end{theorem}

\begin{proof}
By the deterministic time-hierarchy theorem, since
$c\,t(n)\log(c\,t(n))=o(g(n))$ for every fixed $c$, choose a language
$L\in\mathrm{DTIME}(g)\setminus\bigcup_{c\ge1}\mathrm{DTIME}(c\,t)$
\cite{arorabarak}. Take
\[
W=Z=\{0,1\}^*,\qquad I=\mathrm{id},\qquad
R=H=Y=\{0,1\},
\]
with singleton message and action sets, identity decoder, zero--one loss, and
passive queried value $q(U_0(x))=\chi_L(x)$. The unique valid report selector is
$\chi_L$. Take $\Pi_{\mathrm{eff}}$ to be the total computable report policies
and $\Pi_{\mathcal B}=\Pi_{\mathrm{DTIME}(t)}$ from
Definition~\ref{def:dtime-class}. A valid bounded report policy would decide
$L$ in $\mathrm{DTIME}(t)$, contradicting the choice of $L$. Thus an effective
section exists but no section lies in the declared time class.
\end{proof}

\begin{proposition}[Degrees of the computability examples]\label{prop:example-degrees}
For the halting-stage selector game with zero--one loss,
\[
d_{\mathrm{pt}}=d_{\mathrm{obs}}=0,
\qquad
d_{\mathrm{eff}}=1.
\]
For the deterministic-time hierarchy game with
$\Pi_{\mathcal B}=\mathrm{DTIME}(t)$,
\[
d_{\mathrm{pt}}=d_{\mathrm{obs}}=d_{\mathrm{eff}}=0,
\qquad d_{\mathcal B}=1.
\]
\end{proposition}

\begin{proof}
The explicit set-theoretic selector for the halting-stage selector game gives
$d_{\mathrm{obs}}=0$, and world-dependent choices give $d_{\mathrm{pt}}=0$.
Every computable policy fails on at least one world, so its worst-case zero--one
loss is $1$; the constant-zero policy is computable, so the class is nonempty and
$d_{\mathrm{eff}}=1$.

For the time-hierarchy game, the characteristic function of $L$ is an effective
zero-loss selector, so the first three degrees vanish. Every policy in
$\mathrm{DTIME}(t)$ fails on some input, hence has worst-case loss $1$; constant
machines belong to the class, so it is nonempty. Thus $d_{\mathcal B}=1$.
\end{proof}

\subsection{Arithmetic complexity of recursive capability}

\begin{definition}[Concrete primitive-recursive game language]\label{def:AHdef}
Let $\mathcal L_{\mathrm{PR}}$ be a fixed effective syntax of well-typed game
descriptions. A description contains a finite list of carrier sorts, each equal
to $\N$ or to a fixed finite tagged set, together with indices from a standard
effective enumeration of primitive-recursive structural maps. Every index has
explicit source and target sort annotations, and the grammar permits only
sort-compatible composition and application. Syntactic well-formedness is
decidable. The index set $\mathrm{CAP}_{\mathrm{PR}}$ is viewed as a subset of
$\N$ by declaring malformed strings false; equivalently, it is the subset of
valid syntax codes satisfying the capability predicate. Exact zero--one loss uses
decidable equality. A policy is coded by a fixed primitive-recursive pairing of a
pair of ordinary partial-recursive program indices, and membership in
$\mathrm{CAP}_{\mathrm{PR}}$ requires both policy components to be total on their
entire declared domains and to form an exact valid selector. Product carriers
such as $\N\times\N_{\ge1}$ are represented in the $\N$ sort by a fixed
primitive-recursive pairing, with positive coordinates encoded by a successor.
Totality on counterfactual report--action pairs is part of this language;
requiring action totality only on the realised report graph would define a
different index set. The universal variable tags and all finite tuples of
computation stages are encoded by fixed primitive-recursive pairing functions.
\end{definition}

\begin{theorem}[Recursive capability is $\Sigma^0_3$-complete]\label{thm:AH}
The index set $\mathrm{CAP}_{\mathrm{PR}}$ is $\Sigma^0_3$-complete.
\end{theorem}

\begin{proof}
\emph{Membership.} Encode the report and action programs by one index $p$ and
let the universal variable $x$ carry one of three tags:
\begin{enumerate}[label=(\alph*)]
\item $x=\langle0,z\rangle$ asks for a stage at which the report program halts
on $z$;
\item $x=\langle1,z,r\rangle$ asks for a stage at which the action program
halts on $(z,r)$;
\item $x=\langle2,w\rangle$ asks for stages computing the selected report at
$I(w)$ and the corresponding action, followed by verification of the exact
validity equality at $w$.
\end{enumerate}
The well-formedness predicate for $G$ is decidable. Primitive-recursive
structural maps are total by syntax, and their values and the final equality are
decidable once the policy computations have converged. A single stage code $s$
can pair all finite computation witnesses needed by the tag. Hence a recursive
predicate $R$ satisfies
\[
G\in\mathrm{CAP}_{\mathrm{PR}}
\iff
\mathsf{WF}(G)\land\exists p\,\forall x\,\exists s\ R(G,p,x,s).
\]
Since $\mathsf{WF}$ is decidable, this places the index set in $\Sigma^0_3$.

\emph{Hardness.} Let $W_k$ be the $k$th c.e. set. The set
$\mathrm{REC}=\{k:W_k\text{ is recursive}\}$ is $\Sigma^0_3$-complete
\cite{soare}. Let $E(k,n,t)$ be the primitive-recursive predicate saying that
$n$ first enters $W_k$ at stage $t\ge1$. Uniformly from $k$, construct
\[
W=\N\times\N_{\ge1},\quad Z=\N,\quad I(n,t)=n,
\quad R=M=\N,\quad A=\{\ast\},\quad H=Y=\{0,1\},
\]
with $D(r)=r$, $q=\mathrm{id}$, constant decoder $e(r)=0$, and
\[
\ell(r,y)=
\begin{cases}0,&e(r)=y,\\1,&e(r)\ne y.\end{cases}
\]
Define the primitive-recursive dynamics
\[
U_k((n,t),r,\ast)=
\begin{cases}
1,&r=0\text{ and }E(k,n,t),\\
1,&r>0\text{ and }\neg E(k,n,r),\\
0,&\text{otherwise}.
\end{cases}
\]
Report $0$ is valid throughout the fibre over $n$ exactly when
$n\notin W_k$; a positive report $r$ is valid exactly when $E(k,n,r)$. Thus
\[
\Adm_k(n)=
\begin{cases}
\{(0,\ast)\},&n\notin W_k,\\
\{(r,\ast):r>0\land E(k,n,r)\},&n\in W_k.
\end{cases}
\]
If $W_k$ is recursive, decide membership, output $0$ on nonmembers, and search
for the unique positive first-entry stage on members. Conversely, a total
computable selector $\rho$ yields the total characteristic function
\[
\chi_{W_k}(n)=
\begin{cases}1,&\rho(n)>0,\\0,&\rho(n)=0,\end{cases}
\]
so $W_k$ is recursive. The map from $k$ to the displayed tuple of
primitive-recursive structural indices is computable. Therefore
$k\in\mathrm{REC}$ iff $G(k)\in\mathrm{CAP}_{\mathrm{PR}}$, proving many-one
hardness.
\end{proof}

\section{Extensive-form causal descent}\label{sec:feedback}

The causal tier adds temporal structure. The arena is fixed before the strategy
is chosen: nature moves, information sets, choice sets, and losses are part of
the input. A strategy then selects actions at the predictor information states.
In perfect-information arenas, backward induction characterises capability
because each reached decision node can be solved using the information available
there. Compatibility obstructions arise only when choices are shared across
information-limited contexts; Corollary~\ref{cor:team-decision} supplies that
imperfect-information case. Declared covers may still exhibit
presentation-relative matching failure, as active sensing shows.

\begin{definition}[Finite extensive prediction arena]\label{def:arena}
A finite extensive prediction arena consists of a finite rooted tree
$\mathcal T$ with root $\varepsilon$. Every non-root node has a unique
predecessor. Terminal nodes have no children. The nonterminal nodes are partitioned as $\mathcal T_N\sqcup \mathcal T_P$ into nature and
predictor nodes. Every nonterminal node $h$ has a finite nonempty set
$\operatorname{Ch}(h)$ of outgoing labels, with one child $h\cdot d$ for each
$d\in\operatorname{Ch}(h)$. Write
$\operatorname{Ch}(h)=M_N(h)$ at nature nodes and
$\operatorname{Ch}(h)=C_h$ at predictor nodes.

Let $W\ne\varnothing$. Define world-consistency sets recursively by
\[
W_\varepsilon=W,
\]
\[
W_{h\cdot m}=\{w\in W_h:\nu_h(w)=m\}
\quad(h\text{ nature}),
\qquad
W_{h\cdot c}=W_h
\quad(h\text{ predictor}),
\]
where $\nu_h:W_h\to M_N(h)$. At each predictor node there is a loss
$\lambda_h:W_h\times C_h\to[0,\infty]$. Empty $W_h$ are allowed; such nodes are
not reached by any world. Predictor nodes are their own information states, so
the proved arena has perfect information. Arena losses may take the value
$+\infty$, as do the static prediction-game losses of Definition~\ref{def:game}.
\end{definition}

\begin{definition}[Strategies, reachability, and play loss]\label{def:reach}
A strategy is $\pi\in\mathrm{Strat}(\mathcal T)=\prod_{h\in \mathcal T_P}C_h$. Define
$\operatorname{Reach}(w,h;\pi)$ recursively: the root is reached; a nature child
$h\cdot m$ is reached iff $h$ is reached and $\nu_h(w)=m$; a predictor child
$h\cdot c$ is reached iff $h$ is reached and $\pi(h)=c$.

For a partial strategy $s$ on a prefix-closed patch containing $h$, define
$\operatorname{Reach}(w,h;s)$ by the same recursion. Every predictor ancestor of
$h$ belongs to the patch, so all required choices are defined. A node $h$ is
$\pi$-consistent when $\pi$ selects every predictor-edge label on the unique
root-to-$h$ path.

The aggregate finite-play loss is
\[
L_{\mathcal T}(\pi,w)=
\max\left(
\{\lambda_h(w,\pi(h)):
 h\in \mathcal T_P,\ \operatorname{Reach}(w,h;\pi)\}\cup\{0\}
\right).
\]
A strategy is $\eps$-capable iff
$L_{\mathcal T}(\pi,w)\le\eps$ for every world.
\end{definition}

\begin{lemma}[Path characterisation of reachability]\label{lem:path-reach}
For a total strategy $\pi$, world $w$, and node $h$,
$\operatorname{Reach}(w,h;\pi)$ holds iff $w\in W_h$ and $\pi$ selects every
predictor-edge label on the unique root-to-$h$ path. The same equivalence holds
for a partial strategy on a prefix-closed patch containing $h$.
\end{lemma}

\begin{proof}
Induct along the unique path. At a nature edge, the recursive definition of
$W_{h\cdot m}$ adds exactly the condition $\nu_h(w)=m$. At a predictor edge,
reachability adds exactly the condition that the strategy selects its label.
Prefix closure supplies every ancestor coordinate in the partial case.
\end{proof}

\begin{proposition}[Finite causal unfolding]\label{prop:finite-unfolding}
Every strategy and world determine a unique finite play.
\end{proposition}

\begin{proof}
At a reached nature node, $\nu_h(w)$ selects one child; at a reached predictor
node, $\pi(h)$ selects one child. Finiteness and acyclicity force termination,
and deterministic selection at each node gives uniqueness.
\end{proof}

\begin{definition}[Backward solvability]\label{def:solv}
For fixed $\eps$, define $\mathrm{Good}_\eps(h)$ at a predictor node by
\[
\mathrm{Good}_\eps(h)=
\{c\in C_h:\forall w\in W_h,\ \lambda_h(w,c)\le\eps\}.
\]
Define $\mathrm{Solv}_\eps(h)$ by backward recursion. It holds at terminal nodes
and at every node with $W_h=\varnothing$. At a nonempty nature node it holds iff
it holds at every child with nonempty world-consistency set. At a nonempty
predictor node it holds iff
\[
\exists c\in\mathrm{Good}_\eps(h),
\qquad
\mathrm{Solv}_\eps(h\cdot c).
\]
\end{definition}

\begin{theorem}[Perfect-information backward induction]\label{thm:backward}
A finite perfect-information arena has an $\eps$-capable strategy iff
$\mathrm{Solv}_\eps(\varepsilon)$ holds.
\end{theorem}

\begin{proof}
Assume root solvability. Since the arena is finite, use finite choice to select,
for every solvable predictor node $h$ with nonempty $W_h$, a witness
$c_h\in\mathrm{Good}_\eps(h)$ whose child is solvable. Set $\pi(h)=c_h$ at
those nodes and choose an arbitrary element of $C_h$ at all remaining predictor
nodes. Induction down any $\pi$-consistent path shows that every
node on the path with nonempty world set is solvable. Hence each reached
predictor choice is good, and the strategy is capable.

Conversely, let $\pi$ be capable. Use backward induction on height for nodes that
are $\pi$-consistent and have nonempty $W_h$. At a nature node, each child with
nonempty world set is reached by the worlds selecting that move and is solvable
by induction. At a predictor node, Lemma~\ref{lem:path-reach} shows that every
world in $W_h$ reaches $h$; capability therefore gives
$\pi(h)\in\mathrm{Good}_\eps(h)$, and induction gives solvability of
$h\cdot\pi(h)$. The root is consequently solvable.
\end{proof}

\begin{corollary}[Node-local sufficiency in perfect-information arenas]
\label{cor:no-perfect-D}
If $\mathrm{Good}_\eps(h)\ne\varnothing$ at every predictor node with
$W_h\ne\varnothing$, then the arena is $\eps$-capable.
\end{corollary}

\begin{proof}
Backward induction on height shows every node solvable: all relevant nature
children are solvable by induction, and at a predictor node any good choice has
a solvable child by the induction hypothesis. Apply
Theorem~\ref{thm:backward}.
\end{proof}

\begin{remark}[Perfect-information case]
In a perfect-information arena, node-local nonemptiness is enough. Once every
reached predictor node has at least one choice that is good for all worlds still
consistent with that node, backward induction assembles those choices into a
global strategy. Compatibility becomes a separate obstruction only when choices
are shared across information-limited nodes or external cover objects.
\end{remark}
\begin{definition}[Backward-induction congruence]
\label{def:record-congruence}
Let $\mathcal T$ be a finite perfect-information arena. A
\emph{backward-induction congruence} is an equivalence relation $\sim$ on nodes
together with a rank $r:\mathcal T\to\N$ which strictly decreases along every
edge, such that equivalent nodes have equal rank and:
\begin{enumerate}[label=(\roman*)]
\item they have the same status: terminal, nature, or predictor;
\item their world-consistency sets are equal;
\item equivalent predictor nodes have the same choice set and identical loss
tables;
\item equivalent nature nodes have the same nature alphabet and identical
nature maps on their common world set;
\item children reached under any common outgoing label are equivalent.
\end{enumerate}
The quotient graph has one vertex per equivalence class and the induced labelled
edges.
\end{definition}

\begin{theorem}[Quotient backward induction]
\label{thm:record-quotient-bi}
For a backward-induction congruence, $\mathrm{Good}_\eps$ at predictor nodes and
$\mathrm{Solv}_\eps$ at all nodes are constant on equivalence classes. The
quotient graph is acyclic, and backward induction on it decides capability. For
explicit tables, a safe running-time bound is
\[
O\!\left(
\sum_{[h]\,\mathrm{predictor}}|W_{[h]}|\,|C_{[h]}|
+\sum_{[h]\,\mathrm{nature}}
(|W_{[h]}|+|M_N([h])|)
+|E_{\mathrm{quot}}|
\right).
\]
\end{theorem}

\begin{proof}
At equivalent predictor nodes, equality of world sets, choice sets, and loss
tables gives equality of $\mathrm{Good}_\eps$ directly. Equal rank and strict
rank decrease along edges make the quotient acyclic. Prove constancy of
$\mathrm{Solv}_\eps$ by induction on rank. Terminal status and empty-world
status agree. At equivalent nature nodes, equality of nature maps identifies
the nonempty-world children label by label, and successor congruence plus the
induction hypothesis gives the same universal solvability condition. At
equivalent predictor nodes, equality of good choices and successor classes gives
the same existential condition. Root solvability therefore can be evaluated on
the quotient and is equivalent to capability by
Theorem~\ref{thm:backward}.

Computing each predictor good set uses the first displayed sum. Nature-world
routing and nonempty-child tests are bounded by the second, and propagation
through the quotient uses each quotient edge once. No reduction factor is
asserted unless the quotient is actually smaller than the original table.
\end{proof}


Under the node-local presentation by the sets $\mathrm{Good}_\eps(h)$, there is
no compatibility-only obstruction. Cover-incidence presentations can assign a
cover-relative $\mathsf D_\eps$ label, as in
Theorem~\ref{thm:sensing-negative}.

\begin{definition}[Causal patches and partial strategies]\label{def:patch}
Let $\mathsf{Patch}(\mathcal T)$ be the poset of nonempty prefix-closed node
sets. Every patch contains the root; intersections of patches are therefore
nonempty and prefix closed. A union cover is a family $\{U_i\subseteq U\}$ with
$U=\bigcup_iU_i$. Pulling such a cover back along $V\subseteq U$ gives the cover
$V=\bigcup_i(U_i\cap V)$.

Put
\[
\mathscr S(U)=\prod_{h\in U\cap \mathcal T_P}C_h,
\]
with coordinate restriction. Because the empty lower set is omitted,
$\mathsf{Patch}(\mathcal T)$ is the basis of nonempty lower sets for the prefix
order, equivalently Alexandroff opens under the order convention in which arrows
point from larger patches to smaller patches, rather than the full open-set
lattice.
\end{definition}

\begin{lemma}[Partial-strategy sheaf]\label{lem:strategy-sheaf}
The assignment $U\mapsto\mathscr S(U)$ is a sheaf for union covers.
\end{lemma}

\begin{proof}
A matching family agrees at every predictor node occurring in two patches.
Define the glued value using any patch containing the node. Matching gives
well-definedness, and coordinate-wise construction gives existence and
uniqueness.
\end{proof}

\begin{remark}[Partial plans]
A patch is an initial piece of the decision tree, and a section over a patch is a
partial contingent plan. The sheaf condition says that compatible partial plans
are not merely consistent pairwise; they glue uniquely to a plan on the union.
Activated validity adds the operational convention that a loss obligation is
checked only at nodes actually reached by the partial plan.
\end{remark}

\begin{definition}[Activated-validity presheaf]\label{def:activated}
For a patch $U$ define
\[
\mathscr V_\eps(U)=
\left\{s\in\mathscr S(U):
\forall h\in U\cap \mathcal T_P\,\forall w\in W_h,
\operatorname{Reach}(w,h;s)
\Rightarrow
\lambda_h(w,s(h))\le\eps
\right\}.
\]
A partial strategy may choose a successor outside its patch; obligations at
nodes not activated by that choice are not evaluated. No hereditary implication
from successor loss to predecessor loss is assumed or needed.
\end{definition}

\begin{theorem}[Activated validity is a subsheaf]\label{thm:activated-sheaf}
The assignment $\mathscr V_\eps$ is a subsheaf of $\mathscr S$.
\end{theorem}

\begin{proof}
If $V\subseteq U$ and $h\in V$, every ancestor of $h$ lies in $V$, so the run
to $h$ computed from $s|_V$ is the run computed from $s$. Validity therefore
restricts. For a matching family on a union cover, let $s$ be its unique ambient
glue. If $s$ reaches $h$, choose a patch containing $h$. Prefix closure places
all ancestors there, and matching makes the local and global runs to $h$
identical. Local activated validity supplies the loss bound.
\end{proof}

\begin{corollary}[Full patch-site collapse]\label{cor:patch-collapse}
For the full causal patch category of Section~\ref{sec:feedback}, the full tree
is initial and activated validity is a subfunctor. Hence
\[
\Gamma_{\mathrm{set}}(\mathscr V_\eps)
\cong\mathscr V_\eps(\mathcal T),
\]
and full-site $\mathsf D_\eps$ is impossible.
\end{corollary}

\begin{proof}
Every patch receives the unique restriction arrow from the full tree. Apply
Theorem~\ref{thm:initial-collapse} and
Theorem~\ref{thm:activated-sheaf}.
\end{proof}

\begin{theorem}[Set-theoretic causal descent]\label{thm:causal-descent}
For $\Pi\subseteq\mathrm{Strat}(\mathcal T)$ write
$\Capability^{\mathrm{arena}}_\eps(\mathcal T,\Pi)$ when some
$\pi\in\Pi$ satisfies $L_{\mathcal T}(\pi,w)\le\eps$ for every world. For the
full patch $B=\mathcal T$,
\[
\Capability^{\mathrm{arena}}_\eps(\mathcal T,\Pi)
\iff
\Pi\cap\mathscr V_\eps(B)\ne\varnothing.
\]
For every union cover of $B$, global valid strategies are in bijection with
matching families of valid partial strategies.
\end{theorem}

\begin{proof}
The first statement expands the aggregate reached-node loss. The second is the
sheaf equalizer of Theorem~\ref{thm:activated-sheaf}.
\end{proof}

\begin{proposition}[Finite effective gluing]\label{prop:finite-effective-gluing}
In a finite effectively tabulated arena with a finite effectively indexed cover,
a uniformly computable matching family has a computable global glue.
\end{proposition}

\begin{proof}
For each predictor node, search the finite cover for the first patch containing
it and output that patch's computable value. Matching makes the result
independent of the selected patch. This finite coordinate procedure computes the
unique glue.
\end{proof}

\begin{proposition}[Finite height-one recovery]\label{prop:heightone}
Let a static selector problem have finite $\im(I)$ and finite choice set $C$.
Form an arena with a nature root, nature moves $\im(I)$, and
$\nu_\varepsilon(w)=I(w)$. Its child $h_z$ is a predictor node with choice set
$C$ and loss
\[
\lambda_{h_z}(w,c)=
\begin{cases}0,&V(w,c),\\1,&\neg V(w,c).\end{cases}
\]
Then
\[
\mathscr V_0(\mathcal T)
\cong
\{s:\im(I)\to C:s(z)\in\Adm_V(z)\ \forall z\}.
\]
\end{proposition}

\begin{proof}
The root sends exactly the worlds in $I^{-1}(z)$ to $h_z$, so zero activated
loss there is membership in $\Adm_V(z)$.
\end{proof}

\subsection{Active sensing and presentation relativity}

Active sensing is a compact case study in presentation relativity. The full
patch site contains a decision skeleton whose validity set is empty, so the same
incapability appears as local failure in that presentation. The branch cover
omits that global object. Each branch can then be solved locally by choosing an
action that leaves the branch, but the branch-local solutions cannot be matched:
matching forces one common root action before the later observations are known.

\begin{definition}[Two-bit sensing arena]\label{def:sensing}
Let $W=\{0,1\}^2$. At the root choose sensor $a\in\{0,1\}$ and set root loss
identically to zero. The chosen edge leads to nature node $n_a$; nature reveals
$o=x_a$ and moves to predictor node
$h_{a,o}$, where report $r\in\{0,1\}$ is compared by zero--one loss with a
specified target. Let
$U_{a,o}=\downarrow h_{a,o}$ and let
$B=\bigcup_{a,o}U_{a,o}$ be the decision skeleton. Terminal leaves carry no
additional strategy coordinate or loss. This is a perfect-information arena; the
four reporting nodes are distinct predictor information states. Figure~\ref{fig:sensing} displays the arena and its branch incidence diagram. Restriction identifies
$\mathscr V_0(\mathcal T)$ with $\mathscr V_0(B)$.
\end{definition}

\begin{proposition}[Positive sensing]\label{prop:sensing-positive}
For target $x_a$, the arena has a zero-loss strategy.
\end{proposition}

\begin{proof}
Choose sensor $0$ and report the observed bit at each reached node $h_{0,o}$.
\end{proof}

\begin{figure}[htbp]
\centering
\[
\begin{array}{c@{\qquad\qquad}c}
\begin{array}{ccccccc}
&&&\varepsilon&&&\\[-1mm]
&&\scriptstyle a=0\swarrow&&\searrow\scriptstyle a=1&&\\
&n_0&&&&n_1&\\[-1mm]
\scriptstyle o=0\swarrow&&\searrow\scriptstyle o=1&&
\scriptstyle o=0\swarrow&&\searrow\scriptstyle o=1\\
h_{0,0}&&h_{0,1}&&h_{1,0}&&h_{1,1}
\end{array}
&
\begin{array}{ccccccc}
U_{0,0}&&U_{0,1}&&U_{1,0}&&U_{1,1}\\
&\searrow&&\swarrow& &\searrow&\swarrow\\[-1mm]
&&U_0&&U_1&&\\
&&&\searrow\!\!\swarrow&&&\\[-1mm]
&&&R&&&
\end{array}
\end{array}
\]
\caption{Active sensing. Left: the predictor selects a sensor and nature reveals
the selected bit. Right: the seven-object branch-cover incidence diagram, with
four branch patches, two nature-prefix intersections, and the root intersection.}
\label{fig:sensing}
\end{figure}

Let $\mathcal C_{\mathrm{cov}}$ be the category whose objects are all nonempty
finite intersections of the four branch patches $U_{a,o}$, with one arrow from a
larger intersection to each contained smaller intersection. Restrict
$\mathscr V_0$ to this seven-object incidence category.

\begin{theorem}[Full-site and cover-relative sensing outcomes]
\label{thm:sensing-negative}
For target $x_{1-a}$:
\begin{enumerate}[label=(\roman*)]
\item the full patch site has outcome $\mathsf L_0$;
\item the literal cover-incidence diagram has seven objects: the four branches
$U_{a,o}$, the two same-sensor intersections
$U_a=\{\varepsilon,n_a\}$, and the root intersection
$R=\{\varepsilon\}$. Every local set is nonempty, but no matching family exists,
so the cover-relative outcome is $\mathsf D_0$.
\end{enumerate}
The restriction $\mathscr V_0(U_a)\to\mathscr V_0(R)$ is a bijection on the
root choice. Consequently the seven-object intersection diagram is naturally
equivalent for matching-family calculations to the reduced five-object star
with four branch objects and the root.
\end{theorem}

\begin{proof}
For (i), the decision skeleton $B$ has empty activated-validity set. Either root
choice activates a second-stage node whose information fixes only the sensed bit
while the target unsensed bit takes both values. Since $B$ is an object of the
full patch site, the full-site outcome is local emptiness.

For (ii), on $U_{a,o}$ choose root action $1-a$; then $h_{a,o}$ is unreachable,
so either report gives a valid local strategy. The overlap objects contain no
predictor node beyond the root and have identically zero root loss, hence are
locally inhabited. A matching family agrees through the overlap arrows and
therefore fixes one root action $a^*$. On each $U_{a^*,o}$ that action reaches
$h_{a^*,o}$ for two worlds with opposite unsensed bits, so no report is valid for
both. This contradicts local validity of the matching family.
\end{proof}

\begin{proposition}[Cover-relative boundary]
\label{prop:cover-boundary}
Let $F$ be a restriction-stable validity subfunctor on a patch category with an
initial full object $B$, and let $\mathcal C_{\mathrm{cov}}$ be a cover-incidence
subcategory not containing $B$. If $F(B)=\varnothing$, every object of
$\mathcal C_{\mathrm{cov}}$ has nonempty local set, and the matching-family space
on $\mathcal C_{\mathrm{cov}}$ is empty, then the full-site presentation has
outcome $\mathsf L$ and the cover-incidence presentation has outcome
$\mathsf D$.
\end{proposition}

\begin{proof}
The full-site outcome is $\mathsf L$ because the initial object $B$ has empty
local set. On the cover-incidence subcategory, all local sets are nonempty by
hypothesis, while the compatible-family space is empty; this is precisely the
$\mathsf D$ outcome.
\end{proof}

The cover-relative result is a commitment obstruction: matching forces a common
root action before branch-local obligations are checked. It is obtained after
deleting the full initial object and is not an independent failure of the full
patch sheaf. Proposition~\ref{prop:cover-boundary} gives the general
presentation-relative pattern.

\subsection{Finite imperfect information and one-shot realisation}
\label{subsec:oneshot}

\begin{definition}[Finite imperfect-information arena]\label{def:ii-arena}
A finite imperfect-information arena is a finite arena together with a partition
of predictor nodes into information sets $\mathcal I$. All nodes in
$J\in\mathcal I$ have one common choice set $C_J$, and no information set
contains two nodes on one root-to-leaf path. A pure strategy assigns one value
$\pi(J)\in C_J$ to every information set. No-absentmindedness is the only
information-set hypothesis used in Theorem~\ref{thm:ii-sheaf}; the stronger
global recall conditions of Definition~\ref{def:information-refinement} are
separate assumptions.
\end{definition}

\begin{lemma}[Imperfect-information path characterisation]\label{lem:ii-path}
A node $h$ is reached by $(\pi,w)$ iff $w\in W_h$ and, at every predictor
ancestor $g$ of $h$, the path edge is labelled by $\pi(J_g)$, where $J_g$ is
the information set containing $g$.
\end{lemma}

\begin{proof}
Induct along the unique path as in Lemma~\ref{lem:path-reach}. The condition that
one information set contains no two nodes on a path ensures the play to $h$ never
consults the current information set earlier on that path.
\end{proof}

\begin{definition}[Information-saturated patches]\label{def:ii-patch}
A patch $U$ is information saturated when
$h\in U\cap\mathcal T_P$ implies $J_h\subseteq U$. Intersections and unions of
saturated prefix-closed patches remain saturated and prefix closed. Define
\[
\mathscr S_{\mathcal I}(U)
=\prod_{\substack{J\in\mathcal I\\J\subseteq U}}C_J.
\]
For $s\in\mathscr S_{\mathcal I}(U)$, define
$\operatorname{Reach}_{\mathcal I}(w,h;s)$ by the ordinary arena recursion,
using choice $s(J_g)$ at each predictor ancestor $g$. Prefix closure and
saturation ensure $J_g\subseteq U$. Put
\[
\mathscr V_{\eps,\mathcal I}(U)
=\left\{s\in\mathscr S_{\mathcal I}(U):
\forall h\in U\cap\mathcal T_P\,\forall w\in W_h,
\operatorname{Reach}_{\mathcal I}(w,h;s)
\Rightarrow
\lambda_h(w,s(J_h))\le\eps
\right\}.
\]
\end{definition}

\begin{theorem}[Imperfect-information activated subsheaf]\label{thm:ii-sheaf}
On information-saturated patches, $\mathscr S_{\mathcal I}$ is a sheaf for union
covers and $\mathscr V_{\eps,\mathcal I}$ is a subsheaf.
\end{theorem}

\begin{proof}
A matching family assigns one value to each information set: if one node of an
information set lies in a covering patch, saturation places the whole set there.
Matching makes the value independent of the patch, giving unique coordinatewise
gluing.

For restriction, prefix closure and saturation retain every information set
consulted on the path to a retained node, so reachability and loss there are
unchanged. For gluing, choose a patch containing a reached node. It contains all
ancestor information sets; matching makes the local and glued runs coincide, and
local activated validity supplies the loss bound.
\end{proof}

\begin{remark}
The proof of Theorem~\ref{thm:ii-sheaf} uses neither finiteness nor perfect
recall; these restrictions are retained for the surrounding capability and
complexity results.
\end{remark}


\begin{corollary}[Full-site collapse with imperfect information]
\label{cor:ii-collapse}
On the category of nonempty information-saturated prefix-closed patches, the
full tree is initial, so
\[
\Gamma_{\mathrm{set}}(\mathscr V_{\eps,\mathcal I})
\cong\mathscr V_{\eps,\mathcal I}(\mathcal T).
\]
Thus the full site still has no $\mathsf D_\eps$ outcome.
\end{corollary}

\begin{proof}
Apply Theorem~\ref{thm:initial-collapse} to
Theorem~\ref{thm:ii-sheaf}.
\end{proof}

\begin{proposition}[Behavioural-coordinate sheaf]
\label{prop:behavioural-sheaf}
Assume the information-set choice sets are finite. For an
information-saturated patch $U$, put
\[
\mathscr B(U)=
\prod_{\substack{J\in\mathcal I\\J\subseteq U}}\Delta(C_J),
\]
with restriction by dropping coordinates. Then $\mathscr B$ is a sheaf on the poset of information-saturated patches
for union covers by information-saturated patches, and the Dirac embedding
$\mathscr S_{\mathcal I}\hookrightarrow\mathscr B$ commutes with restriction.
\end{proposition}

\begin{proof}
A matching family assigns the same probability vector to an information set in
every covering patch containing it. Choose that common vector for each
information set in the union. This gives a unique global element of
$\mathscr B(U)$ and restricts to every local member. Applying the Dirac map in
each coordinate plainly commutes with deleting coordinates.
\end{proof}

\begin{remark}
Proposition~\ref{prop:behavioural-sheaf} is a coordinate-gluing statement and
does not require perfect recall. A theorem about expected-loss validity or
Kuhn realisation equivalence would additionally require player ownership,
probabilities for nature, zero-reach conventions, and a specified expected-loss
objective. In particular, a Bell coupling is shared randomisation over complete
pure response tables and is not identified here with independent behavioural
randomisation at information sets.
\end{remark}


\begin{definition}[Information refinement and recall]
\label{def:information-refinement}
Let two imperfect-information arenas have the same tree, nature maps, choice
sets, and losses. Write $\mathcal T\preceq\mathcal T'$ when every information
set of $\mathcal T'$ is contained in an information set of $\mathcal T$.
The arena has \emph{information recall} when any two nodes in one information
set have the same ordered sequence of predecessor information sets. It has
\emph{action recall} when the corresponding predecessor predictor-edge labels
also agree. Perfect recall is the conjunction. This is a global predictor-level recall
condition, stronger than the standard player-relative perfect-recall condition
in multi-player extensive-form games; no player ownership relation is part of
this definition.

For a no-absentminded arena $\mathcal T$, its action-recall refinement
$\mathcal T^+$ splits each information set according to the complete sequence
of predecessor predictor-edge labels, leaving every other arena field
unchanged.
\end{definition}

\begin{proposition}[Capability monotonicity under information refinement]
\label{prop:information-refinement}
If $\mathcal T\preceq\mathcal T'$, pure $\eps$-capability of $\mathcal T$
implies pure $\eps$-capability of $\mathcal T'$.
\end{proposition}

\begin{proof}
Given a strategy on the coarser partition, assign its value at an original
information set to every refined subset. Every world then follows the same
predictor edges in the two arenas, so reached nodes and losses agree.
\end{proof}


\begin{proposition}[Reachable-action collapse criterion]
\label{prop:refinement-collapse}
Let $\mathcal T\preceq\mathcal T'$ and let $\sigma$ be a pure strategy on the
refined arena. There is a coarse strategy $\pi$ inducing the same play as
$\sigma$ at every world if and only if, for every original information set
$J$, all refined classes $J'\subseteq J$ containing a $\sigma$-reachable node
receive one common action. When this condition holds, all worldwise losses are
preserved.
\end{proposition}

\begin{proof}
Suppose first that the condition holds. For each original information set $J$
with a $\sigma$-reachable node, let $c_J$ be the common action on its reachable
refined classes; choose an arbitrary element of $C_J$ otherwise. Define
$\pi(J)=c_J$. Induct down the play at any fixed world. The root and every nature
move agree because the underlying arena data are unchanged. If the plays have
agreed up to a reached predictor node $h\in J$, then the refined class $J'_h$
contains a $\sigma$-reachable node and
$\pi(J)=c_J=\sigma(J'_h)$. Thus the same predictor edge is selected. The full
plays, and hence all reached losses, agree.

Conversely, suppose a coarse strategy $\pi$ reproduces every $\sigma$-play.
If a refined class $J'\subseteq J$ contains a node reached under $\sigma$ at
some world, equality of the two plays at that node forces
$\sigma(J')=\pi(J)$. Therefore all reachable refined classes inside $J$ have
the common value $\pi(J)$.
\end{proof}

\begin{theorem}[Pure action-recall redundancy]
\label{thm:action-recall-equivalence}
Let $\mathcal T$ be a finite no-absentminded arena with information recall in
the global predictor-level sense of Definition~\ref{def:information-refinement}.
For every tolerance,
\[
\mathcal T\text{ is pure $\eps$-capable}
\iff
\mathcal T^+\text{ is pure $\eps$-capable}.
\]
\end{theorem}

\begin{proof}
The forward implication is
Proposition~\ref{prop:information-refinement}. For the converse, let $\sigma$
be a capable pure strategy on $\mathcal T^+$. Call a node
$\sigma$-reachable when it is reached by $(\sigma,w)$ for some world $w$.
First prove by induction on the length of the predecessor information-set
sequence that all $\sigma$-reachable nodes in one original information set $J$
have the same predecessor action sequence.

The assertion is immediate when the sequence is empty. Suppose the common
information-set sequence supplied by information recall is
$J_1,\ldots,J_t$, and let $h,h'\in J$ be $\sigma$-reachable. Their predecessor
nodes in $J_t$ are also reachable. By the induction hypothesis those nodes have
the same action prefix and hence belong to one refined class inside $J_t$.
The deterministic strategy $\sigma$ assigns that class one action. Nature nodes
between that last predictor ancestor and $h$ or $h'$ add no predictor-edge
labels, so appending this common action proves that $h$ and $h'$ have the same
predecessor action sequence. Thus all reachable nodes of $J$ lie in one refined class $J^+$.

The claim is exactly the reachable-action condition of
Proposition~\ref{prop:refinement-collapse}. That proposition supplies a coarse
strategy with the same play and losses at every world. It is therefore
$\eps$-capable.
\end{proof}

\begin{proposition}[Sharpness of information recall]
\label{prop:action-recall-sharp}
There is a finite no-absentminded arena without information recall which is not
$0$-capable although its action-recall refinement is $0$-capable.
\end{proposition}

\begin{proof}
Let $W=\{0,1\}$. The nature root reveals $w$ and moves to a predictor node
$a_w$; the two nodes $a_0,a_1$ are singleton information sets, have binary
choices, and zero loss. Choice $c$ leads to $b_{w,c}$. Put all four nodes
$b_{w,c}$ in one information set $J$ with binary report $r$, and set its loss
to zero exactly when $r=w$.

A strategy in the original arena assigns one report to $J$ and therefore fails
on one world. In the action-recall refinement, $J$ splits into classes $J_0$
and $J_1$ according to the first action. Choose action $w$ at $a_w$ and report
$c$ at $J_c$. The reached report equals the world, giving zero loss. Information
recall fails because nodes over worlds $0$ and $1$ have predecessor information
sets $\{a_0\}$ and $\{a_1\}$, respectively.
\end{proof}

\begin{remark}
Theorem~\ref{thm:action-recall-equivalence} is a pure-strategy, worst-case-loss
statement. It neither identifies mixed and behavioural strategies nor implies
one-shot structure or a complexity bound.
\end{remark}



\begin{definition}[One-shot team arena]\label{def:oneshot}
Let a finite distributed selector problem be given and order its agents as
$J=\{j_1,\ldots,j_m\}$. Define a finite arena as follows.
\begin{enumerate}[label=(\roman*)]
\item The root is a nature node with move set $W$ and nature map
$\nu_\varepsilon(w)=w$; its child is denoted $h_w$ and has
$W_{h_w}=\{w\}$.
\item For $1\le k\le m$, each history
$h_{w,\mathbf c_{<k}}$ is a predictor node for agent $j_k$, with choice set
$C_{j_k}$ and one child for each $c_k\in C_{j_k}$. At level $k=1$ the preceding
action tuple is empty.
\item The information set associated with agent $j_k$ and observation
$z\in\im(I_{j_k})$ is
\[
\mathcal I_{j_k,z}
=\{h_{w,\mathbf c_{<k}}:I_{j_k}(w)=z\}.
\]
All nodes in this set have choice set $C_{j_k}$.
\item Loss is zero at levels $k<m$. At the final predictor node,
\[
\lambda_{h_{w,\mathbf c_{<m}}}(w,c_m)=
\begin{cases}
0,&(c_1,\ldots,c_m)|_{S(w)}\in V_w,\\
1,&\text{otherwise}.
\end{cases}
\]
The final choice leads to a terminal child.
\end{enumerate}
Every agent acts at every world, including worlds at which that agent is not in
$S(w)$; terminal validity ignores coordinates outside $S(w)$. Earlier actions
are not included in later agents' information labels.
\end{definition}

Each distributed agent moves only once, so a player-relative perfect-recall
condition would be automatic after adding player ownership. In the global
single-predictor sense of Definition~\ref{def:information-refinement}, however,
the one-shot arena need not have information recall, since later information
sets may identify histories with different predecessor information sets.

\begin{theorem}[One-shot team-arena realisation]\label{thm:oneshot}
For every finite distributed selector problem:
\begin{enumerate}[label=(\roman*)]
\item pure team strategies are in bijection with full local policy tuples
\[
\prod_{j\in J}C_j^{\im(I_j)};
\]
restriction to active observations gives the canonical CSP assignment;
\item a distributed strategy is valid iff the corresponding team strategy has
zero worst-case terminal loss;
\item the full information-saturated patch site has outcome $\mathsf L_0$ if
and only if the distributed problem is unsolvable; if every realised-scope
relation is nonempty but no strategy exists, the canonical incidence diagram
has outcome $\mathsf D_0$.
\end{enumerate}
\end{theorem}

\begin{proof}
A pure team strategy chooses one value at each information set
$\mathcal I_{j,z}$, exactly a function $\im(I_j)\to C_j$ for every agent. These
functions form the full distributed policy tuple. Restricting them to active
observations gives the CSP assignment, while inactive values do not enter
terminal validity.

At world $w$, the action of agent $j$ is $\pi_j(I_j(w))$. The terminal loss
vanishes exactly when the tuple on active scope $S(w)$ lies in $V_w$, proving
(ii). If the distributed problem is unsolvable, part (ii) makes the activated-validity
set of the full tree empty, so the full-site outcome is $\mathsf L_0$ by
Corollary~\ref{cor:ii-collapse}. Conversely, a valid distributed strategy gives
a valid full-tree strategy by part (ii), and its restrictions inhabit every
patch; hence the full site cannot have $\mathsf L_0$. The final assertion of
(iii) is Theorem~\ref{thm:distributed-equivalence}.
\end{proof}


\begin{definition}[Blind-order one-shot normal form]
\label{def:blind-one-shot}
A finite imperfect-information arena is in blind-order one-shot normal form when
it is equipped with a finite world set $W$, an ordered agent set
$J=\{j_1,\ldots,j_m\}$, finite choice sets $C_j$, finite observation sets
$Z_j$, interfaces $I_j:W\to Z_j$, nonempty scopes $S(w)\subseteq J$, and
relations $V_w\subseteq\prod_{j\in S(w)}C_j$, with the following structure.
The root is a nature node selecting $w$. At level $k$, the predictor nodes are
indexed by
\[
(w,c_{<k})\in W\times\prod_{\ell<k}C_{j_\ell},
\]
and all nodes with the same value of $I_{j_k}(w)$ form one information set,
independently of $c_{<k}$. Intermediate losses vanish, and terminal loss is zero
exactly when
\[
(c_j)_{j\in S(w)}\in V_w.
\]
All choices lead to the corresponding action-prefix child. Isomorphisms of such
arenas preserve the rooted labelled tree, node types, nature branches,
information partition, choice sets, and losses.
\end{definition}

\begin{definition}[Distributed problem of a normal arena]
\label{def:arena-to-distributed}
For a blind-order one-shot arena $\mathcal T$, let
$\mathsf{Dist}(\mathcal T)$ be the finite distributed selector problem given by its
worlds, ordered agents, interfaces, choices, scopes, and terminal relations from
Definition~\ref{def:blind-one-shot}. Write $\mathsf N(P)$ for the arena of
Definition~\ref{def:oneshot} constructed from a finite distributed problem $P$.
Here finite means that the world, active observation, and choice sets and all
relation tables are finite.
\end{definition}

\begin{theorem}[One-shot normal-form equivalence]
\label{thm:one-shot-normal-form}
For every finite distributed selector problem $P$ and every blind-order
one-shot arena $\mathcal T$:
\begin{enumerate}[label=(\roman*)]
\item $\mathsf N(P)$ is in blind-order one-shot normal form and
\[
\mathrm{Strat}(\mathsf N(P))
\cong\prod_{j\in J}C_j^{\im(I_j)};
\]
\item this bijection preserves every worldwise zero--one terminal loss and hence
preserves and reflects capability;
\item $\mathsf N(\mathsf{Dist}(\mathcal T))$ is isomorphic to $\mathcal T$;
\item $\mathsf{Dist}(\mathsf N(P))$ recovers $P$ up to restriction of each
observation set to its realised image and renaming of observations and choices; the world set is recovered verbatim.
\end{enumerate}
Thus finite distributed selectors and blind-order one-shot arenas are equivalent
for pure-strategy capability and worldwise zero--one loss. This is not, by
itself, a diagnostic equivalence between the canonical CSP incidence
presentation and the full arena patch-site presentation. Morphisms and
functoriality require additional structure beyond this object-level statement.
\end{theorem}

\begin{proof}
At each information set $\mathcal I_{j,z}$, a pure strategy chooses one element
of $C_j$. Since these information sets are indexed exactly by
$z\in\im(I_j)$, choices at all information sets are precisely a tuple of maps
$\pi_j:\im(I_j)\to C_j$, proving (i).

Fix a world $w$. At level $k$, the reached node has action prefix equal to the
previously selected values, and the next action is
$\pi_{j_k}(I_{j_k}(w))$. The terminal tuple is therefore the distributed
response tuple. By construction, terminal loss is zero exactly when its
restriction to $S(w)$ belongs to $V_w$. This proves (ii).

For (iii), the normal arena already has one node $(w,c_{<k})$ for every world
and action prefix. The operation $\mathsf{Dist}$ records the same worlds, order,
interfaces, choices, scopes, and terminal relations; applying $\mathsf N$
therefore reconstructs the same predecessor and successor maps, information
sets, and losses. The coordinatewise identification is the required arena
isomorphism. Finally, applying $\mathsf{Dist}$ to $\mathsf N(P)$ reads back the
same data used in the construction. Observations outside $\im(I_j)$ never label
an information set, giving exactly the harmless restriction stated in (iv).
\end{proof}


\begin{definition}[World-scheduled one-pass arena]
\label{def:world-scheduled}
Let $P$ be a finite distributed selector problem. A world schedule assigns to
each $w$ an ordering
\[
\kappa_w=(j_{w,1},\ldots,j_{w,m})
\]
of the full finite agent set, with every agent occurring once. The associated
world-scheduled arena has a nature root selecting $w$, followed by predictor
moves in the order $\kappa_w$. Agent $j$ moves at the information set labelled
$(j,I_j(w))$, independently of previous actions. Intermediate losses vanish,
and terminal loss is zero exactly when the selected tuple restricted to $S(w)$
belongs to $V_w$. The schedule is part of the arena data.
\end{definition}

\begin{proposition}[World-scheduled distributed realisation]
\label{prop:world-scheduled}
For every finite distributed problem equipped with world schedules, pure
strategies of the associated world-scheduled arena are canonically
\[
\prod_{j\in J}C_j^{\im(I_j)},
\]
and distributed validity is equivalent to zero worst-case terminal loss. The
arena is reconstructed from the distributed datum together with the schedules,
up to renaming of action-prefix nodes.
\end{proposition}

\begin{proof}
Every information set is indexed by one pair $(j,z)$, so a pure strategy is
exactly one choice in $C_j$ for each realised observation $z$. At world $w$,
the schedule changes only the order in which the values
$\pi_j(I_j(w))$ are selected; the terminal tuple contains exactly those values.
The terminal loss condition is therefore the distributed validity condition.
The distributed data and $\kappa_w$ determine every world branch, action-prefix
child, information-set label, and terminal loss, proving reconstruction.
\end{proof}

\begin{remark}
Proposition~\ref{prop:world-scheduled} is a sufficient structure-preserving
normal form, strictly allowing world-dependent orders. It is not a converse for
arbitrary imperfect-information arenas: ordinary distributed data do not record
schedules, and strategy-dependent signaling can still be represented only by
the extensional compilation of Theorem~\ref{thm:finite-ii-compilation}.
\end{remark}

\begin{theorem}[Finite pure-strategy distributed compilation]
\label{thm:finite-ii-compilation}
Let $\mathcal T$ be any finite imperfect-information arena and fix a tolerance
$\eps$. There is a finite distributed selector problem $P_{\mathcal T,\eps}$
whose strategy space is canonically the pure-strategy space of $\mathcal T$ and
such that distributed capability is equivalent to $\eps$-capability of the
arena.
\end{theorem}

\begin{proof}
Use one distributed agent $a_J$ for every information set $J\in\mathcal I$.
Give it the singleton observation space $Z_{a_J}=\{\ast\}$ and choice set $C_J$.
Use the same world set $W$, take every agent active at every world, and define
\[
V_w^\eps=
\left\{(c_J)_{J\in\mathcal I}:
L_{\mathcal T}((c_J)_{J\in\mathcal I},w)\le\eps\right\}.
\]
A distributed strategy chooses one value in $C_J$ for every singleton
interface, exactly a pure arena strategy. Membership in $V_w^\eps$ is, by
definition, the arena loss bound at $w$. Universal validity over worlds is
therefore equivalent to $\eps$-capability.
\end{proof}

\begin{remark}
Theorem~\ref{thm:finite-ii-compilation} is extensional. Its world relations can
mention the complete pure-strategy tuple and can be exponentially large under
explicit tabulation. It preserves strategy space, worldwise validity, and
capability, but not temporal locality, the original arena tree, or the diagnostic
outcome attached to a chosen patch presentation. By contrast,
Theorem~\ref{thm:one-shot-normal-form} characterises the restricted class for
which the distributed presentation reconstructs the arena itself.
\end{remark}

\begin{theorem}[Distributional one-shot realisation of Bell locality]
\label{thm:bell-oneshot}
For a finite Bell scenario, let $\mathcal T_{\mathrm{Bell}}$ be the one-shot
information arena whose parties are agents, local settings are observations, and
local outcomes are choices. For every finite behaviour $p$, the following are
equivalent:
\begin{enumerate}[label=(\roman*)]
\item $p$ has a local hidden-variable model;
\item $\operatorname{Coupl}(p)\ne\varnothing$;
\item there is $\mu\in\Delta(\mathrm{Strat}(\mathcal T_{\mathrm{Bell}}))$
such that evaluation at every setting tuple has distribution $p_{\mathbf x}$;
\item the Bell star diagram has a compatible global family.
\end{enumerate}
\end{theorem}

\begin{proof}
Theorem~\ref{thm:bell} proves the equivalence of (i), (ii), and (iv). By
Theorem~\ref{thm:oneshot}(i),
\[
\mathrm{Strat}(\mathcal T_{\mathrm{Bell}})
\cong\prod_iO_i^{X_i}=\Omega.
\]
Under this bijection, evaluation of a pure strategy at $\mathbf x$ is exactly
$\pi_{\mathbf x}$. The induced affine bijection between the two probability
simplices therefore sends the distribution of mixed-strategy evaluation to
$(\pi_{\mathbf x})_*\mu$. The marginal equations in (iii) are consequently
exactly the defining equations of $\operatorname{Coupl}(p)$, proving
(ii)$\Leftrightarrow$(iii).
\end{proof}

\begin{remark}
Theorem~\ref{thm:bell-oneshot} concerns mixed-strategy marginal realisation, not
the worst-case capability predicate. The deterministic CHSH winning loss and a
Bell behaviour use the same one-shot information skeleton and deterministic
response tables but impose different, respectively support-level and
probability-level, conditions.
\end{remark}

\begin{corollary}[One-shot deterministic CHSH arena]\label{cor:chsh-arena}
The private-information team problem of Corollary~\ref{cor:team-decision} has a
finite one-shot imperfect-information arena realisation. Nature chooses
$(x,y)$, the first agent observes only $x$, the second only $y$, and terminal
loss vanishes exactly when $a\oplus b=x\wedge y$. These are the support
constraints of the CHSH winning predicate; the statement concerns impossibility
of a deterministic strategy that wins every input.
\end{corollary}

\begin{proof}
Apply Theorem~\ref{thm:oneshot} to the two-agent distributed selector.
\end{proof}

\begin{proposition}[Explicit-table complexity boundary]\label{prop:pnp}
On finite perfect-information arenas, backward induction is a bottom-up pass
using $O(\sum_h|W_h|\,|C_h|)$ loss-table checks. Capability for explicitly
represented finite imperfect-information arenas is NP-complete.
\end{proposition}

\begin{proof}
Membership in NP follows by guessing one choice per information set and checking
all explicitly represented world--node loss constraints. For hardness, reduce the NP-complete variant of 3SAT in which each clause has
three distinct variables. Nature selects a clause. The arena then queries, in a
fixed order, the information sets corresponding to the variables occurring in
that clause;
all occurrences of the same Boolean variable in different clauses share one
binary information set, whose chosen value is interpreted as the truth value of
that variable. The terminal loss is zero exactly when the selected clause is
satisfied. This arena has size polynomial in the formula, because each branch
contains only the variables of one clause. A pure strategy is a truth assignment,
and it is zero-capable exactly when it satisfies every clause. Thus the explicit
perfect-information regime is polynomial-time by backward induction, whereas the
explicit imperfect-information regime is NP-complete. The generic blind-order
normal-form construction of Definition~\ref{def:oneshot} may be exponential in
the number of distributed agents and is not used for this explicit-tree hardness
reduction; for compressed blind-order input data, capability is exactly the
finite distributed-selector problem of Theorem~\ref{thm:distributed-NP}.
\end{proof}




\subsection{Static compilation}

\begin{definition}[Static compilation]\label{def:staticcomp}
A static compilation consists of a compiled game with world set $W^\sharp$, a
surjection $j:W\twoheadrightarrow W^\sharp$, and a strategy transformation
$K:\mathrm{Strat}(\mathcal T)\to\Pi^\sharp$ whose image is the declared
compiled policy class and which preserves aggregate losses:
\[
L^\sharp(K(\pi),j(w))=L_{\mathcal T}(\pi,w).
\]
\end{definition}

\begin{proposition}[Compiled selector semantics]\label{prop:FBcompile}
Under a static compilation, sequential and compiled capability are equivalent.
\end{proposition}

\begin{proof}
A sequential witness maps to a compiled policy. For $u\in W^\sharp$, choose
$w$ with $j(w)=u$ and use loss preservation. Conversely, every declared
compiled policy lies in the image of $K$; choose a preimage strategy and use
loss preservation at every world.
\end{proof}

\begin{proposition}[Composition of static compilations]
\label{prop:static-compilation-composition}
Static compilations compose. If
\((j_1,K_1):A\to B\) and \((j_2,K_2):B\to C\) preserve the corresponding
worldwise aggregate losses, then
\[
(j_2\circ j_1,\ K_2\circ K_1):A\to C
\]
is a static compilation. Identity maps give identity compilations, so static
compilations form a category on any fixed class of compatible source and target
presentations.
\end{proposition}

\begin{proof}
The composite of surjections is a surjection, and the composite strategy map has
image in the declared target policy class. For every strategy \(\pi\) and world
\(w\),
\[
L_C(K_2(K_1(\pi)),j_2(j_1(w)))
=L_B(K_1(\pi),j_1(w))
=L_A(\pi,w).
\]
The identity case and associativity are ordinary identities and associativity of
function composition.
\end{proof}

\begin{remark}
A static compilation is a capability equivalence. Preservation of the diagnostic
outcome additionally requires transport of the declared local objects,
admissibility sets, and realised global-family spaces as part of the compilation
data.
\end{remark}

\subsection{Operational guarded feedback and universal concurrency}

A finite causal poset represents universally executed cuts whose order is
partial. Protocols in which an action determines whether a later cut occurs are
instead represented by the branching arenas of Section~\ref{sec:feedback}.

\begin{definition}[Finite-poset guarded protocol]
\label{def:operational-feedback}
A finite-poset guarded protocol consists of a finite strict poset
$(T,\prec)$, a finite nonempty world set $W$, and, for each $\tau\in T$,
finite nonempty sets $E_\tau,Z_\tau,R_\tau,A_\tau,M_\tau,H_\tau$. Put
$C_\tau=R_\tau\times A_\tau$ and
\[
\mathbf{Hist}_{<\tau}
=\prod_{\sigma\prec\tau}
(Z_\sigma\times M_\sigma\times A_\sigma\times H_\sigma),
\]
where this predecessor-indexed product is a set of functions on the strict
past, and is a singleton when the past is empty. The structural maps are
\[
Q_\tau:W\to E_\tau,
\qquad
I_\tau:E_\tau\times\mathbf{Hist}_{<\tau}\to Z_\tau,
\qquad
D_\tau:R_\tau\to M_\tau,
\]
\[
U_\tau:W\times\mathbf{Hist}_{<\tau}\times M_\tau\times A_\tau
\to H_\tau,
\]
\[
\lambda_\tau:
W\times\mathbf{Hist}_{<\tau}\times C_\tau\to[0,\infty].
\]
The independently supplied map $Q_\tau$ is the accessible exogenous record at
cut $\tau$. Strict-past typing is the guard: no cut reads its own or a future
record. Every cut occurs in every execution.

For an arbitrary complete choice profile $c=(c_\sigma)_\sigma\in
\prod_\sigma C_\sigma$, its syntactic record at $\tau$ is defined by
well-founded recursion, writing $c_\sigma=(r_\sigma,a_\sigma)$:
\[
z_\tau^c(w)=I_\tau(Q_\tau(w),\mathbf h_{<\tau}^c(w)),
\quad m_\tau^c=D_\tau(r_\tau),
\quad
h_\tau^c(w)=U_\tau(w,\mathbf h_{<\tau}^c(w),m_\tau^c,a_\tau),
\]
where $\mathbf h_{<\tau}^c$ collects the predecessor tuples
$(z_\sigma^c,m_\sigma^c,a_\sigma,h_\sigma^c)$. For such a complete profile also define
\[
L_{\mathrm{op}}(c,w)=
\max\left(
\{\lambda_\tau(w,\mathbf h_{<\tau}^c(w),c_\tau):\tau\in T\}
\cup\{0\}\right).
\]
Define
\[
Z_\tau^{\mathrm{act}}
=\{z_\tau^c(w):w\in W,\ c\in\prod_\sigma C_\sigma\}.
\]
A policy is a family
$\pi=(\pi_\tau)_\tau$ with
$\pi_\tau:Z_\tau^{\mathrm{act}}\to C_\tau$.
It records the realised report--action pair. The larger two-component policy
space $R_\tau^{Z_\tau}\times A_\tau^{Z_\tau\times R_\tau}$ has this space as
its observational quotient; counterfactual actions at unchosen reports are not
part of the operational semantics. If report-conditional counterfactual actions
are required, a cut should be split into a report decision followed by an action
decision, or the action carrier should be replaced by a function carrier
$R_\tau\to A_\tau$.
\end{definition}

\begin{definition}[Poset execution and capability]
\label{def:operational-execution}
For a policy $\pi$ and world $w$, an execution is a family
$(z_\tau,c_\tau,m_\tau,h_\tau)_{\tau\in T}$ satisfying, with
$c_\tau=(r_\tau,a_\tau)$,
\[
z_\tau=I_\tau(Q_\tau(w),\mathbf h_{<\tau}),
\qquad c_\tau=\pi_\tau(z_\tau),
\qquad m_\tau=D_\tau(r_\tau),
\]
\[
h_\tau=U_\tau(w,\mathbf h_{<\tau},m_\tau,a_\tau).
\]
Its loss is
\[
L_{\mathrm{op}}(\pi,w)
=\max\left(
\{\lambda_\tau(w,\mathbf h_{<\tau},c_\tau):\tau\in T\}
\cup\{0\}\right).
\]
For $\Pi\subseteq\prod_\tau C_\tau^{Z_\tau^{\mathrm{act}}}$, the protocol is
$\eps$-capable in $\Pi$ when some $\pi\in\Pi$ has
$L_{\mathrm{op}}(\pi,w)\le\eps$ for every world.
\end{definition}

\begin{theorem}[Unique extension-independent execution]
\label{thm:operational-execution}
Every policy and world determine a unique execution. Evaluating its equations in
any linear extension of $(T,\prec)$ gives the same $T$-indexed family and the
same aggregate loss.
\end{theorem}

\begin{proof}
Apply Proposition~\ref{prop:guarded-recursion} to the finite strict poset. At
$\tau$, every
coordinate of $\mathbf h_{<\tau}$ has already been determined, so the displayed
equations determine successively $z_\tau,c_\tau,m_\tau,h_\tau$. Well-founded
induction gives uniqueness. A linear extension orders these evaluations;
since the equation at $\tau$ reads exactly the strict-predecessor coordinates,
its value is the unique recursively determined value in every extension. The
aggregate loss is a function of that indexed family and is therefore unchanged.
\end{proof}

\begin{remark}[Partial order of cuts]
A finite cut poset records concurrency without making execution nondeterministic.
Incomparable cuts may be evaluated in either order, but each cut reads only its
strict causal past. The resulting execution is therefore indexed by the poset,
not by the chosen linear extension used to compute it.
\end{remark}

\begin{definition}[Order-ideal configuration system]
\label{def:order-ideal-configuration}
Let $\mathcal I(T)$ be the finite lattice of order ideals of $T$. For fixed
$(\pi,w)$, a configuration over $C\in\mathcal I(T)$ records the execution
tuples at cuts in $C$. A cut $\tau\notin C$ is enabled when every predecessor
of $\tau$ lies in $C$; the transition $C\to C\cup\{\tau\}$ computes its tuple
from the recorded strict past.
\end{definition}

\begin{proposition}[Concurrency diamond and confluence]
\label{prop:operational-confluence}
If incomparable cuts $\tau$ and $\sigma$ are both enabled at a configuration,
then executing $\tau$ followed by $\sigma$ or $\sigma$ followed by $\tau$
produces the same configuration over $C\cup\{\tau,\sigma\}$. Every maximal path
through the configuration system therefore yields the unique execution of
Theorem~\ref{thm:operational-execution}.
\end{proposition}

\begin{proof}
Simultaneous enablement implies that neither cut is a predecessor of the other.
The computation at $\tau$ reads only coordinates indexed by
$\{\xi:\xi\prec\tau\}\subseteq C$ and writes only the $\tau$ coordinate;
the analogous statement holds for $\sigma$. Thus the first transition changes
no input read by the second, and the two computed tuples are identical in either
order. Adjacent swaps of simultaneously enabled incomparable cuts connect the
maximal paths associated with linear extensions; repeated use of the diamond
property gives the same final record. Equivalently, confluence follows from the
uniqueness theorem.
\end{proof}

\begin{definition}[Linear-extension arena]
\label{def:operational-arena}
Fix a linear extension $L=(\tau_1,\ldots,\tau_m)$. The arena
$\mathcal A_L(P)$ has a nature root selecting $w\in W$. At stage $k$ it has a
predictor node $h^k_{w,c_{<k}}$ for every world and every prefix
$c_{<k}\in\prod_{i<k}C_{\tau_i}$, with choice set $C_{\tau_k}$. Compute the
causal-predecessor record from $w$ and the prefix, ignoring earlier incomparable
cuts, and label the node by
\[
z_{\tau_k}=I_{\tau_k}(Q_{\tau_k}(w),\mathbf h_{<\tau_k}).
\]
Two stage-$k$ nodes are in the same information set exactly when these labels
agree. The choice $c_k$ leads to $h^{k+1}_{w,c_{\le k}}$, or to a terminal at
$k=m$, and node loss is
\[
\lambda_{h^k_{w,c_{<k}}}(w,c_k)
=\lambda_{\tau_k}(w,\mathbf h_{<\tau_k},c_k).
\]
All worlds and prefixes are included before a policy is selected.
\end{definition}

\begin{theorem}[Operational arena realisation]
\label{thm:operational-arena}
For every linear extension $L$, $\mathcal A_L(P)$ is a finite
imperfect-information arena, independent of the policy, and
\[
\mathrm{Strat}(\mathcal A_L(P))
\cong\prod_{\tau\in T}C_\tau^{Z_\tau^{\mathrm{act}}}.
\]
Under this bijection, arena play and operational execution agree at every cut,
\[
L_{\mathcal A_L(P)}(\pi,w)=L_{\mathrm{op}}(\pi,w),
\]
and capability is preserved and reflected for every transported policy class.
Different linear extensions need not give isomorphic trees, but they realise the
same operational policy space and worldwise loss functional. On the
information-saturated patch site of each arena, activated validity is a
subsheaf and the full tree is initial, so full-site
$\mathsf D_\eps$ is impossible.
\end{theorem}

\begin{proof}
Finiteness and no-absentmindedness follow because the data and prefix products
are finite and each information set lies at one stage. The arena contains all
syntactic prefixes, including prefixes that may be unreachable under a fixed
observation-indexed policy. Every label in $Z_\tau^{\mathrm{act}}$ occurs at a
node: a complete syntactic choice profile
witnessing the label restricts to a prefix in any linear extension, and
incomparable earlier choices are ignored. Conversely every node label is in the
defined active image. Hence a pure strategy is exactly one map
$Z_\tau^{\mathrm{act}}\to C_\tau$ per cut.

Fix a policy and world. Induct along $L$. All causal predecessors of
$\tau_k$ occur earlier, and by induction their arena records equal their
operational records. The node label is therefore the operational observation,
and the selected edge is the operational choice. The stage record and loss then
agree by definition. Taking maxima proves loss equality, and quantifying over
worlds proves capability equivalence. Extension independence follows from
Theorem~\ref{thm:operational-execution}. The patch statement is
Theorem~\ref{thm:ii-sheaf} and Corollary~\ref{cor:ii-collapse}.
\end{proof}


\begin{proposition}[Size of a linear-extension arena]
\label{prop:operational-arena-size}
For $L=(\tau_1,\ldots,\tau_m)$, the arena of
Definition~\ref{def:operational-arena} has
\[
1+|W|\sum_{k=1}^{m}\prod_{i<k}|C_{\tau_i}|
+|W|\prod_{i=1}^{m}|C_{\tau_i}|
\]
nodes: the root, all predictor nodes, and all terminals. Thus explicit arena
synthesis can be exponential in $m$. Different linear extensions can have
different numbers of intermediate predictor nodes, although
Theorem~\ref{thm:operational-arena} identifies their operational policy spaces
and worldwise loss functionals.
\end{proposition}

\begin{proof}
At stage $k$, a predictor node is determined by a world and one choice for each
of the $k-1$ earlier cuts in $L$, giving
$|W|\prod_{i<k}|C_{\tau_i}|$ nodes. After the final choice, a terminal is
determined by a world and a complete choice profile, giving the last product.
Adding the nature root proves the formula. For $m=0$, the empty product is
understood as $1$; the formula counts a nature root with $|W|$ terminal leaves,
which may be identified without changing capability. The products can grow
exponentially with the number of cuts. Their partial sums can depend on the
order of unequal choice-alphabet sizes, while the operational equivalence is
already proved in Theorem~\ref{thm:operational-arena}.
\end{proof}

\begin{proposition}[Perfect-recall decoding criterion]
\label{prop:operational-perfect-recall}
In the global arena-level sense of Definition~\ref{def:information-refinement},
the arena $\mathcal A_L(P)$ has perfect recall exactly when, for every cut
$\tau$ and realised observation $z$, all syntactic nodes in the information set
$(\tau,z)$ have the same preceding sequence of predictor information sets and
choices. Equivalently, for each realised label there is a decoding map from
$(\tau,z)$ to that common sequence.
\end{proposition}

\begin{proof}
The preceding information-set/action sequence is precisely the predictor's own
history in the synthesized arena. Perfect recall says this sequence is constant
on each current information set. Constancy defines the decoding map, and a
decoding map implies constancy. The existence of the observation map alone is
insufficient.
\end{proof}

\begin{corollary}[Effective operational evaluation]
\label{cor:effective-operational}
For effectively tabulated finite protocol data and a computable policy, the
unique execution and every linear-extension arena are computable uniformly from
the supplied finite codes. Here effectively tabulated means that finite carrier
sets are given by finite tagged codes with decidable equality and that structural
maps, losses, nature maps, and restrictions are given by finite lookup tables or
primitive-recursive indices on the declared finite domains.
\end{corollary}

\begin{proof}
Well-founded evaluation makes finitely many table lookups. Arena synthesis
enumerates finite world and choice-prefix products and performs the same
lookups.
\end{proof}

\paragraph{World-only decoupling.}

\begin{theorem}[Operational decoupling]
\label{thm:operational-decoupling}
Suppose each observation is independent of the interaction history:
\[
I_\tau(Q_\tau(w),\mathbf h)=\bar I_\tau(Q_\tau(w)).
\]
Then the protocol has an equivalent finite distributed selector with agents
$T$, interfaces $w\mapsto\bar I_\tau(Q_\tau(w))$, choices $C_\tau$, and
world relation
\[
V_w^\eps=
\{c\in\prod_{\tau\in T}C_\tau:
L_{\mathrm{op}}(c,w)\le\eps\},
\]
where $L_{\mathrm{op}}(c,w)$ is the complete-profile loss defined in
Definition~\ref{def:operational-feedback}. This preserves and reflects
$\eps$-capability.

If, in addition, there are predicates
$V_{\tau,w}^\eps\subseteq C_\tau$ such that operational aggregate validity is
exactly
\[
L_{\mathrm{op}}(c,w)\le\eps
\iff
c_\tau\in V_{\tau,w}^\eps\quad\text{for every }\tau,
\]
then the relation factors as $V_w^\eps=\prod_\tau V_{\tau,w}^\eps$ and
capability is simultaneous fibrewise selection. In the corresponding discrete fibre-indexed packaging over cuts and
observations, $\mathsf D_\eps$ cannot occur in ZFC.
\end{theorem}

\begin{proof}
Under the hypothesis, a policy supplies at world $w$ exactly the profile
$(\pi_\tau(\bar I_\tau(Q_\tau(w))))_\tau$. By definition, that profile belongs
to $V_w^\eps$ exactly when its operational loss is at most $\eps$. This is the
distributed validity condition. Under the additional equivalence, membership in
the world relation is coordinatewise, giving the displayed product. Grouping
the unary constraints by each interface fibre gives the ordinary fibrewise
selector conditions; the discrete compatibility result then excludes
$\mathsf D_\eps$.
\end{proof}

\paragraph{Observation-preserving compilation.}

\begin{definition}[Observation-preserving static compilation]
\label{def:observation-compilation}
An observation-preserving static compilation of a guarded protocol is a static
compilation in the sense of Definition~\ref{def:staticcomp}, together with a
policy-independent interface $I^\sharp:W^\sharp\to Z^\sharp$ and maps
$q_\tau:Z^\sharp\to Z_\tau$ satisfying
\[
z_\tau(\pi,w)=q_\tau(I^\sharp(j(w)))
\]
for every cut, policy, and world.
\end{definition}

\begin{proposition}[Disclosure obstruction to observation compilation]
\label{prop:observation-compilation-no-go}
If the domain of the strategy transformation contains two policies that at one
world produce different observations at some cut, then no observation-preserving
static compilation exists on that domain. This does not preclude a
loss-preserving static compilation.
\end{proposition}

\begin{proof}
For fixed $w$ the value $q_\tau(I^\sharp(j(w)))$ is independent of the policy,
so it cannot equal two distinct operational observations. The existing static
compilation definition requires only aggregate-loss preservation and therefore
is not contradicted.
\end{proof}


\begin{remark}[Active sensing as guarded operational feedback]
\label{rem:sensing-operational}
The sensing arena of Definition~\ref{def:sensing} is the two-cut chain
$\tau_0\prec\tau_1$. At $\tau_0$ the observation is trivial and the physical
action chooses $a\in\{0,1\}$. At $\tau_1$ the observation generator reads that
action from the predecessor record and returns $x_a$; the current choice is the
binary report. Hence the two stage observations can differ at one world under
different first-stage policies, so
Proposition~\ref{prop:observation-compilation-no-go} excludes an
observation-preserving static compilation. A loss-preserving contingent-plan
compilation is not excluded. The full-site $\mathsf L_0$ and cover-relative
$\mathsf D_0$ outcomes remain those of
Theorem~\ref{thm:sensing-negative}.
\end{remark}

\paragraph{Trace-relative past factorisation.}

\begin{proposition}[Trace-relative past factorisation]
\label{prop:trace-relative-past}
Fix a cut $\tau$ and a syntactically valid choice assignment
$\mu=(c_\sigma)_{\sigma\prec\tau}$ on its causal past. Recursively evaluate the
predecessor history $\mathbf h_{<\tau}^\mu(w)$ using these fixed choices and
put
\[
\mathsf P_\tau^\mu=E_\tau\times\mathbf{Hist}_{<\tau},
\qquad
P_\tau^\mu(w)=(Q_\tau(w),\mathbf h_{<\tau}^\mu(w)),
\]
\[
J_\tau^\mu(e,\mathbf h)=I_\tau(e,\mathbf h),
\qquad
I_\tau^\mu(w)=I_\tau(Q_\tau(w),\mathbf h_{<\tau}^\mu(w)).
\]
Then
\[
I_\tau^\mu=J_\tau^\mu\circ P_\tau^\mu.
\]
The construction is independent of a policy after $\mu$ is fixed.
\end{proposition}

\begin{proof}
The strict past of $\tau$ is downward closed, so fixed predecessor choices and
well-founded recursion uniquely determine every predecessor record. The
factorisation equation is the definition of the three displayed maps. The
exogenous map $Q_\tau$ was specified independently of the observation interface
and of every policy, so this is a past factorisation of the form required in
Definition~\ref{def:past-factor}.
\end{proof}

\begin{remark}
The same idea applies to a fixed branch of a contingent arena after restricting
to worlds that reach that branch; no separate formal object called a contingent
tree protocol is used here.
\end{remark}

\begin{corollary}[Trace-slice obstruction]
\label{cor:trace-slice-obstruction}
Fix $\tau$ and $\mu$. If two worlds have the same value of
$P_\tau^\mu$ and
\[
\bigcap_{i=0}^1
\{c\in C_\tau:
\lambda_\tau(w_i,\mathbf h_{<\tau}^\mu(w_i),c)\le\eps\}
=\varnothing,
\]
then no policy whose execution realises $\mu$ on both worlds is
$\eps$-capable.
\end{corollary}

\begin{proof}
Equal past records give one common observation. Such a policy must therefore
choose one common $c\in C_\tau$ at the two worlds. The empty intersection says
that every choice violates the stage bound at one of them, and aggregate loss is
at least that stage loss.
\end{proof}

\begin{remark}[Guardedness boundary]
The poset protocol models acyclic, universally executed dependencies; contingent
occurrence is represented by the arena tree. Cyclic simultaneous equations are
not operational tree executions. The Scott-continuous least-fixed-point model in
Appendix~\ref{sec:domain-feedback} gives a separate semantics from the
operational poset results.
\end{remark}

\subsection{A bounded-height countable effectivity instance}

The preceding arena theory is finite. The following narrowly bounded extension
has countably many nodes but no infinite plays and requires no limit-payoff or
transfinite semantics.

\begin{definition}[Countable bounded-height arena]
\label{def:countable-bounded-arena}
A countable bounded-height imperfect-information arena has the data of
Definitions~\ref{def:arena} and~\ref{def:ii-arena}, except that the rooted tree,
world set, and branching sets may be countable and there is a fixed finite bound
on the length of every root-to-leaf path. Predictor choice sets remain finite.
Reachability is defined recursively along each finite path, and aggregate loss
is the maximum of the finitely many reached-node losses. Patches are nonempty
information-saturated prefix-closed node sets.
\end{definition}

\begin{proposition}[Bounded-height activated subsheaf]
\label{prop:bounded-activated-sheaf}
On a countable bounded-height arena,
$\mathscr S_{\mathcal I}$ is a sheaf and
$\mathscr V_{\eps,\mathcal I}$ is a subsheaf for union covers.
\end{proposition}

\begin{proof}
A matching family gives one value to every information set occurring in the
union: choose a node of the information set that lies in the union, then choose a
covering patch containing that node. Saturation places the whole information set
in that patch, and matching makes the value
independent of the choice. This defines the unique glued strategy.

If a partial strategy is restricted to a smaller saturated prefix-closed patch,
every information set consulted on a path to a retained node is retained, so
reachability and loss at that node are unchanged. For a glued matching family
and a reached node $h$, choose a cover member containing $h$. Prefix closure and
saturation place every consulted ancestor information set there; matching makes
the local and global runs identical. Local validity therefore supplies the loss
bound. Neither argument uses finiteness of the total node set.
\end{proof}

\begin{definition}[Effective cover representation]
\label{def:effective-cover-representation}
Information sets are indexed by natural numbers, and finite choice-alphabet
elements are represented by fixed finite codes. A cover $(U_i)_{i\in\N}$ is
effectively represented when the predicate $J\subseteq U_i$ is decidable
uniformly in $(i,J)$. A locator is a total computable map $L$ satisfying
$J\subseteq U_{L(J)}$. A matching family is uniformly computable when a partial
recursive evaluator $\Phi(i,J)$ halts with value $s_i(J)$ on the represented
domain $J\subseteq U_i$.
\end{definition}

\begin{theorem}[Uniform effective gluing with a locator]
\label{thm:uniform-effective-gluing}
Let a countable bounded-height arena have computably indexed information sets
and finite choice alphabets. Let $(U_i)_{i\in\N}$ be a computably indexed cover
by information-saturated patches. Suppose there is a computable locator
$L$ such that $J\subseteq U_{L(J)}$ for every information set $J$, and suppose
a matching family $(s_i)_i$ is uniformly computable in the sense that
$(i,J)\mapsto s_i(J)$ is computable whenever $J\subseteq U_i$. Then its unique
global glue is computable. If every $s_i$ is activated-valid, so is the glue.
\end{theorem}

\begin{proof}
Define
\[
s(J)=s_{L(J)}(J).
\]
The locator and the uniform evaluator compute this value; the decidability of
the cover predicate is part of the representation convention and is not needed
once such a locator is supplied. If another cover
member contains $J$, matching says that its value agrees with the displayed
one, so $s$ is the unique coordinatewise glue. Activated validity follows from
Proposition~\ref{prop:bounded-activated-sheaf}.
\end{proof}


\begin{definition}[Finite invalidity witness]
\label{def:finite-invalidity-witness}
Index binary strategy coordinates by $\N$ and give $2^\N$ its standard product
topology. A finite binary string $\sigma$ is an invalidity witness when every
total strategy in the basic clopen cylinder $[\sigma]$ fails zero-loss activated
validity. A represented arena has uniformly c.e. finite invalidity when such
strings can be enumerated uniformly from its name and their cylinders cover the
entire invalid-strategy set.
\end{definition}

\begin{proposition}[WKL upper bound for represented causal selection]
\label{prop:causal-WKL-upper}
Consider represented countable bounded-height arenas with computably indexed
binary information-set choices and uniformly c.e. finite invalidity witnesses.
On the promised domain where a zero-loss global strategy exists, selecting one
is Weihrauch-reducible to $\mathsf{WKL}$. Uniformly finite alphabets admit a
fixed computable coding into Cantor space; the binary hypothesis avoids that
coding.
\end{proposition}

\begin{proof}
A total strategy is a point of $2^\N$. Enumerating the witness strings gives a
negative name for the valid set
\[
\mathscr V_{0,\mathcal I}(\mathcal T)
=2^\N\setminus\bigcup_{\sigma\ \mathrm{enumerated}}[\sigma],
\]
a co-c.e. closed subset of Cantor space. The promise says this closed set is
nonempty. One invocation of closed choice $\mathsf C_{2^\N}$ returns a valid
strategy, and Theorem~\ref{thm:fibsel-weihrauch} gives
$\mathsf C_{2^\N}\equiv_{\mathrm W}\mathsf{WKL}$.
\end{proof}

\begin{remark}
Proposition~\ref{prop:causal-WKL-upper} is an upper bound. The node-local
height-two arena of Theorem~\ref{thm:K-arena} enforces coordinatewise DNR
conditions and does not encode arbitrary prefix correlations, so no matching
WKL-hardness claim is made for that bounded-height class.
\end{remark}

\begin{definition}[Height-two halting-stage arena]
\label{def:K-arena}
Let $W=\N\times\N$. The root is a nature node with moves $n\in\N$ and
$\nu_\varepsilon(n,s)=n$. Its child $h_n$ is a predictor node with choice set
$\{0,1\}$ and singleton information set; each choice leads to a terminal node.
Put
\[
\lambda_{h_n}((n,s),c)=
\begin{cases}
0,&\neg\operatorname{Comp}(n,s,c),\\
1,&\operatorname{Comp}(n,s,c).
\end{cases}
\]
This is a primitive-recursively presented countable arena of height two.
\end{definition}

\begin{theorem}[Causal halting-stage effectivity obstruction]
\label{thm:K-arena}
For the height-two halting-stage arena:
\begin{enumerate}[label=(\roman*)]
\item its zero-loss global strategy space is exactly the class $\mathcal S$ of
Theorem~\ref{thm:Pi01}; hence it is nonempty and effectively closed but has no
computable member;
\item for the branch cover
\[
U_n=\{\varepsilon,h_n,h_n\!\cdot0,h_n\!\cdot1\},
\]
there is a matching family of activated-valid partial strategies such that each
individual section is computable, but no activated-valid matching family is
uniformly computable;
\item the arena has an $X$-computable zero-loss strategy exactly when $X$ has
PA degree. In particular, it has low and hyperimmune-free strategies.
\end{enumerate}
\end{theorem}

\begin{proof}
A global strategy is a binary function $g$. At node $h_n$, zero loss for every
world $(n,s)$ is equivalent to
\[
\forall s\,\neg\operatorname{Comp}(n,s,g(n)),
\]
which is exactly the defining condition for $g\in\mathcal S$. Part (i) follows
from Theorem~\ref{thm:Pi01}.

The patches $U_n$ cover the tree and distinct patches intersect only at the
nature root, so their strategy restrictions have no common predictor coordinate.
Use the classical selector $f$ from Theorem~\ref{thm:Kgame} and let the section
on $U_n$ choose the single bit $f(n)$. Each such finite constant section is a
computable object, and the family is matching and valid. If an activated-valid
matching family were uniformly computable, evaluating its section at $h_n$
uniformly in $n$ would produce a computable member of $\mathcal S$, contrary to
part (i). This proves (ii). Part (iii) is Corollary~\ref{cor:PAdegree}, since the
global strategy condition is the same class $\mathcal S$.
\end{proof}

\begin{remark}
Theorem~\ref{thm:uniform-effective-gluing} isolates the required uniformity.
Each branch-local section in Theorem~\ref{thm:K-arena} is finite and therefore
computable as an individual object. What fails is a single computable procedure
that supplies the correct finite section for every branch. The example is a
bounded-height countable instance, not a general complexity classification of
infinite arenas. In particular, the full patch is computable as a set of nodes
but has no computable valid section.
\end{remark}

\begin{corollary}[Division of labour]\label{cor:division}
Within the proved framework:
\begin{enumerate}[label=(\roman*)]
\item single-observer canonical diagrams have no $\mathsf D$;
\item perfect-information arenas are decided by backward induction and have no
full-site $\mathsf D$ under node-local nonemptiness;
\item finite distributed selectors canonically generate CSP compatibility
obstructions, while nested countable distributed selectors can exhibit the
effectivity obstruction of Theorem~\ref{thm:nested-distributed-K};
\item every finite distributed instance has a blind-order one-shot realisation,
and supplied world schedules permit world-dependent move orders;
\item the explicit finite perfect-information and one-shot imperfect-information
regimes are separated by the complexity boundary of Proposition~\ref{prop:pnp};
\item universally executed partially ordered cuts have a confluent operational
semantics and linear-extension-relative arena realisations, while contingent
stage occurrence remains tree-indexed.
\end{enumerate}
\end{corollary}


\begin{proof}
Clause (i) is Corollary~\ref{cor:distributed-one} together with discrete
compatibility. Clause (ii) is Theorem~\ref{thm:backward},
Corollary~\ref{cor:no-perfect-D}, and full-site collapse. Clause (iii) combines
Theorem~\ref{thm:distributed-equivalence} with
Theorem~\ref{thm:nested-distributed-K}. Clause (iv) is
Theorems~\ref{thm:oneshot} and~\ref{thm:one-shot-normal-form} together with
Proposition~\ref{prop:world-scheduled}. The complexity statement in (v) is
proved in Proposition~\ref{prop:pnp} by backward induction and a direct
polynomial 3SAT reduction for explicit imperfect-information arenas. Clause (vi) is
Theorems~\ref{thm:operational-execution} and~\ref{thm:operational-arena} together
with Proposition~\ref{prop:operational-confluence}.
\end{proof}

\section{Conjectural extensions}\label{sec:conjectures}

\begin{conjecture}[Structure-preserving repeated-move representation]
\label{conj:imperfect}
Theorem~\ref{thm:finite-ii-compilation} gives an extensional distributed
compilation of every finite pure-strategy imperfect-information arena, while
Proposition~\ref{prop:world-scheduled} gives a structure-preserving sufficient
class with supplied schedules. Characterise broader repeated-move perfect-recall
protocols admitting such presentations, develop expected-loss behavioural
semantics only after player and nature probabilities are specified, and
establish functoriality only after explicit arena and distributed-problem
morphisms have been fixed.
\end{conjecture}

\begin{conjecture}[Functoriality of contextual packaging]\label{conj:context}
After specifying morphisms of finite empirical support models and constrained
diagrams, the downward-closed packaging extends to a functor that preserves and
reflects the global-support criterion.
\end{conjecture}

\begin{conjecture}[Arena functoriality]\label{conj:arena-functor}
After fixing morphisms that preserve nature dynamics, predictor information,
choices, and losses, the finite-arena construction extends to a functor into a
category of sheaves on causal-patch sites and reflects global capability.
\end{conjecture}

\begin{conjecture}[Effective gluing beyond bounded height]
\label{conj:effective-gluing}
Theorem~\ref{thm:uniform-effective-gluing} treats bounded-height arenas with a
computable cover locator and a uniformly computable matching family. Extend the
result to a specified class of computably infinite, locally finite arenas with
represented strategy spaces and an explicit no-infinite-play or limit-payoff
semantics.
\end{conjecture}

\begin{remark}[No infinite-arena hierarchy claimed]
Except for the bounded-height countable results of
Definitions~\ref{def:countable-bounded-arena}--\ref{def:K-arena}, the arena
theorems in Section~\ref{sec:feedback} are finite. Arenas with unbounded finite
plays or infinite plays require termination or limit-payoff semantics,
representations of branching and strategies, and separate choice and
effectivity analyses. No arithmetical, hyperarithmetical, or analytic complexity
classification of such arenas is asserted.
\end{remark}







\section{Conclusion}\label{sec:conclusion}

The static tier separates observation-uniform choice, distributed compatibility,
effective realisation, and resource bounds. Initial-object collapse explains
why a single-observer canonical presentation has no compatibility
level, while distributed selectors generate CSP incidence diagrams in which
$\mathsf D$ outcomes can occur.
The extensional obstruction datum transports the diagnostic outcome when local
success and all nested global spaces are preserved. Terminal compression shows
that global spaces alone cannot preserve the distinction between $\mathsf L$
and $\mathsf D$. Isomorphism of the full relational constraint datum preserves
the solution space. Compactness and Helly theory give finite local certificates;
join trees give a positive local-to-global theorem; finite source integration
instantiates the compatibility diagnostic and has a $\#\mathrm P$-complete
counting problem; the affine cokernel and bijective-network fixed-point criteria
give complete certificates on their stated subclasses.
At the probabilistic level, Bell coupling failure is equivalent to a Farkas dual
certificate. The halting-stage selector space is a nonempty $\Pi^0_1$ class
without computable members, with oracle capability exactly at the PA degrees.
Countably nested distributed selectors realise arbitrary binary-tree path spaces;
when the tree has no computable path they give an overlapping-scope
$\mathsf K$ obstruction. Their input-uniform selection problem, and negatively
represented binary fibre selection, are Weihrauch-equivalent to $\mathsf{WKL}$.
Recursive capability is $\Sigma^0_3$-complete for the concrete
primitive-recursive game language of Definition~\ref{def:AHdef}. Computational-class separation and
deterministic-time resource separation are stated distinctly.

The causal tier fixes the arena independently of the strategy. Backward
induction characterises finite perfect-information capability, while patch
sheaves represent contingent partial strategies. Active sensing demonstrates
that obstruction labels depend on whether one uses the full patch site or a
branch cover. Blind-order one-shot arenas and finite distributed selectors are
equivalent for pure-strategy capability and worldwise loss, but not necessarily
for the presentation-relative diagnostic; supplied world schedules extend this to
world-dependent move orders. More general finite pure-strategy
imperfect-information arenas admit an extensional distributed compilation,
although a general structure-preserving converse remains open. Simplex-valued
behavioural coordinates glue as a sheaf, without being identified with
shared-randomness Bell mixtures. On the Bell one-shot information skeleton,
mixed-strategy marginal realisation is exactly local coupling. A finite causal
poset of universally executed cuts has a unique confluent operational execution;
each linear extension gives a policy-independent arena preserving aggregate
loss, while contingent occurrence remains natively tree-indexed. World-only
observations give an extensional distributed representation, and trace-relative
records instantiate past factorisation. Under global information recall, pure action recall is capability-redundant;
finite acyclic record congruences permit quotient backward induction. Finally, the height-two halting-stage arena shows that
pointwise computable local sections need not form a uniformly computable family;
with a computable locator and uniform local evaluator, effective gluing is
computable. The common lesson is that the location of failure is part of the
presentation: the same yes--no capability question may conceal local
infeasibility, incompatibility of local data, noncomputability, or resource
inadequacy, and the declared diagram determines which distinction is visible.

\appendix
\section{Domain-theoretic feedback}\label{sec:domain-feedback}

\begin{proposition}[Well-founded guarded recursion]\label{prop:guarded-recursion}
Let $(T,\prec)$ be well founded. For each $\tau\in T$, let $X_\tau$ be a set,
let $\mathrm{Dep}(\tau)\subseteq\{\sigma:\sigma\prec\tau\}$, and let
\[
F_\tau:W\times\prod_{\sigma\in\mathrm{Dep}(\tau)}X_\sigma\to X_\tau.
\]
For every world $w$ there is a unique family $(x_\tau)_{\tau\in T}$ satisfying
\[
x_\tau=F_\tau(w,(x_\sigma)_{\sigma\in\mathrm{Dep}(\tau)}).
\]
\end{proposition}

\begin{proof}
Well-founded recursion applies because every dependency of $\tau$ precedes
$\tau$. Uniqueness follows by well-founded induction from the determinism of the
maps $F_\tau$.
\end{proof}

Unguarded feedback motivates the following domain-theoretic formulation.

\begin{definition}[Parameterized feedback operator]\label{def:feedback}
Let $D$ be a pointed dcpo. A parameterized feedback operator is a map
\[
\Phi:\mathrm{Strat}\times W\times D\to D
\]
such that $F_{\pi,w}(d)=\Phi(\pi,w,d)$ is Scott-continuous for every
$(\pi,w)$.
\end{definition}

\begin{proposition}[Least fixed-point semantics]\label{prop:leastfp}
Every $F_{\pi,w}$ has the least fixed point
\[
\mu F_{\pi,w}=\bigvee_{n\ge0}F_{\pi,w}^{n}(\bot).
\]
\end{proposition}

\begin{proof}
This is the Kleene fixed-point theorem \cite{abramskyjung}. Scott continuity
preserves the supremum of the ascending iteration chain. If $x$ is any fixed
point, induction gives $F_{\pi,w}^n(\bot)\sqsubseteq x$ for every $n$, so the
supremum is least.
\end{proof}

\begin{proposition}[Query-invariant feedback semantics]
\label{prop:query-invariant}
Let $Q:D\to Y$ and suppose all fixed points of $F_{\pi,w}$ have the same
$Q$-value. Then $q_\Phi(\pi,w)=Q(\mu F_{\pi,w})$ is independent of fixed-point
choice.
\end{proposition}

\begin{proof}
The least fixed point is itself a fixed point, and all fixed points have the same
$Q$-image by hypothesis.
\end{proof}

\begin{remark}
With an interface and decoder, this query value defines a static validity
relation to which the selector lemma applies. No computability of the least
fixed point is asserted without effective representations of $D$ and $\Phi$.
\end{remark}


\section*{Formal verification}

The formal artifact associated with this article contains a 36-part reproduction
gate: twenty-four core Lean files and twelve deterministic Python experiments.
The Lean artifacts check the discrete selector, robust-selector, section,
world-pair, obstruction, discrete-ECD, finite-representation, parity, and
$\mathsf L/\mathsf D/\mathsf K/\mathsf R$ priority-classification cores.
The file \texttt{theorem\_D\_tri.lean} proves the parity contradiction underlying
Theorem~\ref{thm:Dtri}; it does not formalise the six-object category. The
separate Mathlib artifact \texttt{research/mathlib/k\_game.lean} verifies a
co-c.e. family with nonempty fibres and no total computable selector; it is not
part of the core gate.

The compactness/Helly certificates, initial-object collapse, distributed-selector
and NP-completeness theorems, nested distributed effectivity,
constraint-datum invariance and saturated-data recovery, obstruction transport
and compression boundary, join-tree, source-integration and counting lower-bound, affine cokernel, and
bijective-network holonomy theorems, categorical packaging, the binary
fibre-selection and nested prefix Weihrauch classifications, finite Bell
coupling and its Farkas certificate, $\Pi^0_1$/PA-degree formulation, REC
hardness reduction, Ackermann and deterministic-time separations,
extensive-form backward induction, perfect- and imperfect-information subsheaf
theorems, sensing, one-shot normal form, general finite pure-strategy
compilation, world-scheduled realisation, causal-poset operational execution and
linear-extension arena realisation, trace-relative factorisation, refinement-collapse, action-recall, arena-size,
and quotient-backward-induction results,
behavioural-coordinate gluing, mixed Bell realisation, bounded-height effective
gluing and WKL upper bound, height-two causal effectivity,
private-information team and division-of-labour
results, stochastic-kernel Bell construction, and static-compilation results are
proved in the text without dedicated Lean artifacts. Object-level
encodings and conjectures are likewise outside the machine-checked claim. The
command \texttt{make reproduce} runs the available gate. The recorded environment
uses Lean~4.33.0 and Python~3.13 on Linux; the separate K-game requires the
Mathlib project described in \texttt{research/mathlib/README.md}. Successful execution verifies the formal statements in the source files. It
does not verify the modelling adequacy of the article's additional
constructions.

\begin{center}
\begin{tabular}{@{}p{0.30\linewidth}p{0.42\linewidth}p{0.18\linewidth}@{}}
\toprule
Manuscript content & Artifact & Scope\\
\midrule
selector/robust/section cores & \texttt{theorem\_S.lean}, \texttt{theorem\_SE.lean}, \texttt{theorem\_ESS.lean} & checked core\\
triangle & \texttt{theorem\_D\_tri.lean} & parity core only\\
finite representation & \texttt{theorem\_Rep\_fin.lean} & combinatorial core\\
co-c.e. selector & \texttt{mathlib/k\_game.lean} & separate Mathlib check\\
extensive causal tier & none & manuscript proof\\
REC and time hierarchy & none & manuscript proof\\
\bottomrule
\end{tabular}
\end{center}

\begin{thebibliography}{99}

\bibitem{abramsky}
S.~Abramsky and A.~Brandenburger. The sheaf-theoretic structure of non-locality
and contextuality. \emph{New Journal of Physics} \textbf{13} (2011), 113036.

\bibitem{abramskyjung}
S.~Abramsky and A.~Jung. Domain theory. In S.~Abramsky, D.~M. Gabbay, and
T.~S.~E. Maibaum, editors, \emph{Handbook of Logic in Computer Science},
volume~3, pages 1--168. Oxford University Press, 1994.

\bibitem{beerifagin}
C.~Beeri, R.~Fagin, D.~Maier, and M.~Yannakakis. On the desirability of acyclic
database schemes. \emph{Journal of the ACM} \textbf{30} (1983), 479--513.

\bibitem{chsh}
J.~F. Clauser, M.~A. Horne, A.~Shimony, and R.~A. Holt. Proposed experiment to
test local hidden-variable theories. \emph{Physical Review Letters}
\textbf{23} (1969), 880--884.

\bibitem{fine}
A.~Fine. Hidden variables, joint probability, and the Bell inequalities.
\emph{Physical Review Letters} \textbf{48} (1982), 291--295.

\bibitem{jockusch1989}
C.~G. Jockusch, Jr. Degrees of functions with no fixed points. In
J.~E. Fenstad et al., editors, \emph{Logic, Methodology and Philosophy of
Science VIII}, pages 191--201. North-Holland, 1989.

\bibitem{jockuschsoare}
C.~G. Jockusch and R.~I. Soare. $\Pi^0_1$ classes and degrees of theories.
\emph{Transactions of the American Mathematical Society} \textbf{173} (1972),
33--56.

\bibitem{kleene}
S.~C. Kleene. \emph{Introduction to Metamathematics}. Van Nostrand, 1952.

\bibitem{pearl}
J.~Pearl. \emph{Causality: Models, Reasoning, and Inference}. Second edition,
Cambridge University Press, 2009.

\bibitem{rogers}
H.~Rogers. \emph{Theory of Recursive Functions and Effective Computability}.
McGraw-Hill, 1967.

\bibitem{soare}
R.~I. Soare. \emph{Recursively Enumerable Sets and Degrees}. Springer, 1987.

\bibitem{arorabarak}
S.~Arora and B.~Barak. \emph{Computational Complexity: A Modern Approach}.
Cambridge University Press, 2009.

\bibitem{eckhoff}
J.~Eckhoff. Helly, Radon, and Carath\'eodory type theorems. In
P.~M. Gruber and J.~M. Wills, editors, \emph{Handbook of Convex Geometry},
pages 389--448. North-Holland, 1993.

\bibitem{weihrauch}
K.~Weihrauch. \emph{Computable Analysis: An Introduction}. Springer, 2000.


\bibitem{brattkagherardi}
V.~Brattka and G.~Gherardi. Weihrauch degrees, omniscience principles and weak
computability. \emph{Journal of Symbolic Logic} \textbf{76} (2011), 143--176.

\bibitem{osbornerubinstein}
M.~J. Osborne and A.~Rubinstein. \emph{A Course in Game Theory}. MIT Press,
1994.

\bibitem{pitowsky}
I.~Pitowsky. \emph{Quantum Probability---Quantum Logic}. Springer, 1989.


\bibitem{schrijver}
A.~Schrijver. \emph{Theory of Linear and Integer Programming}. Wiley, 1986.

\end{thebibliography}

\end{document}
